\chapter{可泛化前馈重建}
\label{ch:13}

\paragraph*{本章阅读锚点}
本章的锚点是从逐场景优化转向前馈初始化。原始 3DGS 每个场景从头训练；前馈方法用网络直接预测相机、深度、点图或高斯，再交给优化器细化。


3D Gaussian Splatting 的原始方法依赖逐场景优化——为每个新场景运行数万次迭代的梯度下降，训练时间长达数十分钟甚至数小时。这一范式虽然能产生高质量的重建结果，却严重限制了 3DGS 在需要快速响应的实际场景中的适用性。机器人导航需要即时感知环境，增强现实（AR，Augmented Reality）需要实时融合用户视角，自动驾驶需要秒级构建驾驶场景的三维表示。一个自然的问题由此产生：能否训练一个通用模型，在前馈（feed-forward）推断中直接从输入图像预测三维高斯表示，将重建时间从分钟级压缩到秒级甚至毫秒级？

这一方向被称为可泛化前馈重建（Generalizable Feed-Forward Reconstruction），核心思想是用大规模多视角数据集训练一个深度网络，使其学习从二维观测到三维场景的通用映射。与原始 3DGS"为每个场景从头优化"的做法不同，前馈方法将场景的先验知识编码在网络的权重中，推理时仅需一次前向传播（single forward pass）即可完成重建。这一范式不仅大幅提升了重建速度，还赋予模型跨场景泛化能力：训练阶段未见过的场景也能被有效重建。

本章综述了 125 篇相关论文，按五个主题展开：

\textbf{多视角前馈重建}（第~\ref{sec:13_1}~节，71 篇）：在已知相机位姿的条件下，从稀疏或稠密多视角图像直接预测 3D 高斯基元。这是前馈重建中最成熟的方向，从像素对齐到体素对齐、从固定密度到自适应分配、从纯几何到融合生成先验，已经形成了完整的技术谱系。

\textbf{无标定与自由视角}（第~\ref{sec:13_2}~节，26 篇）：去除对已知相机位姿的依赖，从无标定图像同时估计相机参数和三维高斯表示。这一方向将前馈重建的适用范围从受控采集扩展到随意拍摄。

\textbf{单图到 3D 前馈方法}（第~\ref{sec:13_3}~节，8 篇）：从单张图像重建三维场景，面临单目设置下的固有歧义性，需要利用生成先验和几何约束来补偿信息不足。

\textbf{大规模重建模型}（第~\ref{sec:13_4}~节，9 篇）：将前馈重建扩展到数十甚至上百张输入视角的 360\textdegree{}场景级重建，核心挑战在于如何平衡长序列建模的计算开销与重建质量。

\textbf{特定领域前馈方法}（第~\ref{sec:13_5}~节，11 篇）：将前馈 3DGS 适配到事件相机、射频信号、人体渲染、体积可视化等特殊领域，展示了这一范式的广泛适用性。

\begin{figure}[htbp]
\centering
\resizebox{\textwidth}{!}{\begin{tikzpicture}[
  route/.style={rectangle, rounded corners=5pt, draw=#1, fill=#1!10,
    text width=3.2cm, minimum height=0.9cm, align=center, font=\small},
  method/.style={rectangle, rounded corners=3pt, draw=#1, fill=#1!6,
    font=\scriptsize, text width=3.0cm, align=center, minimum height=0.6cm},
  arr/.style={-{Stealth[length=4pt]}, techgray!60},
  >=Stealth
]

  \node[route=techgray, text width=4cm, font=\small\bfseries] (root) at (25.3, 5.6) {可泛化前馈重建方法分类};
  \node[route=techblue] (b0) at (7.7, 3.1) {多视角前馈};
  \node[method=techblue] (b0_l0) at (1.5, 0.6) {像素对齐\\{\tiny MVSGaussian, DepthSplat, TranSplat}};
  \node[method=techblue] (b0_l1) at (5.3, 0.6) {注意力聚合\\{\tiny HiSplat, FreeSplat, GS-LRM}};
  \node[method=techblue] (b0_l2) at (10.0, 0.6) {Epipolar注意力\\{\tiny MVSplat, MuGS, IDESplat}};
  \node[route=techgreen] (b1) at (19.5, 3.1) {无标定自由视角};
  \node[method=techgreen] (b1_l0) at (14.9, 0.6) {联合位姿估计\\{\tiny NoPoSplat, Splatt3R, AnySplat}};
  \node[method=techgreen] (b1_l1) at (19.5, 0.6) {基础模型引导\\{\tiny DUSt3R, MASt3R, VGGT}};
  \node[route=techorange] (b2) at (25.3, 3.1) {单图到3D};
  \node[method=techorange] (b2_l0) at (23.6, 0.6) {单图前馈\\{\tiny LRM, InstantMesh, Unique3D}};
  \node[method=techorange] (b2_l1) at (27.3, 0.6) {LRM架构\\{\tiny GS-LRM, Large reconstruction model}};
  \node[route=techpurple] (b3) at (33.2, 3.1) {大规模重建模型};
  \node[method=techpurple] (b3_l0) at (31.2, 0.6) {逐像素高斯\\{\tiny pixelSplat, MVSplat, DepthSplat}};
  \node[method=techpurple] (b3_l1) at (35.1, 0.6) {迭代精炼\\{\tiny ReSplat, GIFSplat, Diff3R}};
  \node[method=techpurple] (b3_l2) at (39.0, 0.6) {大场景扩展\\{\tiny GS-LRM, TokenGS, GlobalSplat}};
  \node[route=techblue] (b4) at (43.1, 3.1) {特定领域};
  \node[method=techblue] (b4_l0) at (43.1, 0.6) {特殊传感器\\{\tiny EventGS, RadarGS}};
  \node[method=techblue] (b4_l1) at (47.0, 0.6) {医学/卫星\\{\tiny Medical-GS, Satellite-GS}};
  \draw[arr] (root) -- (b0);
  \draw[arr] (root) -- (b1);
  \draw[arr] (root) -- (b2);
  \draw[arr] (root) -- (b3);
  \draw[arr] (root) -- (b4);
  \draw[arr] (b0) -- (b0_l0);
  \draw[arr] (b0) -- (b0_l1);
  \draw[arr] (b0) -- (b0_l2);
  \draw[arr] (b1) -- (b1_l0);
  \draw[arr] (b1) -- (b1_l1);
  \draw[arr] (b2) -- (b2_l0);
  \draw[arr] (b2) -- (b2_l1);
  \draw[arr] (b3) -- (b3_l0);
  \draw[arr] (b3) -- (b3_l1);
  \draw[arr] (b3) -- (b3_l2);
  \draw[arr] (b4) -- (b4_l0);
  \draw[arr] (b4) -- (b4_l1);
\end{tikzpicture}}
\caption{可泛化前馈重建方法分类}
\label{fig:ch13_feedforward_taxonomy}
\end{figure}

%=====================================================================
\section{多视角前馈重建}
\label{sec:13_1}

多视角前馈重建是可泛化 3DGS 研究中论文数量最多、技术发展最充分的方向。其基本框架为：给定 $N$ 张带有已知相机位姿的输入图像，训练一个前馈网络，在一次前向传播中预测场景的 3D 高斯基元参数（位置、协方差、不透明度、颜色），用于新视角合成和三维重建。本节按技术演进路线组织，从早期基础方法到最新的自适应表示和生成增强方案。

\begin{figure}[htbp]
\centering
\resizebox{\textwidth}{!}{\begin{tikzpicture}[
  box/.style={rectangle, rounded corners=4pt, draw=#1, fill=#1!10,
    text width=3.2cm, align=center, font=\small, minimum height=1cm},
  sbox/.style={rectangle, rounded corners=3pt, draw=#1!60, fill=#1!5,
    font=\scriptsize, text width=2.4cm, align=center, minimum height=0.5cm},
  arr/.style={-{Stealth[length=6pt]}, thick, techgray},
  darr/.style={-{Stealth[length=6pt]}, thick, techgray, dashed},
  >=Stealth
]

  \node[box=techblue] (f1) at (3.0, 0.6) {多视角输入(N张图像)};
  \node[box=techblue] (f2) at (10.1, 0.6) {特征提取(CNN/ViT)};
  \node[box=techgreen] (f3) at (18.2, 0.6) {跨视角聚合(Epipolar/注意力)};
  \node[box=techorange] (f4) at (28.5, 0.6) {高斯参数预测(位置/协方差/SH/$\alpha$)};
  \node[box=techpurple] (f5) at (37.7, 0.6) {可微分渲染(光栅化)};
  \node[font=\scriptsize, text=techgray] at (10.1, -0.6) {深度先验};
  \node[font=\scriptsize, text=techgray] at (37.7, -0.6) {新视角图像};
  \draw[arr] (f1) -- (f2);
  \draw[arr] (f2) -- (f3);
  \draw[arr] (f3) -- (f4);
  \draw[arr] (f4) -- (f5);
\end{tikzpicture}}
\caption{可泛化前馈重建网络架构}
\label{fig:ch13_feedforward_arch}
\end{figure}

\subsection{基础架构：像素对齐与前馈预测}

前馈 3DGS 的早期工作确立了"特征提取$\to$多视角交互$\to$高斯参数解码"的基本管线。以下两个方法分别代表了 MVS 驱动和深度先验驱动两条技术路线。

\paragraph{MVS 代价体驱动的像素对齐高斯重建}

读者已经从第~\ref{ch:01}~章了解到，原始 3DGS 为每个场景运行数万次迭代的梯度下降来优化高斯参数——从 COLMAP 稀疏点云出发，反复执行自适应密度控制（分裂、克隆、剪枝），整个优化过程耗时数十分钟。MVSGaussian（2024.05）~\cite{2405_12218}将这一范式从"逐场景反复优化"转变为"前馈一次推断"：训练一个通用网络，推理时仅需一次前向传播即可从输入图像直接预测场景的完整 3D 高斯基元参数。其核心洞察是：多视角立体视觉（Multi-View Stereo, MVS，通过多张图像间的特征匹配恢复场景深度的技术）天然提供像素级深度估计，而深度恰好将每个像素映射为三维空间中的一个点，这与像素对齐高斯表示的位置参数完美对应。相比此前采用射线体渲染的可泛化方法（如 IBRNet、ENeRF，每条射线需采样多个点，计算开销大），MVSGaussian 借助 MVS 深度直接获得高斯位置，并利用高斯 splatting 的实时渲染能力，实现了泛化性和效率的统一。

MVSGaussian 的改动可概括为四个互相衔接的阶段：MVS 代价体深度估计 $\to$ 像素对齐高斯表示 $\to$ 混合渲染 $\to$ 粗到细级联与一致性聚合优化。以下逐一展开。

\textbf{（1）深度估计。}利用 Feature Pyramid Network（FPN）从 $N$ 张源视角图像提取多尺度特征。接下来需要确定每个像素的深度。直觉上，如果把相机前方切成一系列平行于成像平面的等距平面（"扫掠平面"），再把其他视角的特征投影到这些平面上，那么正确深度处来自不同视角的特征应该高度一致，错误深度处则不匹配。

这一投影操作通过\emph{可微单应变换}（differentiable homography，一种对平面间坐标映射的参数化方法）实现：给定源视角和目标视角的相机内参 $K$、旋转 $R$、平移 $t$，以及深度候选 $z$，源视角特征被 warp 到目标视角的第 $z$ 个扫掠平面：

\begin{equation}
H_i(z) = K_i R_i \left(I + (R_i^{-1} t_i - R_t^{-1} t_t) \frac{a^T R_t}{z}\right) R_t^{-1} K_t^{-1}
\end{equation}

其中 $a$ 为主轴方向向量。直观理解，$H_i(z)$ 就像一张可变形的橡皮膜，将源视角的特征图像"拉伸"并投影到目标视角在深度 $z$ 处的平面上——深度不同，拉伸方式也不同。

在所有深度平面上，$N$ 个 warp 后特征之间的\emph{方差}被汇总为一个三维张量，即\emph{代价体}（cost volume）——它编码了每个像素在每个深度处的"多视角匹配不一致程度"。代价体经 3D CNN 正则化后输出深度概率分布，深度概率经 softmax 归一化后逐深度加权求和，得到最终深度图：

\begin{equation}
d = \sum_{m=1}^{D} p(z_m) \cdot z_m, \quad p(z_m) = \text{softmax}(\text{CNN}_{\text{3D}}(\text{CostVol}(\{F_i\}_{i=1}^N)))[m]
\end{equation}

即最终深度是所有深度候选的加权平均，权重由网络学习到的匹配概率决定——匹配越好的深度获得越大的概率。

\textbf{（2）像素对齐高斯表示。}估计的深度将每个像素解投影为三维点，直接作为该像素对应高斯的位置 $\mu = \pi^{-1}(x, d)$，其中 $\pi^{-1}$ 为解投影操作，$x$ 和 $d$ 分别为像素坐标和估计深度。

为编码丰富的特征，MVSGaussian 采用了两级特征增强策略。第一级，将源视角特征 warp 到目标视角，通过池化网络 $M$ 聚合多视角信息得到 $f_v = M(\{f_i\}_{i=1}^N)$——聚合方式为：先计算多视角特征的均值和方差，再通过 MLP 生成加权权重，以 soft-argmax 方式融合。第二级，利用 2D UNet 对 $f_v$ 进行空间增强得到 $f_g$，使每个像素的特征具备空间上下文感知能力——这一设计源于 splatting 的本质特性：一个高斯影响周围多个像素的颜色，因此仅编码单像素信息是不够的。

高斯参数通过独立的 MLP 头解码：

\begin{align}
s &= \text{Softplus}(h_s(f_g)), \quad r = \text{Normalize}(h_r(f_g)), \\
\alpha &= \text{Sigmoid}(h_\alpha(f_g)), \quad c = \text{Sigmoid}(h_c(f_g))
\end{align}

其中 $s$、$r$、$\alpha$、$c$ 分别为缩放（正数）、旋转四元数（单位范数）、不透明度（$[0,1]$）和颜色（$[0,1]$）。值得注意的是，MVSGaussian 直接回归 RGB 颜色而非原始 3DGS 的球谐（SH）系数，因为消融实验表明从特征学习球谐系数的可泛化性不佳——SH 系数对视角方向的高度敏感性在跨场景迁移中反而成为负担。

\textbf{（3）混合渲染。}MVSGaussian 的一个关键设计是混合渲染策略，动机来自 splatting 的结构特性。纯高斯 splatting 建立高斯与像素之间的\emph{多对多}映射——一个高斯影响周围多个像素的颜色，一个像素的颜色由多个高斯共同决定。这种复杂的空间交互虽然表达力强，却不利于前馈网络的可泛化学习：网络需要同时学会"如何摆放高斯"和"高斯之间如何交互"，两个问题纠缠在一起增加了学习难度。

为此，作者引入一种简化的\emph{深度感知体渲染}作为互补：利用聚合的多视角特征 $f_v$（未经 UNet 空间增强）解码辐射值 $r$ 和体密度 $\sigma$，沿每条射线\textbf{仅采样一个}深度感知点，像素颜色为 $c = (1 - e^{-\sigma}) \cdot r$。直觉上，这退化为一种一对一的像素对齐 splatting——每个三维点仅贡献一个像素的颜色，映射关系被极大简化。最终渲染图像取高斯 splatting 渲染和体渲染的平均值，兼顾高斯表示的表达力和体渲染的泛化稳定性。消融实验证实，混合渲染显著优于单独使用任一方式。

\textbf{（4）粗到细级联与一致性聚合优化。}整个管线构建为两阶段级联结构：粗阶段采样 $D=64$ 个深度平面估计初始深度并渲染粗略图像，细阶段在此基础上仅采样 $D=8$ 个深度平面进行精炼，粗阶段的深度和渲染结果传递给细阶段作为先验。训练损失为 MSE、SSIM 和感知损失的加权和 $\mathcal{L}^k = \mathcal{L}_{\text{mse}} + \lambda_s \mathcal{L}_{\text{ssim}} + \lambda_p \mathcal{L}_{\text{perc}}$，其中 $\lambda_s = 0.1$，$\lambda_p = 0.05$，总损失为两阶段损失加权求和 $\mathcal{L} = \sum_k \gamma^k \mathcal{L}^k$（$\gamma^1 = 0.5$，$\gamma^2 = 1$），粗阶段以较小权重参与训练。

在逐场景优化阶段，MVSGaussian 提出基于多视角几何一致性的聚合策略。由于前馈模型在多个目标视角独立预测高斯点云，直接拼接会导致大量冗余和噪声点。核心思想是：对同一三维点在不同视角的深度预测计算\emph{重投影误差}。具体来说，像素 $p$ 在参考深度图 $D_0$ 中对应的深度为 $D_0(p)$，将其投影到邻近视角再反投影回来得到 $p'$，则位置误差 $\epsilon_p = \|p - p'\|_2$，相对深度误差 $\epsilon_d = |D_0(p) - d'| / D_0(p)$。通过\emph{动态一致性检查}进行筛选：深度在少数视角中重投影误差极低，或在多数视角中误差较低，即视为可靠。不可靠点被剔除，可靠点保留，为后续优化提供高质量初始化。优化损失为 $\mathcal{L}_{\text{ft}} = (1-\lambda_{\text{ft}})\mathcal{L}_1 + \lambda_{\text{ft}} \mathcal{L}_{\text{D-SSIM}}$，其中 $\lambda_{\text{ft}} = 0.2$。

MVSGaussian 在 DTU、Real Forward-facing、NeRF Synthetic 和 Tanks and Temples 四个数据集上验证。相比可泛化 NeRF 方法（IBRNet、ENeRF、MVSNeRF 等），实现了更好的合成质量和实时渲染速度（300+ FPS）。相比逐场景优化 3DGS，在 Real Forward-facing 数据集上仅需 45 秒优化即超越 3DGS 的 10 分钟结果（13.3 倍加速），同时保持相当的渲染速度。

\paragraph{双分支深度模型与深度-渲染互增益}

与 MVSGaussian 依赖 MVS 代价体进行纯多视角特征匹配的路线不同，DepthSplat（2024.10）~\cite{2410_13862}在 MVS 分支之外引入预训练单目深度先验作为第二分支，构建双分支深度模型。它同时证明了一个有趣的现象：高斯 splatting 渲染损失可作为深度模型的无监督预训练目标——深度估计与视角合成两个任务互相增益。

\textbf{Baseline 与改进动机。}MVSGaussian 等基于 MVS 代价体的方法虽然能有效估计深度，但在遮挡区域、弱纹理区域和反射面等场景中，纯多视角特征匹配本质上会失效——没有可供匹配的可靠特征。DepthSplat 的关键洞察是：预训练的单目深度模型（如 Depth Anything V2）在海量数据上学习了丰富的单目几何先验，能够为这些困难区域提供互补信息。反过来，仅靠单目深度先验也难以获得尺度一致的多视角深度，因此需要两个分支互补。

\textbf{双分支深度架构。}DepthSplat 的深度模型包含两个并行分支：

\textit{多视角特征匹配分支}：使用轻量级共享 ResNet 从 $N$ 张输入图像提取 $s$ 倍下采样特征（$s$ 通过控制步长为 2 的 $3\times3$ 卷积层数灵活调节，如两层对应 $s=4$，三层对应 $s=8$），再通过六层 Swin Transformer 的自注意力和交叉注意力交换跨视角信息。当 $N>2$ 时，每个参考视角仅与距离最近的两个邻近视角执行交叉注意力，兼顾精度与效率。匹配阶段采用平面扫掠立体（plane-sweep stereo）：在一系列均匀采样的深度平面上将邻近视角特征投影回参考视角，寻找最佳匹配。具体地，均匀采样 $D$ 个深度候选 $d_m$，通过差分 warp 和点积相关构建代价体 $C_i \in \mathbb{R}^{H/s \times D}$，代价体中的每个值反映该像素在该深度候选下的跨视角匹配质量。

\textit{单目深度特征分支}：直接使用 Depth Anything V2 的冻结 ViT 骨干（而非训练深度估计头）提取单目特征 $F^i_{\text{mono}} \in \mathbb{R}^{H/s \times C_{\text{mono}}}$。ViT 骨干的 patch 大小为 14，输出分辨率为原图的 $1/14$，通过双线性插值调整至与代价体相同的分辨率。整个分支以冻结方式运行，保留预训练时学到的丰富单目几何先验。

\textbf{特征融合与深度回归。}两个分支的特征在通道维度拼接后，送入 2D UNet 进行融合和深度回归。直觉上，MVS 代价体提供多视角一致的几何约束，单目特征为匹配失效区域补充几何先验，拼接使 UNet 同时看到两种互补信号。深度预测通过 softmax 归一化和加权平均得到。

此外采用层次化匹配架构：粗阶段在低分辨率估计初始深度，精阶段在 $2\times$ 更高分辨率特征上于粗深度邻域内搜索更精细的对应，因已有粗深度指引，精阶段只需搜索更小范围并构建更小的代价体，兼顾效率和精度。最终深度图通过 DPT head 上采样至全分辨率。

\textbf{高斯参数预测。}深度图通过相机参数解投影为高斯中心位置 $\mu_j$。其余参数（不透明度 $\alpha_j$、协方差 $\Sigma_j$、颜色 $c_j$）由额外 DPT head 从拼接的图像、深度和特征信息预测。整个高斯集合通过标准 3DGS 的 splatting 操作渲染新视角。

\textbf{训练损失与互增益机制。}DepthSplat 同时在两个任务上训练：深度估计与视角合成。对于深度估计任务，在逆深度域上使用 $L_1$ 损失加梯度损失：

\begin{equation}
\mathcal{L}_{\text{depth}} = |D_{\text{pred}} - D_{\text{gt}}| + \alpha |\nabla_x D_{\text{pred}} - \nabla_x D_{\text{gt}}| + \beta |\nabla_y D_{\text{pred}} - \nabla_y D_{\text{gt}}|, \quad \alpha=\beta=20
\end{equation}

梯度损失惩罚深度图在 $x$ 和 $y$ 方向上的边缘不匹配，促使预测深度保持锐利的几何边界。

对于视角合成任务，使用 MSE 与 LPIPS 的组合损失：

\begin{equation}
\mathcal{L}_{\text{gs}} = \sum_{m=1}^{M} \left(\text{MSE}(I^m_{\text{render}}, I^m_{\text{gt}}) + \lambda \cdot \text{LPIPS}(I^m_{\text{render}}, I^m_{\text{gt}})\right), \quad \lambda = 0.05
\end{equation}

其中 $M$ 为单次前向传播中渲染的新视角数量，LPIPS 权重 $\lambda=0.05$ 防止感知损失压倒像素级重建精度。

关键发现是：用高斯 splatting 渲染损失在大规模 RealEstate10K 数据集（约 67K YouTube 视频）上无监督预训练的深度模型，迁移到纯深度估计任务后在 TartanAir 和 KITTI 等数据集上显著优于从零训练的模型。这表明三维重建的渲染信号提供了比稀疏深度监督更丰富的几何学习信号，高斯 splatting 与深度估计之间构成了互相增益的正反馈循环。

在 ScanNet 和 DL3DV 数据集上，DepthSplat 实现了当时最佳的前馈重建性能，从 12 张输入视角（$512\times 960$ 分辨率）的完整重建仅需 0.6 秒。消融实验表明，移除单目特征分支后深度和渲染质量均显著下降，尤其是在弱纹理和反射面上；移除代价体后仅靠单目先验也难以获得尺度一致的多视角预测——两个分支缺一不可。

在上述基础架构之上，后续方法沿着多个方向进行了改进。

\paragraph{深度置信度引导的局部匹配与单目修补}

与 DepthSplat 在双分支架构中冻结 Depth Anything V2 骨干提取单目特征不同，TranSplat（2024.08）~\cite{2408_13770}聚焦于多视角特征匹配本身的精度问题——尤其是在包含大量相似区域和视角间非重叠区域的场景中，纯特征点积匹配容易产生误匹配。TranSplat 的核心思路是两步走：先用深度置信度图引导精确的局部匹配，再用单目深度先验修补非重叠区域的深度。

\textbf{粗到精匹配管线。}特征提取阶段采用与 DepthSplat 类似的 CNN 加 Transformer 架构（Swin Transformer 的局部窗口注意力实现跨视角信息交换），并通过 Squeeze-Excitation 层将相机参数注入 CNN 特征提供全局空间信息。同时，从 DepthAnythingV2（base 规格）提取最后一层特征作为深度先验 $D^i \in \mathbb{R}^{H/4 \times W/4 \times C}$ 和相对深度输出。匹配管线分为两个阶段：粗匹配阶段通过平面扫掠立体（$D{=}128$ 个深度候选）和特征点积获取初始深度分布 $d^i_{\text{coarse}}$；精匹配阶段引入深度感知可变形匹配 Transformer（Depth-Aware Deformable Matching Transformer, DDMT），用可变形采样模块（采样 $P{=}4$ 个可变形点）聚合局部空间信息，再由 DMT 根据融合了粗深度、深度特征和相机参数的深度感知特征，自适应地预测采样偏移和深度感知注意力权重 $W^i$，生成深度残差 $\Delta d^i$ 进行精细修正：$d^i_{\text{fine}} = d^i_{\text{coarse}} + \Delta d^i$。直觉上，可变形采样使网络不再仅依赖逐像素对齐的特征点积，而是在局部邻域内搜索更优的匹配位置；DMT 则通过注意力机制将"哪里值得仔细看"的判断融入深度估计过程。

\textbf{Depth Refine U-Net。}粗到精匹配仅覆盖视角间有重叠的区域。对于非重叠区域，TranSplat 设计了 Depth Refine U-Net，将精确的多视角匹配几何一致性与 DepthAnythingV2 的可靠相对深度先验融合：以输入图像和相对深度图作为输入，通过 UNet 架构优化深度分布。这一设计与 DepthSplat 的双分支融合策略不同——DepthSplat 在特征层面融合，TranSplat 在深度估计的输出层面融合。

在 RealEstate10K 和 ACID 数据集上，TranSplat 仅用 200K 训练迭代（MVSplat 使用 300K）即超越所有基线，推理速度为 0.06 秒（与 MVSplat 的 0.08 秒相当）。在零样本跨数据集测试中，TranSplat 在 DTU 数据集上同样展现出色的泛化能力，尤其是在弱纹理区域（如冰箱表面）和重复图案区域（如墙壁）显著优于先前方法。

\paragraph{粗到细三级多尺度高斯预测}

DepthSplat 和 TranSplat 均生成单一尺度的逐像素高斯，所有高斯具有相同的空间分辨率。HiSplat（2024.10）~\cite{2410_06245}观察到这种单尺度表示存在固有局限：无法同时捕捉大尺度结构（如墙面、地板的整体走向）和小尺度纹理细节（如花纹、文字），导致重建结果中出现模糊、裂纹和定位偏移。其核心思路是引入粗到细的层次化策略，生成多尺度高斯基元并通过跨尺度交互模块协同优化。

\textbf{层次化特征提取。}HiSplat 采用 U-Net 编码器-解码器与多视角 Transformer 的混合骨干，编码器和解码器各包含三个 $3{\times}3$ 残差块，通道数逐级为 $[32, 64, 128]$，每个阶段分辨率减半/倍增。编码器末端通过多视角 Transformer 的交叉注意力和自注意力交换跨视角信息，注入到解码流中。此外，引入冻结 DINOv2 特征与各尺度 U-Net 特征相加，增强泛化能力。

\textbf{三级深度与高斯预测。}整个管线分为三个阶段。最粗阶段（Stage 1，$1/4$ 分辨率）使用代价体匹配估计基准深度——与 DepthSplat 和 TranSplat 相同的平面扫掠方式，但仅在最低分辨率执行以避免高分辨率代价体的计算负担。后续两个阶段（Stage 2 和 Stage 3，分别为 $1/2$ 和全分辨率）不再使用代价体，而是通过\emph{误差感知模块}（Error Aware Module）以上一阶段渲染图像与输入图像的差值图作为驱动信号，预测深度偏移量 $\beta_i \in [0,1]$ 和高斯特征。深度偏移受最大深度系数 $\eta \in [0,1]$ 约束：$D_i = D_{\text{stage1}} + \eta \cdot \beta_i \cdot D_{\text{stage1}}$，确保小尺度高斯始终位于大尺度高斯邻域内。直觉上，大尺度高斯构成场景"骨架"，小尺度高斯在骨架附近补充细节或修正结构误差。

\textbf{调制融合模块。}三级高斯通过\emph{调制融合模块}（Modulating Fusion Module）聚合为最终的多尺度高斯集合。该模块将跨尺度特征与误差图拼接后，动态调整前序高斯的不透明度——在需要精细补偿的区域增大细粒度高斯的贡献，在已充分重建的区域保持粗粒度高斯的主导地位——实现多尺度高斯的协同优化。

在 RealEstate10K 上，HiSplat 超越 MVSplat 达 +0.82 dB PSNR，超越 TranSplat 达 +0.52 dB PSNR，首次在两视角重建任务中突破 27 dB PSNR。在 ACID 上同样一致领先，零样本迁移至 DTU 和 Replica 时渲染质量和几何定位精度均显著优于单尺度方法。推理时仅采用最后一阶段的融合高斯，速度为 0.39 秒。

\paragraph{面向长序列的像素级三元融合}

前述方法（DepthSplat、TranSplat、HiSplat）均假设少量输入视角（通常 2--3 张），生成的高斯数量有限。当输入视角数量增加时，像素对齐高斯在重叠区域产生大量冗余。FreeSplat（2024.05）~\cite{2405_17958}针对室内大场景的长序列重建需求，提出了低成本跨视角聚合框架和像素级三元融合方法。其核心挑战在于：随着输入视角数量增加，像素对齐高斯在重叠区域产生大量冗余，需要一种高效的增量式融合策略来消除重复基元。

\textbf{与 DepthSplat 的架构差异。}DepthSplat 使用 Transformer 架构处理固定数量的输入视角，计算量随视角数二次增长。FreeSplat 选择纯 CNN 骨干（EfficientNet/ResNet）替代 Transformer 进行二维特征提取，使计算开销随视角数线性增长，从而支持任意长度的输入序列。在邻近视角间自适应构建代价体（$K{=}128$ 个深度平面，范围 $[d_{\text{near}}{=}0.5, d_{\text{far}}{=}15.0]$ 米），而非全局构建。具体做法是根据位姿邻近度自适应选择 $N$ 个邻近视角，通过平面扫掠和余弦相似度构建代价体，再经 UNet++ 解码器预测深度图。

\textbf{像素级三元融合。}每个像素被解投影为三维高斯三元组 $(\mu_l^t, w_l^t, f_l^t)$——中心坐标、融合权重（$(0,1)$ 之间的标量）和潜在特征向量。（$t{=}1, 2, \ldots, T$）：每步先将全局高斯中心投影到当前视角，通过像素级最近邻搜索和深度差值阈值（比例阈值 $\rho{=}0.05$）筛选建立对应对 $(i, m_i)$，再以加权求和方式融合对齐的高斯中心坐标 $\mu_g^t(m_i) = \frac{w_l^t(i) \cdot \mu_l^t(i) + w_g^{t-1}(m_i) \cdot \mu_g^{t-1}(m_i)}{w_l^t(i) + w_g^{t-1}(m_i)}$，并通过轻量 GRU 网络聚合局部与全局的潜在特征。直觉上，这类似于传统 TSDF Fusion 的增量式深度融合思路，但在高斯基元的潜在特征空间中执行，而非仅操作几何量。在 ScanNet 10 视角长序列实验中，PTF 将渲染速度从 39 FPS 提升至 72 FPS，显著减少了冗余高斯。

FreeSplat 在 ScanNet 和 Replica 上的渲染质量和深度精度均一致优于 pixelSplat 和 MVSplat，同时支持自由视角训练策略，确保在更广视角范围内的鲁棒合成。

\paragraph{加权浮游体去除与深度正则化微调}

FreeSplat++（2025.03）~\cite{2503_22986}将 FreeSplat 的室内大场景重建扩展到全场景级别，解决三个核心问题：长序列处理效率、跨视角浮游体（floater，悬浮在场景空间中的半透明碎片状伪影）噪声、以及前馈质量与逐场景优化之间的差距。

\textbf{与 FreeSplat 的关键差异。}(1)~FreeSplat++ 将解投影分辨率降至全分辨率的 $1/2$（而非 FreeSplat 的全分辨率），使全场景重建中的高斯数量保持在可承受范围内，同时性能几乎不受影响。(2)~提出加权浮游体去除（Weighted Floater Removal, WFR）模块。该模块利用 PTF 过程中累积的融合权重作为可靠度指标：权重过低的高斯表示其在多视角融合中被很少观测到，极可能是浮游体噪声。WFR 据此降低这些高斯的不透明度，作为显式深度融合方法有效抑制噪声深度图引入的浮游体。(3)~设计深度正则化的逐场景微调：利用前馈预测阶段获得的多视角一致深度图作为额外约束，在保持几何精度的同时精炼高斯参数的渲染质量。在 ScanNet 全场景实验中，FreeSplat++ 的前馈预测已在几何精度上优于传统逐场景优化 3DGS；经 10K 次迭代微调后，渲染质量显著提升，而训练时间仅为传统 3DGS（30K 次迭代）的 $1/3$。

\paragraph{纯 Transformer 直接回归逐像素高斯}

GS-LRM（2024.04）~\cite{2404_19702}代表了与上述方法不同的技术路线——不依赖 MVS 代价体或单目深度先验进行显式深度估计，而是用纯 Transformer 架构直接回归逐像素高斯参数。这一设计的关键优势在于简洁性：网络只需学会"看到图像后直接输出三维表示"，无需手工设计几何推理模块。与 DepthSplat 的双分支深度模型不同，GS-LRM 将三维重建建模为标准的序列到序列映射问题。关于 GS-LRM 的完整技术细节（架构设计、训练策略、大规模扩展），详见第~\ref{sec:13_4}~节。

\textbf{架构设计。}输入为 $N$ 张位姿图像，RGB 与 Pl\"ucker 射线坐标（六维参数化，编码每个像素对应的相机射线几何）在通道维度拼接形成 9 通道输入，按 patch 大小 $p{=}8$ 切分为 $d{=}1024$ 维 token。所有视角的 token 拼接后通过 $L{=}24$ 层 Transformer 块（Pre-LayerNorm + 16 头自注意力 + 两层 MLP，MLP 隐藏维度 4096，GeLU 激活），每个输出 token 经单个线性层解码为 12 维高斯参数（3 通道 RGB、3 通道缩放、4 通道旋转四元数、1 通道不透明度、1 通道射线距离），高斯中心位置由射线距离和相机参数计算 $\mathbf{x} = \mathbf{o}_{\text{ray}} + t \cdot \mathbf{d}_{\text{ray}}$。

\textbf{与 DepthSplat 的对比。}GS-LRM 的关键区别在于三点：其一，不构建代价体，依靠 Transformer 自注意力隐式学习多视角对应关系，避免了几何推理模块的手工设计；其二，参数化中缩放初始化为 $\min(\exp(G_{\text{scale}} - 2.3), 0.3)$，不透明度初始化为 $\sigma(G_{\text{opacity}} - 2.0)$（约 0.1），工程上确保训练稳定性；其三，模型同时适用于物体级（Objaverse 730K 物体，4 输入视角）和场景级（RealEstate10K 80K 视频，2 输入视角）重建，展现架构的通用性。

模型含 300M 参数，采用 MSE 损失加 VGG-19 感知损失（权重 $\lambda{=}0.5$），在 64 张 A100 上训练，分两阶段（$256{\times}256$ 预训练 80K 步 + $512{\times}512$ 微调 20K 步）。推理时在单张 A100 上仅需 0.23 秒。物体级模型在 GSO 基准上超越 Instant3D 达 3.98 dB PSNR，场景级模型在 RealEstate10K 上超越 pixelSplat 达 2.2 dB PSNR。

\subsubsection{高斯基元分配策略的演进}

本路线介绍从固定逐像素分配到自适应密度控制的高斯基元分配策略演进。像素对齐方法将每个像素映射为一个三维高斯，导致高斯数量随输入分辨率和视角数量线性增长，产生大量冗余。这一观察催生了多个改进方向。

\paragraph{SplatWeaver: 混合专家驱动的基数高斯密度分配}

像素对齐方法将每个像素映射为一个三维高斯，导致高斯数量随输入分辨率和视角数量线性增长，产生大量冗余。这一观察催生了自适应密度分配方向的研究。SplatWeaver（2026.05）~\cite{2605_07287}同样关注"逐像素一高斯"带来的冗余问题，但与 SparseSplat 基于信息熵的概率采样不同，SplatWeaver 引入了\emph{基数高斯专家}（Cardinality Gaussian Experts）机制——将自适应密度分配建模为混合专家（Mixture of Experts）问题。其核心洞察是：与其为每个像素统一预测固定数量的高斯，不如训练一组专家，每个专家专注于预测 0 到 $M$ 个不等数量的基元（$M{=}3$），再由路由机制自适应决定每个空间位置调用哪个专家。

\textbf{专家路由与高频先验}。SplatWeaver 的路由流程为：先用 DINOv2 编码器和多视角几何 Transformer 提取像素级特征 $F_n \in \mathbb{R}^{H \times W \times D}$，再由像素级路由器 $R$（两层线性层加 ReLU）对每个像素的特征输出路由 logits，经 Gumbel-Softmax（Straight-Through 估计器）获得离散 one-hot 路由概率 $p_n(i,j)$，选择 top-1 专家。每个被选中的专家仅预测高斯位置 $\mu$ 和潜在特征 $F_l$，剩余属性（尺度、旋转、不透明度、颜色）则通过注意力聚合 $K$ 个近邻隐藏高斯的特征来预测——这一设计与 SparseSplat 的三维局部预测思路异曲同工。为保证路由的可靠性，SplatWeaver 引入基于离散小波变换（DWT）的高频先验引导模块：对输入图像做 DWT 得到三个高频子带 $(\text{LH}, \text{HL}, \text{HH})$，计算高频能量图 $\text{HF} = \sqrt{\text{LH}^2 + \text{HL}^2 + \text{HH}^2}$，该能量图与 3DGS 密集重建的高斯分布高度一致。论文观察到高频能量高的区域通常对应结构细节丰富的区域，需要更多高斯基元。因此，高频先验被注入像素特征 $F_n^f = F_n \odot \sigma(\Gamma(\text{LH}, \text{HL}, \text{HH})) + F_n$（其中 $\Gamma$ 为变换层，$\sigma$ 为 Sigmoid），同时作为路由正则化的监督信号：像素按高频能量排序后，高能量像素被鼓励路由到更大基数的专家。

\textbf{训练与实验}。SplatWeaver 在九个数据集上联合训练（Hypersim、ARKitScenes、BlendedMVS、ScanNet++、CO3D-v2、Objaverse、Unreal4K、WildRGBD、DL3DV），无需标定位姿输入。路由正则化损失中包含空专家（null expert，预测零个高斯），正则化权重 $\lambda_{3\%} = 2\%$、$\lambda_{2\%} = 2\%$、$\lambda_{1\%} = 20\%$。在 DL3DV 基准上，SplatWeaver 在 4/8/16/24 视角设置下一致优于 DepthSplat 和 NoPoSplat；在 Mip-NeRF 360 上零样本泛化同样表现优异，同时使用的高斯基元数量显著少于像素对齐方法。

\paragraph{SparseSplat: 信息熵引导的自适应采样与局部属性预测}

SplatWeaver 通过混合专家路由实现密度分配，但路由决策仍依赖学习到的路由器权重。SparseSplat（2026.04）~\cite{2604_03069}则从信息论角度出发，不再给每个像素硬性分配一个高斯，而是先用信息熵判断"哪里需要高斯、需要多少"，再用局部邻域而非全局特征来预测每个高斯的属性。

\textbf{基线与问题}：MVSGaussian 和 DepthSplat 代表的前馈 3DGS 普遍采用"逐像素一高斯"（pixel-aligned）范式——输入图像有多少像素，就生成多少高斯基元。这一看似自然的设计带来了两个结构性缺陷。一是\emph{分布错配}——像素对齐强制生成均匀密度的点云，整片白墙和纹理丰富的花丛分配到相同密度的高斯，而优化型 3DGS 的做法截然相反：在无纹理区域用少量大型高斯覆盖大片面积，在高频纹理区域用密集的小型高斯捕捉细节。二是\emph{感受野错配}——主干网络（如 DPT）拥有全局感受野，从整张图像中提取上下文信息，但预测头却只从单个像素的特征向量回归高斯属性。这种"用全局视野做局部决策"的设计偏离了 3DGS 优化的本质：经典 3DGS 的梯度仅从高斯投影覆盖的局部像素区域回传，一个高斯的属性本质上由它自身及其三维邻域决定，而非整张图像的全局信息。

\textbf{核心改动}：针对这两个错配，SparseSplat 提出三阶段管线。(1) \emph{特征与深度提取}：复用预训练的 DepthSplat 主干（冻结权重），提取二维特征图 $F \in \mathbb{R}^{C \times H \times W}$ 和深度图 $D \in \mathbb{R}^{H \times W}$。主干可替换，方法与具体深度估计骨干无关。(2) \emph{基于信息熵的自适应采样}：抛弃逐像素分配，改为根据局部信息丰富度概率性地采样像素，仅在被选中的位置生成三维锚点——纹理越复杂的地方越可能被选中。(3) \emph{三维局部属性预测}：对每个锚点查询其在三维空间中的 $K$-近邻，仅基于邻域特征回归完整高斯属性，将预测器的感受野对齐到 3DGS 优化的局部性本质。

\textbf{自适应采样}：直觉很简单——纹理丰富的区域需要更多高斯来拟合细节，平坦区域只需要少量大型高斯。SparseSplat 用 Shannon 熵来量化这种"需要多少高斯"的直觉。具体而言，对每张输入图像，在 $N \times N$（$N{=}7$）局部窗口内统计灰度直方图，计算每个像素处的 Shannon 熵，生成信息密度图 $E \in \mathbb{R}^{H \times W}$：

\begin{equation}
E(u,v) = -\sum_{i=0}^{L-1} p_i \log p_i
\end{equation}

其中 $L$ 为灰度级数（如 256），$p_i$ 是窗口内灰度级 $i$ 的归一化频率。高熵意味着窗口内灰度值变化剧烈（高频纹理），低熵意味着灰度均匀（平坦区域）。将熵图归一化后乘以温度系数 $\tau$，得到逐像素采样概率：

\begin{equation}
P_\tau(u,v) = \text{clip}\!\left(\frac{\tau \cdot E(u,v)}{\log L},\; 0,\; 1\right)
\end{equation}

即"熵越高，该像素被采样的概率越大"。$\tau$ 是控制全局稀疏度的旋钮：值越高所有位置被采样的概率都增大，点云越密；值越低则越稀疏。对每个像素独立掷一个 $[0,1]$ 均匀随机数，若小于 $P_\tau(u,v)$ 则选中该像素，再利用预测深度图和相机参数将其反投影到三维空间，构成稀疏锚点点云 $\mathcal{P}$。

\textbf{三维局部属性预测}：解决了"在哪里放高斯"后，SparseSplat 用一个严格遵循局部性原则的轻量预测器来回答"高斯属性是什么"。对每个锚点 $p_i \in \mathcal{P}$，通过 FAISS 在三维空间中检索 $K$-近邻 $\{p_{n_1}, \ldots, p_{n_K}\}$，并提取对应的二维特征。这里的关键设计是双投影策略：几何特征 $g \in \mathbb{R}^{d_g}$（三维坐标、法向量、视线方向）和图像特征 $v \in \mathbb{R}^{d_v}$ 分别通过可学习投影函数 $\phi_g$ 和 $\phi_v$ 映射到统一隐空间后拼接，为每个点构建融合几何与外观的表示 $f = [\phi_g(g);\, \phi_v(v)] \in \mathbb{R}^{2d_h}$。聚合中心点与 $K$ 个邻居的拼接特征后，通过 geo-aware attention 加权聚合（让网络自适应地关注最重要的邻居），再经轻量 MLP 回归完整高斯属性：不透明度 $\alpha \in \mathbb{R}$、尺度 $\mathbf{s} \in \mathbb{R}^3$、旋转四元数 $\mathbf{q} \in \mathbb{R}^4$、球谐系数 $\mathbf{c} \in \mathbb{R}^{N_{\text{sh}} \times 3}$。这一设计的核心优势在于：预测器的感受野（三维局部邻域）与任务的内在属性（局部优化）天然对齐——高斯的属性只取决于自身及其空间近邻，不需要整张图像的全局信息。

\textbf{设计动机回溯}：两个错配的分析直接催生了对应的设计选择。分布错配→基于熵的概率采样：用信息论的经典工具量化局部复杂度，使高斯密度与场景信息量成正比。感受野错配→三维局部预测器：回到 3DGS 优化的局部性本质，用 KNN 邻域取代全局上下文。值得注意的是，尽管使用了相同的冻结深度骨干，SparseSplat 在深度估计存在误差的场景中（如天空-建筑边界）反而比 DepthSplat 更鲁棒——论文将此归因于三维局部预测器能够从邻域点获取校正信息。

\textbf{训练与关键参数}：端到端仅用 RGB 监督，损失函数为 MSE 与 LPIPS 的加权组合：$L = (1-\lambda)L_{\text{MSE}} + \lambda L_{\text{LPIPS}}$。温度系数 $\tau$ 是核心超参数，单个训练模型通过调节 $\tau$ 即可在不同稀疏度之间切换，无需重训练。在 DL3DV-140 测试集上：$\tau$ 对应 150k 高斯时（DepthSplat 使用 688k 高斯，即 SparseSplat 仅用约 22\%），PSNR 达 24.20，与 DepthSplat 的 24.17 持平；极端压缩至 10k 高斯（约 1.5\%）时，PSNR 仍有 21.29，推理速度（105ms）甚至快于 DepthSplat（128ms）。

\paragraph{F4Splat: 致密化分数引导的可预算分配}

SparseSplat 通过温度系数 $\tau$ 隐式控制密度，但无法直接指定目标基元数量。F4Splat（2026.03）~\cite{2603_21304}沿续了自适应密度的思路，但引入致密化分数引导的分配策略（densification-score-guided allocation），将密度分配建模为可学习的逐区域分数预测问题，从而实现显式的基元数控制。具体而言，F4Splat 在标准像素对齐前馈管线的参数预测头中附加一个致密化分数头，为每个空间区域预测一个标量分数，表示该区域所需的高斯密度——分数越高，该区域被分配更多高斯基元。这一设计与 SparseSplat 的关键差异在于密度控制方式：SparseSplat 通过温度系数 $\tau$ 概率性地采样像素，是隐式的密度调节；F4Splat 则显式预测密度需求，并允许用户在推理时指定目标高斯总预算 $K$，模型据此自适应缩放各区域的致密化分数，无需重训练即可在不同基元数量间切换。在实现上，F4Splat 在推理时根据致密化分数对候选位置进行排序，保留分数最高的 $K$ 个位置生成高斯，同时根据空间复杂度和多视角重叠度自适应分配——简单区域和高度重叠视角间分配较少基元，几何复杂且重叠度低的区域获得更多基元。实验表明，F4Splat 在无需标定位姿的设定下，以显著更少的高斯基元实现了优于先前前馈方法的视角合成质量。

\paragraph{TokenGS: 可学习 token 解耦高斯与输入图像}

前述自适应密度方法虽然缓解了冗余问题，但高斯位置仍受输入像素或射线约束。TokenGS（2026.04）~\cite{2604_15239}完全脱离了像素对齐这一路线——无论是像素对齐的 MVSGaussian/DepthSplat（\S\ref{sec:13_1}），还是自适应密度的 SparseSplat（\S\ref{sec:13_1}），TokenGS 都抛弃了"高斯位置由输入像素/射线决定"这一根本假设，使高斯的位置和数量都与输入图像解耦。

\textbf{基线与问题}：像素对齐前馈 3DGS 存在三个根本性限制。第一，\emph{射线约束}——将高斯均值参数化为沿相机射线的深度值，使模型无法内部修正噪声位姿和多视角不一致性，且难以扩展到动态场景（需要随时间 warp 点/像素）。第二，\emph{数量耦合}——高斯数量与输入分辨率和视角数绑定，例如 32 张 $512 \times 512$ 图像产生超过 800 万个粒子，且将同一输入重复 $N$ 次就产生 $N$ 倍粒子，无论场景本身复杂度如何。第三，\emph{测试时微调破坏先验}——直接优化高斯参数容易退化网络学到的先验，尤其在少视角场景中。

\textbf{核心改动}：TokenGS 针对上述三个限制做了对应的三个关键改变。(1)~\emph{解耦射线}：不再沿相机射线预测深度，而是直接在全局坐标系中回归三维均值坐标 $\boldsymbol{\mu}$，如同前馈 SfM 模型预测点图——但仅用自监督渲染损失训练，无需点图真值监督。(2)~\emph{解耦数量}：高斯数量成为独立超参数 $N_t$，完全脱离输入分辨率和视角数的约束——每个可学习 token 解码 64 个高斯基元，因此总高斯数为 $N_t \times 64$。(3)~\emph{保护先验}：测试时仅微调 token 嵌入而非高斯参数本身，网络权重和图像特征保持冻结，从而在适应特定场景的同时保留学到的全局先验。

\textbf{实现原理——编码器}：编码器采用 ViT 架构。每张输入视角被分块化为 $p \times p$ 的 patch 序列，分别从图像 patch 和对应的 Pl\"ucker 坐标提取嵌入后求和，展平为编码器输入 token 序列 $A \in \mathbb{R}^{N_I \times C}$。这里 Pl\"ucker 坐标是射线的一种六维参数化——每条射线用其方向向量和矩向量联合表示，能同时编码射线的方向和空间位置，为网络提供比单纯相机位姿更丰富的射线几何信息。序列经标准 ViT 层（允许跨视角注意力）和层归一化后，输出图像 token $B \in \mathbb{R}^{N_I \times C}$。

\textbf{实现原理——解码器}：解码器借鉴 DETR（端到端目标检测器）的设计思想——正如 DETR 用一组可学习的「目标查询」通过交叉注意力从图像特征中提取检测框，TokenGS 用一组可学习的「3DGS token」从图像 token 中提取高斯参数。具体而言，初始化 $N_t$ 个可学习 token 嵌入 $T_{\text{IN}} \in \mathbb{R}^{N_t \times C}$，经 Transformer 解码器块处理——每块包含对图像 token $B$ 的交叉注意力、3DGS token 之间的自注意力、以及逐 token MLP——产生输出嵌入 $T_{\text{OUT}}$。每个输出嵌入通过线性层回归 64 个高斯基元的 14 维属性。

\textbf{实现原理——属性映射与 KV 共享}：属性映射遵循物理约束——XYZ 坐标使用 $f(x) = \sigma(x)(\text{sigmoid}(x) - 1)$（其中 $\sigma$ 为 SiLU 激活），该函数的关键特性是输出范围涵盖负值和正值且梯度信号均匀，适合全局坐标回归；颜色和不透明度用缩放 tanh 映射至 $[0,1]$，尺度用截断指数保证正值，旋转用单位归一化约束为有效四元数。此外，由于图像 token 数远多于 GS token（$N_I \gg N_t$），所有解码器交叉注意力层共享图像 token 的键-值投影，仅计算一次，将 $O(N_I \cdot D_{\text{dec}})$ 的内存分配降为单次 $O(N_I)$。默认架构为 6 层编码器、24 层解码器、通道维度 $C = 1024$。

\textbf{设计动机}：直接回归三维坐标的核心动机是解耦高斯位置与相机射线。这带来三重收益：(i)~支持外推和场景补全——高斯可被分配到所有输入视角之外的区域，实验表明在视角外推基准上 TokenGS 无需深度标注即显著超越基线；(ii)~对噪声位姿更鲁棒——坐标预测不受射线方向噪声的直接干扰，注入最多 10 度旋转噪声时 TokenGS 相对 GS-LRM 的优势随噪声幅度单调增大；(iii)~消除深度预测网络中常见的尖刺伪影（spiky artifacts）——因为高斯位置不再受限于射线。

\textbf{零梯度问题与可见性损失}：直接坐标回归引入了一个新的训练挑战——位于所有相机视锥之外的高斯不贡献渲染监督，梯度为零。这些"非活跃"高斯退化训练稳定性、浪费模型容量，并表现为场景周围的漂浮噪点。TokenGS 不采用显式三维监督（会重新引入射线偏差），而是提出可见性损失（visibility loss）——一种柔性的软约束，通过将高斯质心投影到所有监督视角的图像平面，惩罚那些落在所有视角可见范围之外的高斯。具体而言，将第 $m$ 个高斯的质心投影到第 $i$ 个监督视角的图像平面并归一化至 $[-1, 1]$，得到 $(\tilde{u}_m^i, \tilde{v}_m^i)$，然后定义：

\begin{equation}
\mathcal{L}_{\text{vis}} = \sum_m \min_{I_i \in \mathcal{I}_{\text{sup}}} \max\!\left(0,\; 0.5 - \max(|\tilde{u}_m^i|, |\tilde{v}_m^i|)\right)
\end{equation}

直觉上，该公式计算每个高斯到最近可见视角边界的"距离"——当高斯投影在至少一个视角内部（$|\tilde{u}|, |\tilde{v}| < 0.5$）时损失为零，否则产生梯度将其拉回可见区域。实践中损失值被裁剪至 1.0 以避免虚假梯度。消融实验表明可见性损失带来 0.4 dB 的 PSNR 提升。总训练损失为 $L = L_{\text{render}} + \lambda_1 L_{\text{SSIM}} + \lambda_2 L_{\text{vis}}$，其中 $\lambda_1 = 0.2$，$\lambda_2 = 1.0$。

\paragraph{GlobalSplat: 显式对齐的全局 token 场景表示}

TokenGS 通过可学习 token 直接回归全局三维坐标，实现了位置与输入的解耦。GlobalSplat（2026.04）~\cite{2604_15284}遵循与 TokenGS 不同的"先对齐，后解码"（align first, decode later）原则，通过双分支迭代注意力架构在潜在空间中显式构建全局场景表示，再从中解码三维几何。TokenGS 通过可学习 token 直接回归全局三维坐标，而解码前的对齐仍依赖 Transformer 的隐式学习。GlobalSplat 则将"先对齐"显式化为一个独立的潜在场景表示构建阶段——先用双分支迭代注意力架构（显式解耦几何分支和外观分支，各自拥有独立的查询 $Q_G$、$Q_A$）将所有输入视角的特征融合到一组固定数量的全局潜在场景 token 中，跨视角对应关系在这一阶段通过交叉注意力被显式解决。只有在全局场景表示构建完成后，才开始解码显式三维几何。

这一设计的关键优势在于场景 token 的数量与输入视角数完全无关——无论输入 8 张还是 36 张图像，最终都输出相同数量的全局 token（默认 2K--32K），由此解码的高斯基元数量也恒定。为防止表示膨胀，GlobalSplat 采用粗到细的训练课程（coarse-to-fine capacity curriculum）：训练初期仅允许少量 token 被解码为高斯，随着训练推进逐步增加解码容量，使模型自然学会在场景复杂区域集中分配基元。在 RealEstate10K 和 ACID 数据集上，GlobalSplat 仅用 16K 高斯（约 4MB 存储）即可实现有竞争力的渲染质量（24 视角下 28.5 dB PSNR），推理仅需 78 毫秒，峰值 GPU 内存仅 1.79 GB——相比密集像素对齐管线减少超过 99\% 的基元数量。

\paragraph{AnchorSplat: 几何锚点驱动的锚定高斯表示}

与 GlobalSplat 的全局 token 表示不同，AnchorSplat（2026.04）~\cite{2604_07053}提出锚点对齐的高斯表示，与 TokenGS 的可学习 token 方案不同，AnchorSplat 以外部提供的三维几何先验（稀疏点云、体素或 RGB-D 点云）作为锚点，直接在三维空间中构建高斯表示。具体而言，AnchorSplat 将多视角图像特征反投影到锚点位置形成锚定特征体（anchored feature volume），再由高斯解码器从中预测高斯参数——这一流程使高斯表示完全独立于图像分辨率和视角数量。相比 TokenGS 完全依赖网络学习 token 的三维位置，AnchorSplat 利用几何先验提供了更可靠的空间锚定，减少了浮游体和跨视角不一致。此外，AnchorSplat 引入即插即用的 Gaussian Refiner 模块，通过数次前向传播精调中间高斯的属性（位置、尺度、颜色），提升几何一致性而无需重训练整个模型。在 ScanNet++ v2 基准上，AnchorSplat 以约 20 倍少于像素对齐方法（如 AnySplat）的高斯数量实现了更优的渲染质量，重建时间约为前者的一半。

\paragraph{VolSplat: 体素对齐的三维自适应密度控制}

前述方法从 token 和锚点角度脱离像素对齐。VolSplat（2025.09）~\cite{2509_19297}用体素对齐取代像素对齐——从预测的三维体素网格直接解码高斯参数，避免了易错的二维特征匹配，实现了基于三维场景复杂度的自适应密度控制。与 GlobalSplat 在潜在空间中用 token 表示全局场景不同，VolSplat 的体素表示保持了显式的三维空间结构。其管线分为三个阶段：(1) 从多视角图像提取二维特征并构建逐视角代价体，通过深度预测模块估计各视角深度图；(2) 利用深度图将二维特征解投影至三维空间，构建体素特征网格，再由稀疏 3D U-Net 在体素空间中精炼特征；(3) 从精炼后的体素特征直接预测每个占据体素的高斯参数。这一设计解决了像素对齐的两个核心局限：其一，像素级采样在遮挡、弱纹理和运动视差区域产生不一致的采样模式，体素聚合则将多视角信息融合到统一的 3D 容器中，消除了视角偏差的密度分布；其二，像素对齐强制均匀密度，体素对齐则允许根据体素内三维场景的复杂度自适应分配高斯密度——场景结构复杂的体素获得更多基元，平坦区域则更稀疏。在 RealEstate10K、ScanNet 和 ACID 上的实验证实了 VolSplat 在几何一致性和渲染质量上的一致优势。

\paragraph{PointSplat: 三维几何驱动的后处理剪枝}

上述三种方法（GlobalSplat、AnchorSplat、VolSplat）分别从全局 token、几何锚点和体素空间三个维度突破了像素对齐的表示瓶颈，但都聚焦于前馈重建阶段。PointSplat（2026.04）~\cite{2604_09903}则从另一个角度切入——结合三维几何驱动的剪枝和 Transformer 精炼，仅依赖高斯基元的三维属性进行重要性排序，无需二维图像参与剪枝阶段。与 SparseSplat 在前馈管线中自适应采样不同，PointSplat 面向的是后处理场景——给定一个已生成的过参数化 3DGS 模型，如何高效压缩。其核心贡献有两个：(1) 几何驱动的剪枝分数——仅基于高斯的不透明度和体积（三维空间覆盖）计算重要性分数，平衡空间覆盖与信息密度，避免了传统方法需要光线-像素交点计算的昂贵开销；(2) 双分支编码器（Dual-Branch Encoder）——将几何属性（位置、尺度、朝向）与高维外观特征（球谐系数）显式分离，通过独立分支分别处理并重新加权，避免高维外观特征对低维几何属性的支配（feature imbalance）。剪枝后的高斯经 Transformer 精炼模块全局协调参数，修复因剪枝引入的模糊纹理和断裂几何边缘。在 ScanNet++ 和 Replica 上的实验表明，PointSplat 在多种稀疏度级别下一致实现有竞争力的渲染质量，且无需逐场景优化。

\paragraph{EcoSplat: 基元数量可控的前馈框架}

前馈 3DGS 产生的高斯基元数量随输入分辨率和视角数线性增长，如何在不损失渲染质量的前提下实现效率可控的紧凑表示，是一个重要的工程和研究问题。EcoSplat（2025.12）~\cite{2512_18692}是首个效率可控的前馈 3DGS 框架，其核心目标是让模型在推理时能根据任意指定的目标基元数量 $K$ 自适应调整输出，同时最大化渲染质量。与 SparseSplat 通过温度系数间接控制密度不同，EcoSplat 提供显式的基元数控制。其技术路线为两阶段训练：第一阶段像素对齐高斯训练（Pixel-aligned Gaussian Training, PGT），学习基础的像素对齐基元预测能力；第二阶段重要性感知微调（Importance-aware Gaussian Finetuning, IGF），让模型学会为每个基元排序并自适应调整参数。

IGF 阶段的关键设计是将目标基元数 $K$ 编码为可学习重要性嵌入注入参数预测头，同时基于光度梯度和几何法线变化率生成伪真值重要性掩码——高变化区域（纹理丰富、几何边界）的高斯被完整保留，低变化区域则通过单步 K-means 聚类合并冗余。重要性感知不透明度损失（$\mathcal{L}_{\text{io}}$，基于二元交叉熵）引导模型自适应抑制不重要基元的不透明度。推理时仅需保留不透明度最高的 $K$ 个高斯，并可按各视角高频分量比例分配逐视角基元配额。训练中还采用渐进式学习策略（Progressive Learning on Gaussian Compaction, PLGC），逐步扩大目标基元数的采样范围，提升训练稳定性。在 RealEstate10K 和 ACID 的稠密视角设置下，EcoSplat 在严格的基元数约束下一致超越先前方法。

\paragraph{TinySplat: 三重冗余消除的存储压缩}

EcoSplat 聚焦于基元数量控制，而前馈 3DGS 输出的存储压缩同样至关重要。TinySplat（2025.06）~\cite{2506_09479}提出完整的前馈压缩框架，关注的是存储压缩——如何在保持渲染质量的前提下将前馈 3DGS 的数据体积压缩到极致。其核心洞察是前馈 3DGS 生成的像素对齐高斯特征图中蕴含三类不同性质的冗余，可分别消除：(1) 几何冗余——通过视角投影变换（View-Projection Transformation, VPT）消除：VPT 将高斯几何参数从世界坐标系变换到相机空间，利用像素对齐高斯的空间正则性仅编码相对于像素中心的偏移量，并近似地将旋转和缩放参数变换至相机坐标系以增强局部相关性，使几何参数在新的坐标空间中呈现更紧凑的分布。(2) 感知冗余——通过可见性感知基约简（Visibility-Aware Basis Reduction, VABR）消除：VABR 根据各球谐基函数在有效观测方向上的平均响应值计算可见性权重，通过加权协方差矩阵的特征分解提取主成分基，将高维 SH 系数投影到紧凑的可见性感知子空间中——仅保留对渲染贡献最大的方向分量。(3) 空间冗余——通过现成视频编解码器处理。整个压缩框架无需训练，直接应用于前馈 3DGS 的输出。在多个基准数据集上，TinySplat 实现超过 100 倍的高斯数据压缩，以仅 6\% 的存储大小达到与最佳压缩方法相当的渲染质量，编码时间仅需 25\%，解码时间仅需 1\%。

\paragraph{CodecSplat: 管线内特征级端到端压缩}

EcoSplat~和~TinySplat~均在后处理阶段进行压缩或基元选择。CodecSplat（2026.05）~\cite{2605_25563}与前两者采用后处理压缩的策略根本不同——CodecSplat~将压缩整合进前馈高斯生成管线内部，在二维特征到三维高斯的映射之前就施加压缩。其核心洞察是：前馈网络在将中间特征映射为高斯基元之前，已经将场景信息组织为规则、视角对齐且几何感知的二维特征网格——这一结构化表示比最终不规则的~3D~高斯基元更适合压缩。因此，CodecSplat~在编码端对中间二维高斯生成特征进行熵编码生成场景比特流；在解码端，从比特流重建潜在特征，再通过轻量级深度预测头和高斯属性预测头映射为三维高斯基元。这一管线内压缩的关键优势在于：压缩与高斯生成共享同一个紧凑的中间表示，避免了先转换为不规则~3D~高斯再压缩的效率损失。在实现上，CodecSplat~采用深度引导的多视角特征精炼模块和分层学习特征编解码器，支持可调的率失真行为。在~DL3DV~数据集上，CodecSplat~用仅~20--108~KiB~的数据实现~23.56--26.36~dB~PSNR；在~RealEstate10K~上，仅需~3.37--12.51~KiB~即可实现~24.76--27.05~dB~PSNR——数据量比对前馈生成的~3D~高斯直接压缩小约一个数量级。

\subsubsection{几何增强与深度估计改进}

本路线介绍从前馈管线中深度估计环节进行几何增强的改进方法。\S\ref{sec:13_1}~节前述方法已确立了"特征提取$\to$多视角交互$\to$深度估计$\to$高斯参数解码"的基础管线。然而，深度估计作为整个管线中连接二维特征与三维位置的关键桥梁，其精度直接决定了高斯基元的几何质量。本节综述的方法不改变管线整体结构，而是聚焦于深度估计这一核心环节，从不同角度对其进行改进：迭代深度概率精炼、单目先验与多视角特征的深度融合、Nyquist 采样率约束下的几何属性修正、以及深度表示本身的重新审视。

\paragraph{级联多轮 warp 的深度概率迭代提升}

DepthSplat 等方法通过单次平面扫掠（plane-sweep stereo）构建代价体并经 softmax 归一化获得深度概率分布，整个过程只执行一次 warp 操作——即将源视角特征投影到目标视角的一系列深度平面上计算匹配得分。IDESplat（2026.01）~\cite{2601_03824}观察到这种单次 warp 存在固有的不稳定性：一次投影只能捕获"当前深度候选下特征匹配得怎样"，但无法充分利用多视角几何中固有的极线约束（epipolar constraint，同一三维点在不同视角中的投影必须落在对应极线上的几何关系）信号——因为匹配质量的判断需要多个深度的交叉验证，而单次 warp 的视野受限。

针对这一瓶颈，IDESplat 提出深度概率提升单元（Depth Probability Boosting Unit, DPBU）。其核心操作是\emph{级联多轮 warp，以乘法方式融合极线注意力图}。具体而言，第一轮 warp 产生初始极线注意力图 $A_1$（编码源-目标视角在各深度处的匹配关系），第二轮 warp 以 $A_1$ 引导的更新特征为输入产生 $A_2$，两轮注意力图按元素相乘：$A_{\text{fused}} = A_1 \odot A_2$。直觉上，乘法融合意味着"两轮都认为匹配好的深度才获得高概率"——交叉验证的效果。DPBU 可堆叠多个形成迭代深度估计过程：每个 DPBU 在前一轮输出的深度概率分布基础上，重新采样深度候选并执行新一轮 warp，逐步淘汰低概率深度候选，集中概率质量于高似然区域。随着迭代推进，深度估计从粗粒度的多峰分布逐渐收敛为尖锐的单峰分布，高斯中心位置的预测精度随之提高。

在效率方面，IDESplat 的参数量仅为 DepthSplat 的 10.7\%、内存开销仅 70\%，却在 RealEstate10K 上实现 0.33 dB PSNR 提升；在跨数据集实验中，于 DTU 上相对 DepthSplat 提升 2.95 dB PSNR，展现了强大的泛化能力。这一结果表明，在深度估计环节投入迭代计算比简单地增大骨干网络参数量更有效——迭代地精炼深度概率比一次性预测整个深度图更适合匹配问题的固有结构。

\paragraph{单目-多视角特征轻量适配器}

与 DepthSplat 直接冻结预训练单目深度骨干（Depth Anything V2 的 ViT）提取单目特征的策略不同，MonoSplat（2025.05）~\cite{2505_15185}提出了一种更灵活的单目先验利用方式：不冻结骨干，而是设计轻量级的单目-多视角特征适配器（Mono-Multi Feature Adapter），以最少的可训练参数将单目特征转换为多视角表示。这一设计的核心优势在于效率：适配器的可训练参数远小于从头训练整个骨干，却在跨场景泛化上表现卓越。

其动机来自一个实际观察：预训练单目深度模型在大规模数据上学到的视觉先验极为丰富（能处理各种光照、纹理、遮挡条件），但直接冻结使用时这些先验被"锁死"在单目预测范式中——模型只见过从单张图像预测深度的任务，不天然具备多视角几何推理能力。MonoSplat 的适配器通过轻量注意力机制实现两个目标：(1) 跨视角对齐——将不同视角的单目特征映射到统一的多视角特征空间中；(2) 先验保持——在适配过程中保留有价值的单目几何先验不被破坏。适配后的多视角特征与原始单目特征在集成高斯预测模块（Integrated Gaussian Prediction）中融合，共同解码高斯基元的几何和外观参数。这一设计的核心优势在于效率：适配器的可训练参数远小于从头训练整个骨干，却在跨场景泛化上表现卓越——因为丰富的单目先验被有效迁移到了多视角推理任务中。

\paragraph{MVS 与单目深度的投影-采样融合}

前述方法（DepthSplat、IDESplat）主要针对固定的视角间距（基线）设计深度估计策略，但当输入视角的基线差异较大时——有些视角相距很近（小基线，匹配容易但三角化精度低），有些相距很远（大基线，三角化精度高但匹配困难）——单一策略难以同时应对。MuGS（2025.08）~\cite{2508_04297}（Multi-Baseline Gaussian Splatting）正是针对这一问题，提出融合多视角立体视觉（MVS）和单目深度估计（MDE）特征的统一框架。

MuGS 的核心贡献是投影-采样机制（projection-and-sampling mechanism）用于深度深度融合。直觉上，MVS 代价体提供多视角一致的几何约束（在纹理丰富、基线适中的区域精度高），而单目深度估计提供全局一致的相对深度先验（在弱纹理和大基线区域更可靠）。投影-采样机制的工作方式为：将单目深度估计的相对深度投影到 MVS 的深度候选空间中，用单目深度引导深度候选的采样密度——在单目模型高置信的区域多采样以获取精度，在低置信区域少采样以避免误匹配——由此构造一个精细概率体（fine probability volume），引导特征图的深度回归。此外，MuGS 引入参考视角损失（reference-view loss），以参考视角的渲染质量作为额外的几何约束信号，提升优化效率。

MuGS 在从简单物体（DTU）到复杂室内外场景（RealEstate10K）的多种基线设置下均达到最佳性能，并在 LLFF 和 Mip-NeRF 360 上展现了有前景的零样本泛化能力。

\paragraph{采样率引导的低通滤波几何修正}

SurfelSplat（2026.04）~\cite{2604_08370}从信号处理的基本原理出发，为前馈 3DGS 的几何精度问题提供了一个新视角。作者观察到，像素对齐高斯基元的空间频率天然超过 Nyquist 采样率——即高斯基元在图像平面上的投影密度过高，相邻基元之间的空间间隔不足以准确恢复每个基元的几何属性（法向量、曲率等）。这导致虽然渲染的 RGB 图像看起来合理，但提取的几何表面却充满噪声和断裂。

为解决这一问题，SurfelSplat 提出基于采样率引导的低通滤波策略：首先根据输入分辨率和基元密度计算局部 Nyquist 频率，然后对高斯基元的几何属性施加空间低通滤波，滤除超出采样率的高频分量。滤波后的基元被投影到所有输入视角以获取跨视角特征相关性，通过专门设计的特征融合网络回归具有精确几何的高斯 surfel 表示。实验表明，SurfelSplat 在 DTU 重建基准上与最先进的优化方法取得可比结果，而预测仅需 1 秒——相比逐场景优化实现 100 倍加速。

\paragraph{2D 高斯 surfel 与表面连续性约束}

上述方法均基于 3D 高斯基元进行重建。SurfSplat（2026.02）~\cite{2602_02000}则提出了一个根本性的表示变更：用 2D 高斯（Gaussian surfel）替代 3D 高斯作为前馈重建的基元。其动机来自一个关键观察——3D 高斯基元本质上是各向同性的椭球体，虽然表达力强，但在表示薄壁表面（如墙壁、桌面）时存在固有缺陷：一个 3D 高斯需要在表面两侧都分配概率质量，导致重建结果在放大观察时呈现离散的、颜色偏移的点云，而非连续的几何表面。2D 高斯（surfel）则天然是扁平的表面元素，提供更强的各向异性和更高的几何精度。

然而，将 2D 高斯引入前馈管线需要解决两个新问题。第一，\emph{表面连续性}——独立的 2D 高斯可能在表面边界处产生间隙。SurfSplat 引入表面连续性先验，约束相邻 surfel 之间的法向量一致性，确保重建的几何表面连贯。第二，\emph{强制 $\alpha$-混合}——标准 3DGS 的 $\alpha$-混合允许透明度累积产生半透明效果，但这在几何重建中会引入模糊。SurfSplat 强制执行不透明度阈值，使每个 surfel 的贡献更加明确，从而重建出更清晰的表面边界。

此外，SurfSplat 引入高分辨率渲染一致性（High-Resolution Rendering Consistency, HRRC）这一新评估指标，专门设计用于评估高分辨率下的重建质量——标准指标在低分辨率下可能掩盖几何缺陷，但放大后这些缺陷变得明显。在 RealEstate10K、DL3DV 和 ScanNet 上，SurfSplat 在标准指标和 HRRC 上均一致优于先前方法。

\paragraph{Pointmap 正则化抑制深度边界碎片}

Shi et al.（2025.06）~\cite{2506_05327}回到了一个更基本的问题：深度图本身是否是前馈 3DGS 中最合适的几何表示？深度图通过 unprojection 将每个像素映射为三维空间中的一个点，过程简单高效。然而，作者指出深度图存在一个被忽视的结构性缺陷——\emph{物体边界处的深度间断}。当物体与背景之间存在深度跳变时（如人站在墙前），物体边缘处的像素从前景深度突然切换到背景深度，导致 unprojection 后的点云在边界处产生碎片化的"飞溅"——一些点飘浮在前景和背景之间，破坏了几何表面的连贯性。

为解决这一问题，作者提出 PM-Loss，一种基于预训练 Transformer（如 DUSt3R）预测的 pointmap 的几何平滑性正则化损失。Pointmap 是一种替代深度的三维表示——为每个像素预测一个三维坐标而非标量深度值。关键洞察在于：虽然 pointmap 本身的逐点精度通常不如多视角深度估计（因为单目或双目基础模型的几何精度有限），但 pointmap 在物体边界处具有天然的结构连续性——预训练 Transformer 在大规模数据上学到了"物体是连贯的整体"这一先验，因此即使在深度发生跳变的边界处，pointmap 的预测也不会产生碎片化的飞溅。PM-Loss 利用这一特性作为正则化信号：在深度图 unprojection 为点云后，用 pointmap 的局部几何结构约束点云在边界处的平滑性，抑制碎片化效应。

值得注意的是，PM-Loss 是一种即插即用的正则化损失，可无缝集成到各种前馈 3DGS 架构中（如 DepthSplat、MVSGaussian 等），在无需改变网络结构的前提下显著提升渲染质量。

\subsubsection{迭代精炼与反馈机制}

本路线介绍通过迭代精炼与反馈机制突破单次前向传播容量限制的方法。单次前向传播的容量有限，多个工作探索了迭代精炼策略来突破这一限制。

\paragraph{渲染误差驱动的循环迭代精炼}

ReSplat（2025.10）~\cite{2510_08575}修改了 DepthSplat 的单步预测范式——用循环迭代精炼替代一次性推断——渲染和损失函数保持不变。直觉上，这类似于画家先快速铺出画面整体色调，再逐轮添加细节：每轮用"渲染得对不对"作为反馈，指引下一轮该改哪里。

\textbf{Baseline 与动机。}DepthSplat 等单步前馈方法受限于网络容量：预测一次就输出最终结果，无法像逐场景优化那样反复迭代改进。然而逐场景优化太慢，两全其美的思路是——用前馈网络模拟迭代的"多轮修正"，但整个过程仍然是梯度无关的前馈推断，不涉及测试时梯度回传。

\textbf{紧凑初始化。}ReSplat 以 DepthSplat 架构为骨干，但在 $16\times$ 下采样空间（$4\times$ 空间压缩的深度图）预测初始高斯——仅产生 $M = N \times HW/16$ 个基元，是像素对齐方法的 $1/16$。为弥补信息损失，引入六层交替的 $k$NN 注意力（$k=16$，捕捉三维局部结构）和全局注意力（建立远距离依赖），将点云特征 $f_j \in \mathbb{R}^{C_0}$（$C_0=256$ 小模型 / $512$ 基础模型）增强为上下文感知特征 $f_j^*$，再由两层 MLP 解码高斯参数。每个高斯的完整参数 $g_j^0 \in \mathbb{R}^{C_1}$（$C_1=59$，包含位置、协方差、不透明度、球谐系数）与隐藏状态 $z_j^0 = f_j^*$ 构成循环网络的初始输入。

\textbf{渲染误差作为反馈信号。}每次迭代中，先用当前高斯渲染输入视角，再比较渲染图像与真实图像的差异。这一误差融合了像素空间和特征空间两个视角——像素误差经线性投影加 LayerNorm 对齐通道维度后，与预训练 ResNet-18 三阶段特征（$1/2$、$1/4$、$1/8$ 分辨率，统一下采样至 $1/4$）的差异相加：

\begin{equation}
E^t = \{ \hat{F}_i^t - F_i \}_{i=1}^N + \mathrm{proj}(\{ \hat{I}_i^t - I_i \}_{i=1}^N),
\end{equation}

其中 $C_2=256$ 为特征维度——直觉上，$E^t$ 同时回答了两个问题："颜色对不对"和"纹理细节对不对"。

\textbf{误差传播与高斯更新。}像素对齐的误差 $e_j^t$ 只能告诉每个高斯"你正下方的像素画得怎样"，但渲染是一对多的——一个高斯影响周围多个像素。为此 ReSplat 对所有 $N \times HW/16$ 个渲染误差执行全局注意力，使每个高斯接收来自全部渲染像素的信息：

\begin{equation}
\tilde{E}^t = \mathrm{GlobalAttn}(E^t) = \{\tilde{e}_j^t\}_{j=1}^{N \times HW/16}.
\end{equation}

更新模块由四个 $k$NN 注意力块组成——选择 $k$NN 而非全局注意力是因为高斯更新本质上是局部操作（移动位置、调整大小）——接收拼接的高斯参数 $g_j^t$、隐藏状态 $z_j^t$ 和全局聚合的渲染误差 $\tilde{e}_j^t$，经四层 MLP 预测增量更新：

\begin{equation}
g_j^{t+1} = g_j^t + \Delta g_j^t, \quad z_j^{t+1} = z_j^t + \Delta z_j^t.
\end{equation}

整个过程不涉及梯度回传——渲染误差只是网络的输入信号，不是损失函数的梯度。

\textbf{训练策略。}两阶段训练：第一阶段训练初始重建模型，损失为渲染损失加深度平滑正则（$\lambda=0.01$, $\beta=0.5$）——深度平滑项不需要真值深度，仅惩罚与图像梯度不一致的深度跳变；第二阶段冻结初始模型，仅训练循环更新模块，使用指数递增权重监督所有迭代的渲染输出：

\begin{equation}
\mathcal{L}_{2\text{nd}} = \sum_{t=0}^{T-1} \gamma^{T-1-t} \sum_{v=1}^{V} \mathcal{L}_{\text{render}}(\hat{I}_v^t, I_v), \quad \gamma=0.9,
\end{equation}

$\gamma=0.9$ 使后期迭代的监督权重逐步增大，引导模型从粗到细精炼。模型在 4 次迭代后收敛，推理时支持灵活调整迭代次数以权衡速度与精度。在 DL3DV、RealEstate10K 和 ACID 数据集上，ReSplat 以 $1/16$ 的高斯数量超越像素对齐方法，参数量仅 77M（vs 像素对齐方法 559M），速度比逐场景优化快 $100\times$。

\paragraph{冻结扩散先验引导的前馈迭代精炼}

ReSplat 的循环迭代完全依赖渲染误差作为反馈信号——"画得不对就改哪里"——这一策略在前馈框架内已十分有效，但本质上只利用了场景自身的观测信号。一个自然的追问是：能否在迭代过程中引入外部生成先验，为模型提供"应该画成什么样"的额外指导？GIFSplat（2026.02）~\cite{2602_22571}正是沿着这一思路，在 ReSplat 的纯前馈迭代框架之上引入了冻结的扩散先验，在不牺牲前馈效率的前提下实现场景自适应精炼。

GIFSplat 的架构由三个互相衔接的组件构成：初始化器、迭代高斯头和生成先验蒸馏模块。

\textbf{初始化。}GIFSplat 以 AnySplat（\S\ref{sec:13_2} 节中的无位姿前馈方法）架构为初始化器，从稀疏无位姿输入视角预测初始 3DGS。AnySplat 的选择有两重考量：其一，它同时预测相机参数和三维高斯，因此 GIFSplat 不需要已知位姿；其二，AnySplat 的可微分体素化模块将逐像素基元聚类到体素网格中，产生的初始高斯数量与场景复杂度而非像素数相关，为后续迭代提供了紧凑的起点。

\textbf{迭代高斯头。}迭代精炼的核心挑战在于信号粒度的转换——渲染误差天然是\emph{像素级}信号（每个像素的颜色差异），而需要更新的对象是\emph{高斯级}参数。GIFSplat 通过软赋予权重（soft assignment weights）实现这一转换：计算每个像素受哪些高斯影响及影响程度（即 splatting 过程的逆操作），据此将像素级渲染特征差异池化（pool）至高斯级别，构成\emph{观测线索}（observation cue）。直觉上，这回答了"哪些高斯的渲染结果与真实图像不一致，不一致到什么程度"。

\textbf{生成先验蒸馏。}GIFSplat 的核心创新在于引入冻结的扩散先验 Difix（一种预训练的三维场景扩散修复模型）作为第二个信息源。每次迭代中，Difix 对当前渲染的中间图像进行增强——补充缺失的细节、修正伪影——然后将增强前后的图像特征差异同样池化至高斯级别，构成\emph{生成先验线索}（generative prior cue）。关键设计在于，Difix 完全冻结，不参与训练也不参与梯度回传：它仅作为一个"先验提供者"，通过前向推理为迭代过程注入关于三维场景应有外观的知识。

\textbf{残差更新。}观测线索和生成先验线索在高斯级别拼接后，驱动窗口注意力模块（window attention，限制每个高斯仅关注局部邻域以保持效率）预测逐高斯残差更新。整个迭代过程为 $T$ 步纯前向更新，不涉及测试时梯度回传或逐步增长的视角集扩展。在 DL3DV 上 GIFSplat 实现了最多 2.1 dB PSNR 提升，在 RealEstate10K 和 DTU 上同样一致超越最佳前馈基线，同时保持秒级推理时间。

\paragraph{梯度反馈驱动的可泛化迭代重建}

G3R（2024.09）~\cite{2409_19405}从大规模城市场景重建的实际需求出发，提出了一种介于纯前馈推断和逐场景优化之间的折中策略。其核心思想是将逐场景优化的高保真度与数据驱动先验的快速推断相结合。其核心思想是：学习一个重建网络，利用可微渲染的\emph{梯度反馈信号}（而非像 ReSplat 那样利用渲染误差本身作为网络输入）来迭代更新三维场景表示。

具体而言，G3R 的迭代过程不同于 ReSplat 和 GIFSplat 的"前馈残差更新"范式——它直接将可微渲染的梯度信号用于更新三维表示。直觉上，ReSplat 用神经网络学习"如何根据渲染误差修改高斯"（数据驱动的更新策略），而 G3R 直接利用可微渲染的解析梯度来更新（优化驱动的更新策略），但整个更新过程由一个训练好的网络控制更新幅度和方向。这一设计将逐场景优化的高保真度（来自梯度信号）与数据驱动先验的快速推断（来自网络的泛化能力）相结合。在城市驾驶和无人机数据集上，G3R 实现了与 3DGS 逐场景优化可比甚至更优的渲染质量，同时加速 10 倍以上，且对大视角变化更加鲁棒。

Diff3R（2026.04）~\cite{2604_01030}与前述方法的不同在于：它不满足于让前馈网络直接输出最终高斯参数，也不局限于在推理时做若干步梯度精炼，而是将\emph{训练过程本身}设计为一个双层优化问题——内循环模拟测试时优化（TTO，Test-Time Optimization），外循环据此调整网络，使其学会为 TTO 输出最优的初始化。这一设计源于一个核心观察：前馈网络的单步输出 $\theta_0 = g_\theta(\{I_{\text{ctx}}, P_{\text{ctx}}\})$ 虽然能捕获场景的大致几何和外观，但天然缺少高频细节。而标准 TTO 在稀疏视角下又极易灾难性过拟合——因为没有足够的多视角约束，高斯可以轻易地"作弊"以完美重建输入视角，却严重破坏底层几何。Diff3R 的解决思路是让前馈预测和 TTO 协同训练，而非各自独立运作。

\textbf{双层优化框架。}
直观上，Diff3R 的训练过程可以理解为一种''模拟考试''：每次训练迭代中，网络先输出初始高斯 $\theta_0$，然后内循环对其执行若干步 3DGS 优化以拟合上下文视角，外循环再根据精炼后 $\theta^*$ 在新视角上的渲染质量来更新网络权重。内循环的损失函数为近端正则化的光度损失：

\begin{equation}
\mathcal{L}_{\text{in}}(\theta) = \frac{1}{2} \| r(\theta) \|_2^2 + \frac{1}{2} (\theta - \theta_0)^\top \Lambda (\theta - \theta_0),
\end{equation}

其中 $r(\theta) = f(\theta) - I_{\text{ctx}}$ 为渲染残差，$\Lambda$ 为正则化权重。近端项有两个关键作用：一是约束高斯不偏离网络先验过远，防止过拟合；二是确保收敛点具有关于 $\theta_0$ 的梯度，使外循环的反向传播成为可能。外循环则在新视角渲染损失 $\mathcal{L}_{\text{outer}}(\theta^*) = \mathcal{L}_{\text{photo}}(f(\theta^*), I_{\text{nv}}) + \lambda_\text{outer} \mathcal{L}_{\text{reg}}(f(\theta^*), I_{\text{nv}})$ 上更新网络权重 $\theta$。

\textbf{IFT 隐式梯度与 PCG 求解。}
双层优化的核心难点在于计算 $\partial \theta^* / \partial \theta_0$——即"网络输出的初始值发生了微小变化，经过内循环优化后最终参数会如何改变"。一个朴素的做法是展开所有优化步骤逐一求导，但这需要为每一步保留完整的计算图，内存代价极高。Diff3R 利用隐函数定理（Implicit Function Theorem, IFT）绕过这一困难：IFT 的核心思想是，如果在收敛点处内循环梯度为零（$\nabla_\theta \mathcal{L}_{\text{in}}(\theta^*) = 0$），则可以在这个最优性条件上直接求导得到解析梯度，无需回溯整个优化轨迹。对最优性条件关于 $\theta_0$ 求导可得：

\begin{equation}
\frac{\partial \theta^*}{\partial \theta_0} = (H_\Lambda + I)^{-1}, \quad H_\Lambda = J^\top J,
\end{equation}

其中 $J = \partial r(\theta^*) / \partial \theta$ 为渲染残差的雅可比矩阵，$H_\Lambda = J^\top J$ 是用 Gauss-Newton 近似的海森矩阵——即将二阶导数项近似为雅可比矩阵的乘积，从而避免计算昂贵的二阶渲染导数。然而，高斯参数量通常达百万级，直接实例化并求逆矩阵 $(J^\top J + I)$ 在计算上不可行。Diff3R 将其转化为线性系统 $(J^\top J + I)v = \nabla_{\theta^*} \mathcal{L}_{\text{outer}}$，利用预条件共轭梯度（Preconditioned Conjugate Gradient, PCG）算法求解。PCG 只需矩阵-向量乘积而无需显式构造矩阵，最终梯度即为 $\partial_{\theta_0} \mathcal{L}_{\text{outer}} = v$。值得注意的是，该线性系统在结构上与 Levenberg-Marquardt（LM，莱文贝格-马夸特）优化的法方程相同（$\Lambda$ 在 LM 中扮演阻尼因子的角色），因此可直接复用高效的 CUDA LM 求解器加速实现。正则化参数 $\Lambda$ 的两个极端值也提供了清晰的物理直觉：$\Lambda \to \infty$ 时梯度直接穿透内循环（网络学到的即为最终输出），$\Lambda \to 0$ 时梯度消失（高斯完全忘记初始化，学习信号被切断）。

\textbf{不确定性感知正则化。}
上述分析揭示了全局标量 $\Lambda$ 的固有局限：场景中不同区域的预测置信度差异巨大，一个全局值要么过度约束高置信区域、要么切断低置信区域的梯度流。Diff3R 为此引入逐高斯逐属性的正则化权重 $\Lambda \in \mathbb{R}^{N_\theta}$，由网络与 $\theta_0$ 同时预测——若网络对某高斯的位置或颜色不确定，则预测较小的惩罚值，允许该参数在 TTO 中自由调整；反之则强锚定。$\Lambda$ 的梯度同样通过 IFT 推导：$\partial_\Lambda \mathcal{L}_{\text{outer}} = -v \odot (\theta^* - \theta_0)$（逐元素积）。其直觉十分简洁：当外循环期望的更新方向 $v$ 与 TTO 实际移动方向 $(\theta^* - \theta_0)$ 一致时，梯度为负，减小 $\Lambda$ 以释放更多调整自由度；反之则增大 $\Lambda$ 以加强约束。这种自适应正则化无需显式的不确定性标注，仅通过新视角渲染监督即可端到端学习。

\textbf{实现与实验。}
Diff3R 的设计是架构无关的，可适配任意前馈 3DGS 骨干。实验中基于 DepthSplat（pose-given 设定）和 Depth Anything v3（pose-free 设定）构建，在 DPT 头的中间特征上附加小型 MLP 预测正则化权重 $\Lambda$。训练采用两阶段策略：第一阶段纯前馈训练（等效于 $\theta^* = \theta_0$），获得合理初始化后再激活内循环 TTO 微调高斯 DPT 头和不确定性模块。为实现跨属性（位置、尺度、球谐函数等）的均衡正则化，Diff3R 采用 LM 启发式的参数自适应缩放：有效近端项为 $\Lambda \cdot \mathrm{diag}(M)$，其中 $M \in \mathbb{R}^{N_\theta}$ 跟踪 $J^\top J$ 对角线的移动平均，使梯度幅度大的属性（如高斯中心位置）获得更强的正则化。PCG 求解器基于 3DGS-LM 的高效 CUDA 核实现，雅可比矩阵 $J$ 按渲染残差总数归一化以保证梯度稳定。训练开销为每步约 4--5 秒（原始前馈训练为 4 秒），但推理时无需额外开销。在 RealEstate10K 上 Diff3R 较前馈基线在 TTO 后提升超过 0.25 dB PSNR，在 ScanNet++ 上提升超过 0.35 dB PSNR，消融实验证实双层优化框架和不确定性正则化均不可或缺。

\paragraph{多锚点元梯度的优化感知训练}

Diff3R 通过双层优化让前馈网络学会为测试时优化（TTO）输出更好的初始化——但这一方法需要修改训练过程本身（引入内循环模拟和 IFT 梯度计算），训练开销显著增加。ForeSplat（2026.05）~\cite{2605_22020}提出了一个更轻量的替代方案：不改变前馈网络的推断方式，而是通过一种特殊的训练规则，使前馈模型\emph{无意中}产生更适合下游优化的初始化。这一思路的核心洞察是：与其让前馈网络直接输出最优结果，不如让它输出一个"好的起点"，把剩余的精炼工作交给更强大的优化器完成。

\textbf{动机与核心洞察。}标准前馈 3DGS 的训练目标是零步渲染误差（zero-step rendering error）——网络输出直接渲染的结果应尽可能接近真实图像。然而，实际使用中往往会对前馈输出执行若干步逐场景精炼优化（predict-then-refine 范式）。问题在于：标准训练完全不考虑"这个初始化对后续优化是否友好"——网络可能找到一个局部最优的高斯配置，渲染质量不错但优化器难以从中进一步改进（例如，高斯可能集中在错误的位置，虽然渲染结果碰巧可以，但优化器的梯度方向相互冲突）。

\textbf{MetaGrad 训练规则。}ForeSplat 的核心是 MetaGrad——一种轻量级多锚点元梯度训练规则。其工作流程为：(1) 前馈网络输出初始高斯 $\theta_0$；(2) 展开一段短的内循环精炼轨迹（若干步标准 3DGS 优化），得到精炼后参数 $\theta^*$；(3) 在精炼轨迹的多个时间点（"锚点"）采样中间状态 $\{\theta^{t_1}, \theta^{t_2}, \ldots\}$；(4) 聚合这些锚点处的一阶梯度（而非昂贵的二阶梯度），反向传播到前馈网络的预测头，作为优化感知的训练信号。直觉上，这相当于告诉网络："你不仅要画得对，还要画得容易被改进"。

\textbf{与 Diff3R 的关键区别。}ForeSplat 与 Diff3R 的目标相同（让前馈输出成为更好的优化起点），但技术路径不同：Diff3R 通过 IFT 和 PCG 计算精确定义的隐式梯度，理论上更严谨但计算代价更高；ForeSplat 用多锚点的一阶梯度近似，绕过了昂贵的二阶求导，训练更轻量。另一个重要区别是架构兼容性——ForeSplat 已在多种骨干（AnySplat、Pi3X、以及专为边缘部署设计的蒸馏变体）上验证，证明优化感知训练是一种与具体架构无关的通用改进策略。

\textbf{实验效果。}在所有测试架构上，ForeSplat 训练的初始化比标准训练的初始化收敛更快（更少的精炼步数达到相同质量）且达到更高的峰值质量（即使精炼到完全收敛，ForeSplat 的起点仍优于标准训练的起点）。这一结果表明，将部分场景建模负担从容量有限的前馈网络"卸载"给下游优化器，是一种实用且高效的策略。

\subsubsection{特征增强与跨域泛化}
\label{sec:13_1_5}

本路线介绍从前馈管线中特征表示的表达力和跨域适应性角度进行增强的方法。前馈 3DGS 的泛化能力依赖于强大的特征表示，多个工作从特征学习的角度进行改进。本节的方法与前述几何增强方法（\S\ref{sec:13_1}~节）的不同之处在于：几何增强方法聚焦于深度估计这一中间环节的精度，而本节方法关注特征表示本身的表达力和跨域适应性——从多视角对应的建模方式、多尺度约束的引入、高频信息的保持，到多视角信息的压缩与聚合。

\paragraph{显式 MVS 代价体与隐式 Transformer 的混合对应}

\S\ref{sec:13_1}~节中的 DepthSplat 和 TranSplat 分别代表了两条技术路线：显式的 MVS 代价体方法在纹理丰富区域精度高，但在遮挡、弱纹理等歧义区域容易失败；隐式的 Transformer 注意力方法对歧义区域更鲁棒，但收敛较慢。H3R（2025.08）~\cite{2508_03118}的核心洞察是：这两条路线并非互斥，而是互补——显式方法提供几何锚点，隐式方法负责全局精炼。

H3R 提出混合多视角对应框架，包含两个互补组件。

\textbf{体潜在融合}（Volumetric Latent Fusion）通过极线几何约束捕获跨视角对应：在逆深度空间均匀采样 $d$ 个深度平面，通过可微单应变换将邻近视角的潜在特征 warp 到参考视角，构建三维潜在体 $\bm{x}_{j \to i} \in \mathbb{R}^{h \times w \times d \times c}$。H3R 评估了三种代价计算策略——基于相关性的（缩放点积，匹配峰值尖锐但泛化差）、基于差值的（元素差拼接，边缘定位准确）和无代价的（直接拼接 warp 特征，保留完整信息供后续层学习匹配）。实验表明无代价策略在跨数据集泛化上最优，这表明让网络在后续 Transformer 层中自行学习匹配函数比手工设计相似度度量更有效。

\textbf{相机感知 Transformer}（Camera-aware Transformer）将 Pl\"ucker 射线坐标（六维参数化，编码每个像素的射线原点和方向）作为位置编码注入 Transformer 自注意力层，为多视角特征交互提供几何约束。

H3R 的另一个重要发现涉及视觉编码器的选择。通过系统比较多种视觉基础模型作为 3DGS 编码器的性能，作者发现空间对齐的基础模型（如 SD-VAE，以像素级重建为目标训练）在重建任务上\emph{显著}优于语义对齐的模型（如 DINOv2、CLIP，以语义分割/对比学习为目标训练）。SD-VAE 以更少的参数量实现了更优性能，而广泛使用的 DINOv2 反而表现不佳。这一结果揭示了语义表示与空间重建需求之间的根本性失配：语义级预训练优先保证类别级的不变性而牺牲空间精度，而 3DGS 重建需要精确的空间位置编码。H3R 还提供两个扩展变体：H3R-$\alpha$ 支持可变数量输入视角（2--8 张）并利用目标视角位姿生成全面覆盖场景的高斯；H3R-$\beta$ 在 $512 \times 512$ 高分辨率上重建，利用高分辨率输入捕获更多细节。在 RealEstate10K 上，H3R 相比 TranSplat 提升 0.91 dB PSNR；从 2 视角扩展到 8 视角时单调提升 +2.78 dB PSNR，而 MVSplat 反而下降 6.01 dB。

\paragraph{上下文-跨维度-跨尺度三维约束增强}

$C^3$-GS（2025.08）~\cite{2508_20754}在 MVSGaussian 的 MVS 管线基础上，从三个正交维度增强特征学习。其设计动机是：标准管线中的特征编码主要捕获视角间信息，缺乏对 3D 上下文一致性的感知，导致特征表示在空间上不连续。

三个模块分别作用于管线的不同位置。

\textbf{上下文感知注意力}（Context-aware Group Attention, CGA）应用于特征提取阶段：对 FPN 各层级的特征，分别沿高度和宽度执行全局平均池化得到 $T_h \in \mathbb{R}^{C \times H \times 1}$ 和 $T_w \in \mathbb{R}^{C \times 1 \times W}$，拼接后经 Conv1D 和 Sigmoid 生成两个正交注意力图 $A_h$ 和 $A_w$，逐元素乘以原始特征——这一设计以极低计算代价建立了跨行和跨列的长程依赖，突破了标准卷积的局部感受野限制。

\textbf{跨维度注意力}（Cross-Dimension Attention, CDA）在特征解码阶段融合二维图像信息和三维几何一致性：从 MVS 代价体中通过网格采样提取逐像素体素特征 $C_v \in \mathbb{R}^{HW \times 8}$，与聚合后的多视角图像特征 $F_u$ 拼接形成 $F_c \in \mathbb{R}^{HW \times 24}$，以 $F_c$ 为查询、源视角特征为键值执行跨视角注意力——这使解码后的特征同时包含细粒度二维外观和多视角几何一致性信息。

\textbf{跨尺度融合}（Cross-Scale Fusion, CSF）在粗到细级联中连接不同尺度的表示：将粗阶段的高斯特征上采样后与细阶段特征拼接，经轻量 MLP 预测调制权重 $w^{\ell+1}$ 以逐元素调整细阶段高斯的不透明度 $\alpha = \alpha \odot w^{\ell+1}$——仅调整不透明度（而非所有参数）在性能和计算效率间取得最优平衡。三个模块集成到统一渲染管线中，无需额外监督，在 DTU、Real Forward-facing、NeRF Synthetic 和 Tanks and Temples 四个数据集上实现了一致的渲染质量提升。

\paragraph{小波高频保持的轻量适配器}

H3R 和 $C^3$-GS 在管线的特定环节注入几何先验，但均未解决深度网络的一个根本性缺陷：\emph{低通滤波效应}。随着网络深度增加，高频边界信息被逐步平滑，导致高斯尺度预测退化为各向同性球体（$s_x \approx s_y \approx s_z$）——因为网络不确定局部边界方向时，"安全"的预测就是各向同性的。AdaptSplat（2026.05）~\cite{2605_10239}聚焦于这一问题，探索用最少的参数代价补偿高频衰减。

AdaptSplat 采用"最小化适配"设计哲学：以 DINOv3-ConvNeXt 为骨干（其分层卷积天然提供多尺度特征金字塔），搭配标准多视角 Transformer 和 DPT 解码器，不做任何架构级改动。唯一的定制化是频率保持适配器（Frequency-Preserving Adapter, FPA，仅约 1.5M 参数），利用二维离散小波变换（DWT）分解骨干浅层特征，LH/HL/HH 三个高频子带分别沿水平、垂直和对角方向捕获高频能量——这些方向感知的结构先验实质上为每个区域提供了方向结构张量，大幅缩小了高斯缩放矩阵 $\mathbf{S}$ 和旋转矩阵 $\mathbf{R}$ 的假设空间。

FPA 提取的高频先验通过两条路径注入通用管线。(1)~\emph{高频引导注意力位置编码}：将高频信号 $\mathbf{F}_{hf}$ 注入 Transformer 自注意力的 Query 和 Key 空间（$\text{Attention} = \text{Softmax}((\mathbf{Q} + \mathbf{F}_{hf})(\mathbf{K} + \mathbf{F}_{hf})^\top / \sqrt{d})\mathbf{V}$），而非直接修改 Value。这实现了结构相似性与聚合内容的显式解耦——注意力权重的分配受到高频结构引导（优先在结构相似的区域内聚合），而深层语义特征不被浅层高频信号干扰。(2)~\emph{高频自适应多尺度残差调制}：在 DPT 解码器中，高频特征经 Sigmoid 生成空间门控掩码 $M \in (0,1)^{H \times W}$，对多尺度特征执行自适应残差调制 $\mathbf{F}'_i = \mathbf{F}'_{i+1} + \mathbf{F}_i \odot (1 + \gamma \cdot M_i)$（$\gamma$ 为可学习缩放因子），抵消双线性上采样引入的二次高频衰减。训练中还引入焦点频率损失（Focal Frequency Loss）和透明度正则化，损失权重 $\lambda_{rec} = 1.0$，$\lambda_{ffl} = 0.1$，$\lambda_{reg} = 0.01$。在 DL3DV 和 RealEstate10K 上，AdaptSplat 仅以 1.5M 额外参数即超越复杂管线工程方法；零样本迁移至 Tanks and Temples 和 Mip-NeRF 360 同样表现优异。

\paragraph{多尺度空间适配器与存储库无关聚合}

AdaptSplat 通过 FPA 注入高频先验，但其适配策略是单向的——仅从浅层向深层传递信号。MuSASplat（2025.12）~\cite{2512_07165}从参数高效微调的角度出发，设计多尺度适配器 MuSA，在 ViT 的每个 block 中恢复被 patch 化破坏的二维空间结构。

MuSA 的核心思路是将 ViT 的 token 序列重新解释为二维特征图，并通过多核深度卷积捕获局部和全局空间关系。具体而言，每个 ViT block 的输出 $\mathbf{X}^v \in \mathbb{R}^{N \times C}$ 先经 $1 \times 1$ 线性投影降至 $C' = C/4$ 维，重塑为 $C' \times h \times w$ 的空间特征图，然后分别通过核大小为 3、5、7 的三个深度卷积并取平均，捕获从细到粗的多尺度空间依赖。结果经点卷积、GELU 激活和上采样投影回原始维度，以残差方式加回冻结的 token 流——所有参数零初始化以确保微调从预训练行为平稳起步。MuSA 还尝试了 mini-grid 分支（将每个 token 映射回 $4 \times 4$ 的 patch 内部结构），但在实验中该分支权重始终接近零，没有产生可测量的改进——这一负面结果揭示了一个有价值的洞察：适配器的主要收益来自跨 patch 的空间关系恢复，而非 patch 内部细节。

MuSASplat 的另一项创新是特征融合聚合器（Feature Fusion Aggregator, FFA），替代 Spann3R 等方法中的外部存储库。传统存储库需要 $O(V)$ 次读写（$V$ 为视角数），而 FFA 在一次前向传播中同时处理所有视角：轻量 MLP（质量估计器）为每个 token 预测置信度分数，另一个 MLP（边界检测器）从全局池化特征中识别边界视角指示器。置信度低于阈值 $\tau$ 的 token 被掩码，边界视角的权重通过可调系数 $\lambda > 1$ 增强。最终通过一次 softmax 注意力聚合所有视角信息。FFA 将峰值 GPU 内存降低至 0.3 倍，训练速度提升 4.2 倍。

\paragraph{信息瓶颈驱动的多视角特征压缩}

前述方法通常假设少量输入视角（2--8 张），当视角数量增加时，像素对齐高斯的数量随视角数线性增长，导致计算和内存瓶颈。ZPressor（2025.05）~\cite{2505_23734}从信息论的基本原理出发解决这一问题：多视角特征中包含大量冗余信息（联合熵 $H(\mathbf{F}_1, \ldots, \mathbf{F}_K)$ 远小于各视角熵之和 $\sum_i H(\mathbf{F}_i)$），去除冗余可在保留预测能力的同时压缩表示。

ZPressor 基于信息瓶颈（Information Bottleneck, IB）原则，通过三步流程实现压缩。

\textbf{锚点选择}：从 $K$ 个视角的相机位置中，以最远点采样迭代选择 $N$ 个锚点视角——第一个锚点随机选取，后续每个锚点选择与当前锚点集合距离最大的视角——确保锚点在空间上分散、有代表性。

\textbf{分配}：每个支撑视角按欧氏距离分配到最近的锚点。

\textbf{融合}：在每个锚点-支撑簇内，以锚点视角特征为查询、支撑视角特征为键值执行交叉注意力，将支撑视角的信息压缩到锚点视角中。压缩后的紧凑表示 $Z$ 仅包含 $N$ 组增强的锚点特征，直接送入后续高斯预测网络。训练目标为渲染损失（MSE + LPIPS）加 KL 散度正则化项（约束压缩后表示的复杂性），权重 $\beta$ 控制压缩与预测的平衡。ZPressor 是架构无关的，可即插即用地集成到 DepthSplat、MVSplat 等现有模型中，在 80GB GPU 上实现 480P 分辨率下超过 100 张输入视角的处理。在 DL3DV-10K 和 RealEstate10K 上，集成 ZPressor 的模型在适度视角数下一致提升性能，在稠密视角设置下显著增强鲁棒性。

\subsubsection{生成先验与增强策略}
\label{sec:13_1_6}

本路线介绍利用生成模型先验弥补前馈方法高频信息缺口的方法。前馈方法的渲染质量受限于网络容量和训练数据的覆盖范围——高频纹理、遮挡区域和外推视角的信息在稀疏输入中天然缺失。生成模型（尤其是扩散模型，Diffusion Model）在大规模数据上学到了丰富的场景先验，为弥补这一信息缺口提供了新途径。本节的方法与上节（\S\ref{sec:13_1_5}）的特征增强方法形成互补：特征增强方法在原始管线的特征提取和融合阶段改进表示质量，而本节方法在渲染输出之上叠加第二个增强阶段，利用外部生成先验补充高频细节或修复缺陷。

\paragraph{频域-空间双域细节感知与扩散增强}

前馈方法普遍受限于 ViT 骨干的内存开销，输入图像被限制在 $256 \times 256$ 等低分辨率——高分辨率图像中丰富的纹理和结构信息在降采样过程中被直接丢弃。One-Shot Refiner（2026.01）~\cite{2601_14161}的核心问题是：如何在不增加~ViT（Vision Transformer）~骨干计算负担的前提下，利用高分辨率输入中的细节信息？

其解决方案包含两个互相衔接的组件。

\textbf{双域细节感知模块}（Dual-Domain Detail Perception Module, DD-DPM）在频域和空间域联合处理高分辨率输入：频域分支通过二维傅里叶变换获取输入图像的频谱表示，经 MLP 预测重要性分数后 top-k 选择关键频率分量，以可学习复数权重调制后逆变换回空间域——频率域编码全局结构而不依赖严格的空间对应关系，因此能灵活地建模图像细节；空间域分支用轻量 CNN 捕获细粒度纹理。两个分支的输出融合后注入高斯参数预测头。

\textbf{特征引导扩散网络}在渲染阶段进一步增强：每个高斯基元被赋予额外特征向量（而非仅 RGB 和球谐系数），渲染时不仅输出 RGB 图像，还输出对应的高维特征图。这些渲染特征随后送入一步 Stable Diffusion 架构，以高斯特征为辅助输入、参考图像注入 UNet 自注意力模块实现图像间信息交换，在单步去噪中补充高频细节。One-Shot Refiner 还设计了统一的训练策略，使 ViT 几何骨干和扩散增强模块联合优化。在 DL3DV 和 RealEstate10K 上的 512 分辨率评估中，One-Shot Refiner 在 PSNR 和 FID 上均大幅领先基线方法。

\paragraph{极线距离加权注意力的几何约束扩散增强}

One-Shot Refiner 的扩散增强在图像空间操作，但未显式利用多视角几何约束来保证跨视角一致性。在宽基线（wide-baseline）稀疏视角场景中，这一问题尤为严重：视角间重叠少、纹理缺失多，扩散模型产生的增强内容容易在跨视角间不一致。ProSplat（2025.06）~\cite{2506_07670}专门针对宽基线场景，在扩散增强过程中引入显式几何约束。

ProSplat 采用两阶段框架：第一阶段以 DepthSplat 为骨干生成初始三维高斯基元并渲染新视角；第二阶段通过一步扩散模型（基于 img2img SD-Turbo）增强渲染视角。增强阶段的两个关键设计确保跨视角几何一致性。

\textbf{最大重叠参考视角注入}（Maximum Overlap Reference view Injection, MORI）从输入视角中选择与目标视角重叠度最高的视角作为参考——重叠度同时考虑平移距离（欧氏距离的倒数）和视角方向一致性（旋转矩阵第三列的余弦相似度），两者加权求和为综合评分 $\text{Score}^{i,j} = 1/\text{Dist}^{i,j}_{tgt} + (1/2)(\text{Angle}^{i,j} + 1)$。参考视角通过 VAE 编码后注入 UNet 扩展路径，为增强过程提供纹理和颜色补充。

\textbf{距离加权极线注意力}（Distance-Weighted Epipolar Attention, DWEA）在 UNet 瓶颈层对目标-参考特征执行注意力操作时，用极线几何调制注意力权重：通过基础矩阵计算每个目标像素到其在参考视角对应极线的距离 $d$，生成距离图，将全局点积注意力权重按 $\text{attn}_{comb} = \text{attn}_g \cdot \text{Norm}(\exp(-d))$ 进行调制——距极线越近的参考像素获得越大的注意力权重，远距离的无关匹配被抑制。此外，ProSplat 设计了分治训练策略（divide-and-conquer），通过数据集筛选和联合优化使两阶段数据分布对齐。在 RealEstate10K 和 DL3DV-10K 的宽基线设置下，ProSplat 平均提升 1 dB PSNR。

\paragraph{视频扩散模型驱动的全景视角合成}

ProSplat 利用一步扩散模型增强单个渲染视角，但无法保证多个增强视角之间的时序一致性——这对于 360\textdegree{} 全景合成尤为关键。MVSplat360（2024.11）~\cite{2411_04924}选择预训练的 Stable Video Diffusion（SVD）模型作为增强骨干——SVD 在视频数据上训练，天然具备时序一致性约束。

MVSplat360 的核心设计是将前馈 3DGS 模型的输出\emph{重构}到 SVD 的潜在空间。具体而言，前馈模型的渲染不再输出 RGB 图像，而是直接输出 SVD 编码器兼容的潜在特征——这些渲染特征编码了场景的三维几何和外观信息，作为位姿和视觉线索引导 SVD 的去噪过程。SVD 的去噪条件包括：渲染潜在特征（提供三维几何约束）、目标视角的位姿信息（提供相机运动约束），两者共同确保生成的多视角序列在几何上连贯、在视觉上一致。整个模型端到端可训练，仅需要 5 张稀疏输入即可实现 360\textdegree{} 高质量新视角合成。为评估这一挑战性任务，作者在 DL3DV-10K 上构建了新的 360\textdegree{} NVS 基准，MVSplat360 在大范围扫掠和完整 360\textdegree{} 合成任务上均实现了显著优于先前方法的视觉质量。

\paragraph{即插即用的高斯 Splat 几何特征适配器}

前述方法将扩散增强与特定前馈模型紧密耦合——MVSplat360 重构了渲染输出格式，ProSplat 依赖 DepthSplat 的骨干。GeoNVS（2026.03）~\cite{2603_14965}追求的是一种更灵活的设计：能否以即插即用方式为任意前馈几何模型注入几何约束？

GeoNVS 提出高斯 splat 特征适配器（GS-Adapter），其工作原理为：(1) 将输入视角的扩散特征通过 GS-Adapter 提升到 3D 高斯表示中；(2) 利用三维高斯的 splatting 操作渲染几何约束的新视角特征——高斯的位置和协方差天然编码了三维几何信息，splatting 到新视角确保渲染特征遵循正确的透视几何；(3) 将几何约束的新视角特征与扩散模型的原始特征自适应融合，校正几何不一致的表示。与在输入层级注入几何的方法不同，GS-Adapter 在特征空间操作，避免了视角相关颜色噪声对结构一致性的干扰。其实验覆盖 9 个场景和 18 种设置，相对 SEVA 和 CameraCtrl 分别实现 11.3\% 和 14.9\% 的性能提升，平移误差减少最多 2 倍，Chamfer Distance 减少最多 7 倍。

\paragraph{生成增强回灌三维重建的闭环}

前述方法将生成增强作为渲染后的后处理步骤，增强结果不再反馈到三维表示中。Leveling3D（2026.03）~\cite{2603_16211}的关键创新在于构建了一个闭环：增强后的新视角图像被重新送回前馈 3DGS 模型作为额外输入，从而提升三维重建本身的质量——"水平提升"（leveling up）3D 重建。

Leveling3D 提出几何感知水平适配器（geometry-aware leveling adapter），将扩散模型内部的生成知识与前馈模型的几何先验对齐——这一轻量适配器使扩散模型能够在由前馈表示欠约束区域引起的渲染伪影区域上执行几何一致的生成，而非无约束的幻觉填充。训练中引入调色板过滤策略（palette filtering）促进更多样化的分布学习，测试时使用掩码精炼（mask refinement）防止修复区域的边界混乱。更重要的是，增强后的外推新视角可回灌到前馈 3DGS 中作为输入，形成"生成$\to$重建$\to$再生成"的正反馈循环。在新视角合成和深度估计两个任务上，Leveling3D 均达到最佳性能。

\paragraph{置信度感知的扩散修复与三维表示保护}

ConFixGS（2026.05）~\cite{2605_09688}面临的应用场景与前述方法不同——轨迹式驾驶场景中，输入视角沿车辆行驶路径稀疏分布，视角间重叠有限且外推需求大。直接用扩散模型修复前馈输出的缺陷存在风险：扩散模型可能产生多视角不一致的幻觉内容，直接用其监督高斯优化会腐蚀三维表示。

ConFixGS 的核心思想是：将扩散输出视为\emph{候选证据}而非真值。管线分为三个阶段。(1)~初始前馈重建：预训练前馈模型从稀疏支撑视角生成初始高斯 $\mathcal{G}_0$；在每个新视角的邻近轨迹位置构建局部伪视角片段，用 SD-Turbo 增强局部渲染，生成细节丰富的伪目标 $I_n^*$。(2)~置信度估计：对每个伪目标像素，利用 $\mathcal{G}_0$ 的渲染深度作为几何脚手架，将其反投影到三维空间再投影到支撑视角中，与真值 RGB 进行光度一致性比较——重投影误差 $e_n(p)$ 越小（经高斯核平滑后），该像素的置信度 $w_n(p) \in [0,1]$ 越高。无法通过任何支撑视角验证的像素赋予保守基线值 $\gamma$，无深度覆盖的区域直接置零。(3)~置信度引导全局修复：以置信度加权优化高斯参数——低置信度像素的光度梯度和致密化信号均被抑制，防止幻觉内容污染三维表示。在 Waymo、nuScenes 和 KITTI 上，ConFixGS 实现 PSNR 最高提升 3.68 dB、FID 降低近半的效果，展示了"置信度感知融合生成先验与支撑视角一致性"作为鲁棒重建关键原则的有效性。

\paragraph{点 Transformer 直接精炼高斯集合}

前述方法均在二维图像空间或潜在空间操作扩散模型，SplatFormer（2024.11）~\cite{2411_06390}选择了一条根本不同的路线：直接在三维高斯基元集合上操作。这是首个为高斯 splat 设计的点 Transformer 模型。

SplatFormer 接收有限训练视角优化得到的初始 3DGS 集合作为输入（而非二维图像），通过单次前向传播直接精炼高斯参数——每个高斯被视为一个"点"，其属性（位置、尺度、旋转、颜色、不透明度）作为点特征。Transformer 的自注意力在三维高斯之间建立全局交互：精炼过程中，一个高斯的属性调整可以参考空间中其他所有高斯的状态，实现全局协调的修复——例如，消除在分布外视角下暴露的"针刺"伪影需要同时调整多个空间上相邻的高斯的尺度。SplatFormer 成功解决了先前多场景训练方法仅能处理有限数量输入视角的局限，在极端新视角下的渲染质量达到最佳性能，超越了各种 3DGS 正则化技术、稀疏视角多场景模型和扩散增强框架。

\subsubsection{特殊成像条件与鲁棒性}
\label{sec:13_1_7}

本路线介绍在分辨率不匹配、噪声、低光照等非理想条件下维持前馈重建鲁棒性的方法。实际应用中，输入图像可能处于多种恶劣条件——分辨率不匹配、视角数量变化、噪声和低光照等。前馈方法在这些条件下也需要维持重建质量，而非仅在干净、高分辨率的实验室条件下有效。本节与前两节形成递进关系：\S\ref{sec:13_1_5} 改进特征表示的表达力，\S\ref{sec:13_1_6} 用生成先验弥补信息缺口，本节则处理输入端的非理想条件——从抗混叠到退化鲁棒性再到超分辨率重建。

\paragraph{不透明度平衡的带限抗混叠渲染}

前馈 3DGS 方法普遍使用标准的屏幕空间膨胀滤波器，这一设计在训练分辨率下表现正常，但当渲染采样率偏离训练分布时（放大、缩小、改变焦距），会产生严重的渲染伪影——放大时出现网格状侵蚀和虚假透明（退化高斯小于单个像素），缩小时出现锯齿。AA-Splat（2026.03）~\cite{2603_29394}首次系统分析了这一问题的根源：前馈模型生成的大量高斯是退化高斯（degenerate Gaussians），其尺度远小于 Nyquist 采样率所需的最小值。

AA-Splat 提出基于不透明度平衡的带限（Opacity-Balanced Band-Limiting, OBBL）设计，包含两个互补组件。

\textbf{三维带限后滤波器}（3D Band-Limiting Post-Filter, 3D-BLPF）为每个高斯计算其在所有输入视角中的最大采样率 $\hat{\nu}_j = \max(\{\mathds{1}^{(i)}(\mu_j) \cdot f^{(i)} / d_j^{(i)}\}_{i=1}^N)$，据此构造球形各向同性三维高斯 $\mathcal{G}_j^{\text{low}}$ 作为最小尺度约束，通过三维卷积 $\mathcal{G}_j^{\text{reg}} = \mathcal{G}_j \otimes \mathcal{G}_j^{\text{low}}$ 将原始高斯的尺度扩展为 $s_{j,x}^{\text{reg}} = \sqrt{s_{j,x}^2 + (s_j')^2}$（$s_j' = \sigma_s / \hat{\nu}_j$，$\sigma_s = 0.2$ 为滤波大小超参数），确保没有任何高斯退化到 Nyquist 频率之下。

\textbf{不透明度平衡}（Opacity Balancing, OB）解决带限扩展的副作用：扩展后的像素对齐高斯之间重叠增加，后排高斯的可见性下降，在放大渲染时产生不自然的网格伪影。OB 通过将所有高斯的不透明度裁剪至上限 $\tau_{\text{opa}} < 1$，防止任何高斯变得完全不透明，使所有高斯无论深度排序如何都能充分贡献渲染；同时渲染累积 alpha 值用于归一化颜色，补偿裁剪引入的透明度损失。AA-Splat 以 DepthSplat 为骨干实现，在所有测试分辨率（$4\times$ 到 $1/4\times$）上相对基线实现 5.4--7.5 dB PSNR 提升。

\paragraph{Alpha 归一化与三维采样正则化}

像素对齐前馈方法的一个被忽视的缺陷是：当推理时的输入视角数量与训练时不一致，渲染图像会出现系统性过亮。RoSplat（2026.05）~\cite{2605_13093}分析了这一现象的根源：像素对齐高斯按深度排序执行 $\alpha$-混合时，重叠高斯的数量随视角数线性增长，每个额外视角的高斯贡献额外的颜色累积，导致亮度偏移。

RoSplat 的解决方案简洁而有效。

\textbf{Alpha 归一化}：在标准 $\alpha$-混合渲染后，同时渲染累积 alpha 值 $\alpha(\mathbf{x}) = \sum_i \alpha_i^{\text{clip}} G_i^{2D}(\mathbf{x}) \prod_{j<i}(1 - \alpha_j^{\text{clip}} G_j^{2D}(\mathbf{x}))$，以最终颜色除以 alpha 值 $C_{\text{norm}}(\mathbf{x}) = C(\mathbf{x}) / \alpha(\mathbf{x})$——直觉上，视角越多 alpha 累积越高，归一化后颜色恢复到物理正确的亮度水平。这一策略不改变训练过程，仅在推理时应用。

\textbf{三维采样正则化器}解决高分辨率渲染中的孔洞伪影：高斯尺度估计不准确时，在高分辨率下放大渲染会暴露高斯之间的间隙。RoSplat 引入三维空间中的辅助采样约束，鼓励高斯尺度预测的准确性——通过在训练中对高斯三维位置施加采样一致性约束，使高斯更紧密地覆盖表面，减少高分辨率下的间隙。在基准数据集上，RoSplat 显著改善了基线模型在可变输入视角和高分辨率渲染设置下的表现。

\paragraph{退化学习器与语义状态空间的通用鲁棒框架}

AA-Splat 和 RoSplat 处理的是分辨率和视角数量变化，但实际图像还可能受到噪声、低光照、雨雾等多种退化影响。这些退化在前馈管线中引发连锁反应：噪声干扰特征匹配导致深度估计错误，低光照破坏颜色一致性导致外观参数退化，最终重建质量严重下降。RobustGS（2025.08）~\cite{2508_03077}提出一个通用框架，以即插即用方式增强现有前馈方法在多种退化条件下的鲁棒性。

RobustGS 包含两个核心组件。

\textbf{广义退化学习器}（Generalized Degradation Learner）从多视角输入中提取多种退化类型的通用表示和分布——不同于针对单一退化类型的专用模块，退化学习器在训练时接触噪声、模糊、雨雾、低光照等多种退化，学习它们在特征空间中的共享结构，使模型具备退化感知能力。

\textbf{语义感知状态空间模型}（Semantic-Aware State Space Model）在两个层面增强鲁棒性：首先利用提取的退化表示在特征空间增强受损输入（去噪/去雾的方向性先验），然后采用语义感知策略聚合跨视角的语义相似信息——在语义上一致的区域（如同一面墙、同一片天空），跨视角信息互补增强；在语义不一致的区域（退化导致的虚假纹理），避免错误聚合。RobustGS 可以无缝集成到任何预训练前馈管线中，无需重新训练整个模型，在噪声、低光照、雨雾等多种退化条件下均实现一致的质量提升。

\paragraph{残差增强范式的低光照适配}

RobustGS 追求通用退化鲁棒性，DelowlightSplat（2026.05）~\cite{2605_26629}则专门针对低光照场景——这一条件下，图像噪声放大、颜色偏移、特征匹配可靠性下降等问题尤为严重，且这些困难相互叠加。

DelowlightSplat 引入轻量级低光照适配器（Lowlight Adapter），在残差增强范式中工作：适配器不是从零重建正常光照图像，而是预测低光照输入与正常光照之间的残差，通过加法修正恢复图像质量——这一设计避免了对已经微弱的光照信号做大幅变换，更稳定地提升可匹配性。增强后的图像送入基于代价体的多视角推断管线，直接预测干净的三维高斯——而非先做图像增强再重建的级联方案，避免了增强伪影传播到三维表示中。为支持这一方向的系统性研究，作者构建了可控的多视角低光照基准：仅退化上下文视角而保持目标视角干净，确保评估的是模型在低光照输入下的重建能力而非渲染能力。实验表明 DelowlightSplat 在低光照条件下显著优于通用前馈方法和两阶段管线方案。

\paragraph{低分辨率视角直接前馈三维超分辨率}

当输入图像为低分辨率时，传统的三维超分辨率（3DSR）方法是先用二维超分辨率模型提升输入图像分辨率，再用高分辨率图像进行逐场景 3DGS 优化——这一管线有两个根本性局限：高频先验完全继承自二维超分辨率模型，而非从三维数据中学习；且需要逐场景优化，缺乏跨场景泛化能力。

SR3R（2026.02）~\cite{2602_24020}将三维超分辨率重构为\emph{从低分辨率视角直接前馈预测高分辨率 3DGS}的映射问题——使模型在大规模多场景数据上自主学习三维特定的高频几何和外观先验。具体而言，SR3R 是即插即用的框架，可与任何前馈 3DGS 重建骨干配合：骨干提供低分辨率 3DGS 脚手架，SR3R 通过高斯偏移学习（Gaussian Offset Learning）和特征精炼将低分辨率高斯上采样为高分辨率高斯。高斯偏移学习为每个低分辨率高斯预测一组三维空间偏移量，生成多个高分辨率高斯；特征精炼模块在上采样过程中稳定重建并锐化高频细节。这一范式性改变使 SR3R 获得了强大的零样本泛化能力——在三个三维基准上，SR3R 不仅超越最佳三维超分辨率方法，在未见场景上甚至优于逐场景优化方法。

\paragraph{扩散参考引导与纹理感知密度控制}

SR3R 从三维数据中学习高频先验，而 SRSplat（2025.11）~\cite{2511_12040}则探索了一条互补的路线：利用外部高质量参考信息补偿低分辨率输入中的纹理缺失。

SRSplat 的核心洞察是：低分辨率输入的纹理信息不足可以两种方式补偿——外部的场景特定高质量参考图像，以及输入图像内部的纹理线索。为此，SRSplat 首先为每个场景构建场景特定参考画廊（reference gallery），利用多模态大语言模型（MLLM）理解场景描述，再通过扩散模型生成对应的高质量参考图像。参考引导特征增强模块（Reference-Guided Feature Enhancement, RGFE）将低分辨率输入图像的特征与其对应的参考孪生图像的特征对齐并融合——直觉上，参考图像提供了该场景应有但被低分辨率丢失的纹理细节。融合后的多视角特征经解码器预测高斯基元。此外，纹理感知密度控制（Texture-Aware Density Control, TADC）根据低分辨率输入图像的内部纹理丰富度自适应调整高斯密度——纹理丰富的区域分配更多高斯捕捉细节，平坦区域保持稀疏。SRSplat 在 RealEstate10K、ACID 和 DTU 上超越现有方法，并展现了强大的跨数据集和跨分辨率泛化能力。

\subsubsection{高分辨率与高效渲染}

本路线介绍将前馈 3DGS 扩展到高分辨率渲染并实现高效推理的方法。将前馈 3DGS 扩展到高分辨率（如 4K）渲染面临基元数量二次增长的瓶颈：像素对齐方法中高斯数量随输入分辨率线性增长，4K 图像将产生数千万乃至上亿个基元，远超实时渲染的承受能力。

\paragraph{紧凑基元加逐基元纹理贴图}

LGTM（2026.03）~\cite{2603_25745}（Less Gaussians, Texture More）的核心洞察是：高分辨率渲染所需的细节并非必须通过更多高斯基元来表达，而可以通过为每个基元附着纹理贴图来解耦几何复杂度与渲染分辨率。具体而言，LGTM 不预测像素对齐的密集高斯，而是预测一组数量紧凑的高斯基元（数量与场景复杂度相关而非与像素数相关），同时为每个基元预测一个纹理贴图。渲染时，每个高斯不再使用单一颜色或球谐系数，而是根据投影像素在纹理贴图中的采样位置获取颜色——这使得少数紧凑基元即可表达高分辨率纹理细节。这一设计将几何表达（少量基元的三维位置和形状）与外观表达（逐基元纹理贴图中的高频细节）显式解耦，首次在前馈框架下实现 4K 新视角合成，同时使用的高斯基元数量显著少于像素对齐方法。

\paragraph{像素对齐三角形前馈表面重建}

与所有使用高斯基元的前馈方法不同，FTSplat（2026.03）~\cite{2603_05932}选择三角形基元作为重建目标——直接从多视角校准图像预测连续三角形表面。其动机在于：高斯基元虽然渲染质量高，但不直接提供流形几何（manifold geometry），无法直接用于物理模拟和碰撞检测等下游任务。FTSplat 设计了像素对齐三角形生成模块，为每个像素预测三角形顶点的三维坐标，同时引入相对三维点云监督（relative 3D point cloud supervision）增强几何学习的稳定性和一致性。生成的三角形网格可立即用于标准图形引擎和机器人模拟器，无需后处理或逐场景优化，在保持高效重建的同时提供了显式的流形几何。

\paragraph{几何引导多曝光融合的 HDR 重建}

现有前馈方法假设输入图像具有一致的曝光和色调映射，忽略了高动态范围（HDR）重建的需求。InstantHDR（2026.03）~\cite{2603_11298}设计前馈网络从无标定多曝光低动态范围（LDR）图像集合在单次前向传播中重建 HDR 三维场景。技术上，InstantHDR 提出几何引导的外观建模（geometry-guided appearance modeling）实现多曝光融合——将不同曝光的 LDR 图像在统一的三维几何框架中融合为 HDR 表示，而非在二维图像空间中进行曝光合并。同时设计元网络（meta-network）实现可泛化的场景特定色调映射，适应不同场景的相机响应函数。由于缺乏 HDR 场景训练数据，作者构建了 HDR-Pretrain 数据集，包含 168 个 Blender 渲染场景，涵盖多样光照类型和多种相机响应函数。InstantHDR 在单次前向设置下实现与最先进逐场景优化 HDR 方法可比的合成质量，同时获得约 700 倍的重建加速。

\paragraph{解码器仅架构与 KV-cache 隐式表示}

大型视图合成模型（LVSMs）普遍采用编码器-解码器架构——编码器从输入视角提取场景表示，解码器从该表示渲染新视角。DVSM（2026.05）~\cite{2605_29891}重新审视了这一设计，通过对照实验证明：解码器仅架构（将场景隐式表示为 KV-cache）在更少参数量下一致优于编码器-解码器变体。其核心发现是：颜色输入重建网络与仅相机渲染网络之间共享权重，能更好地对齐同一视点的特征，促进图像合成。在此基础上，DVSM 进一步引入基础模型先验（foundation model priors）增强特征表示，配合阶段式 patch 大小调整（stage-wise patch sizing）在效率和质量间取得最优平衡。在多个基准上，DVSM 建立了新的最佳性能，在稠密输入视角下甚至超越逐场景优化的 3DGS。

\paragraph{渲染-比较模块与混合 KV 缓存压缩}

在线新视角合成要求模型从按序到达的观测中逐步构建场景表示，但自回归框架面临一个核心困境：训练时使用真值位姿确保稳定性，但推理时位姿由模型自身预测，分布不匹配导致性能退化。ReCoSplat（2026.03）~\cite{2603_09968}提出渲染-比较（Render-and-Compare, ReCo）模块解决这一问题——ReCo 将当前重建结果从预测视点渲染，与输入观测比较，为模型提供稳定的条件信号以补偿位姿误差，使得即使位姿不准确也能维持连贯的场景重建。为支持长序列在线处理，ReCoSplat 设计混合 KV 缓存压缩策略，结合浅层截断（early-layer truncation，丢弃浅层的细粒度注意力缓存）和块级选择性保留（chunk-level selective retention，仅保留关键帧的深层语义缓存），将缓存大小减少 90\% 以上，支持 100 帧以上的长序列在线重建。ReCoSplat 在有位姿/无位姿、有内参/无内参的多种输入设置下均达到最佳性能。

\subsubsection{语义感知与场景理解}

本路线介绍将语义信息整合进前馈 3DGS、以同时支持新视角合成和场景理解的方法。在三维空间中同时编码外观和语义，使模型不仅能渲染逼真图像，还能执行开放词汇分割、物体定位等场景理解任务。

\paragraph{偏移学习与点级交互的语义高斯}

GSsplat（2025.05）~\cite{2505_04659}提出可泛化语义高斯 splatting 方法，其核心设计是在多任务框架中同时预测高斯的位置、颜色属性和语义属性。架构上，GSsplat 设计混合网络（hybrid network）从输入图像提取颜色和语义两类互补信息，并预测场景自适应的高斯参数——替代传统 3DGS 中的分裂和剪枝流程。为增强高斯的空间感知能力，GSsplat 引入两个关键模块：(1) 偏移学习模块（offset learning module）通过基于分组的监督（group-based supervision）预测高斯位置偏移，使高斯能精确定位到场景结构的关键位置；(2) 点级交互模块（point-level interaction module）通过空间单元聚合（spatial unit aggregation）在三维高斯之间建立信息交换通道，使每个高斯的属性预测能参考空间邻域的上下文。在可变数量多视角输入的评估中，GSsplat 以最快的渲染速度实现了语义合成任务的最佳性能。

\paragraph{双分支解耦与损失引导视角采样的统一重建}

UniForward（2025.06）~\cite{2506_09378}提出统一三维场景和语义场重建的前馈模型，仅从无标定、无位姿的稀疏视角图像预测携带各向异性语义特征的三维高斯。为在统一表示中同时编码场景外观和语义信息，UniForward 将语义特征嵌入三维高斯基元，并通过双分支解耦解码器（dual-branch decoupled decoder）分别预测几何/外观属性和语义特征——解耦设计避免了两类特征在训练中的梯度干扰。训练中，UniForward 提出损失引导的视角采样器（loss-guided view sampler），从易到难逐步选择训练视角——初始阶段使用损失较低（重建较好）的视角确保稳定训练，随着训练推进引入损失较高（重建困难）的视角进行精炼——这一策略消除了先前方法对深度真值或语义掩码的需求。整个模型通过光度损失和从预训练二维语义模型蒸馏的语义损失端到端训练。推理时，UniForward 可实时重建三维场景及对应语义场，渲染的语义特征可在开放词汇范式下解码为稠密分割掩码。

\paragraph{文本引导的多模态语义融合}

TextSplat（2025.04）~\cite{2504_09588}是首个文本驱动的可泛化高斯 splatting 框架，其核心思想是利用文本描述作为语义桥梁，引导多模态信息的融合以增强三维重建质量。框架包含三个并行模块获取互补表示：扩散先验深度估计器（利用预训练扩散模型的几何先验提供准确的深度信息）、语义感知分割网络（从输入图像提取细粒度语义信息）、多视角交互网络（通过跨视角特征交互精炼多视角表示）。三种表示在文本引导语义融合模块（Text-Guided Semantic Fusion Module）中通过注意力机制集成——文本嵌入引导网络关注与语义描述相关的图像区域，使融合后的三维高斯参数同时编码丰富的几何、外观和语义线索。实验表明，文本引导的语义融合显著提升了复杂场景中细粒度细节的重建质量。

\paragraph{多教师蒸馏的通用高斯特征编码}

Chorus（2025.12）~\cite{2512_17817}关注一个被忽视的问题：如何为三维高斯基元学习通用、丰富的特征表示。现有前馈方法通常仅预测高斯的几何和外观参数，而未将通用的视觉理解能力编码到高斯基元中。Chorus 提出多教师预训练框架，从三类互补的二维基础模型蒸馏信号：语言对齐教师（language-aligned teacher，如 CLIP，提供高级语义理解）、通用教师（generalist teacher，提供空间和结构先验）、物体感知教师（object-aware teacher，提供实例级判别能力）。共享的三维编码器通过教师特定的投影头（projector）学习同时捕获从高级语义到细粒度结构的信号。Chorus 在多种任务上评估：开放词汇语义/实例分割、线性/解码器探测、数据高效监督以及基于 LLM 的问答。一个引人注目的发现是：仅使用高斯中心、颜色和估计法线作为输入的编码器变体，在点云基准上甚至超越了使用 39.9 倍更多训练场景的点云基线方法，表明三维高斯基元的空间位置和颜色本身已包含丰富的几何和语义线索。此外，Chorus 提出渲染-蒸馏适配策略（render-and-distill adaptation）支持域外微调。

\paragraph{2DGS 桥接渲染与几何的表面重建}

MeshSplat（2025.08）~\cite{2508_17811}利用二维高斯 surfel（2DGS）作为桥梁，连接新视角合成与几何先验学习，从稀疏视角输入实现可泛化的表面重建。其关键设计思路是：先用前馈网络预测逐视角像素对齐的 2DGS，通过可微分渲染合成新视角图像（提供渲染监督而无需三维真值），再将 2DGS 的位置和朝向约束与几何先验对齐。具体而言，MeshSplat 提出加权 Chamfer 距离损失正则化深度图——在输入视角的重叠区域施加更强的深度一致性约束，确保跨视角几何一致；同时训练法线预测网络，利用单目法线估计器的先验对齐 2DGS 的方向（朝向），使预测的 surfel 朝向与场景表面的法向量一致。这一将渲染质量与几何先验联合优化的设计，在可泛化稀疏视角网格重建任务上达到最佳性能。

\subsubsection{其他架构创新}

本路线介绍从前馈管线其他架构维度进行创新探索的方法。除前述的密度分配、表示设计和几何增强方向外，多篇工作从其他架构维度进行了创新探索。本节介绍 11 条技术路线：图结构建模高斯关系（Gaussian Graph Network）、亚像素级基元检测（Off The Grid）、生成式高斯增密（Generative Densification）、混合 Splat-Voxel 表示（Fuse-and-Refine）、统一体量的城市场景（EVolSplat）、卫星-地面跨视角融合（Cross-View Splatter）、航空稀疏视角重建（AnyCity）、解耦几何与外观的立体视觉（Stereo-GS）、外推视角合成（EVPGS）、空间连贯的高斯序列（Z-Order Transformer）和分摊式自监督 2DGS（AIR）。

\paragraph{图结构建模跨视角高斯关系}

前馈方法简单拼接不同视角的像素对齐高斯作为场景表示，忽略了跨视角高斯之间的空间关系。Gaussian Graph Network（GGN）（2025.03）~\cite{2503_16338}构建高斯图显式建模不同视角高斯组之间的关系——图中的节点为来自各视角的高斯组（按空间邻近性聚类），边连接跨视角的对应高斯组。GGN 将基本图操作（消息传递、特征聚合）重新设计为适用于高斯表示的图运算，使每个高斯能从其相连的高斯组中获取跨视角信息（Gaussian feature fusion），增强多视角一致性。此外，GGN 设计高斯池化层（Gaussian pooling layer）将多个高斯组聚合为更紧凑的表示，在减少高斯数量的同时保留关键信息。在 RealEstate10K 和 ACID 数据集上，GGN 以更少的高斯基元实现了更好的图像质量和更高的渲染速度。

\paragraph{熵驱动的亚像素级基元检测}

像素对齐方法将高斯基元放置在图像像素网格上，这一刚性约束限制了基元的空间分布灵活性——平坦区域的像素不需要高斯，结构边缘处的像素可能需要多个高斯。Off The Grid（2025.12）~\cite{2512_15508}受关键点检测启发，将前馈 3DGS 的基元预测重新设计为亚像素级别的检测问题：解码器学习在每个图像图块（patch）内局部分布基元，取代沿像素网格的刚性排列。同时引入自适应密度机制，根据每个图块的 Shannon 熵分配不同数量的基元——高熵图块（纹理复杂）获得更多基元，低熵图块（平坦区域）获得更少甚至零个基元。整个解码器与预训练三维重建骨干端到端训练，仅使用光度监督而无需任何三维标注。生成的无位姿模型在远少于竞争方法的高斯基元下实现了更高质量的重建，展现了更准确、更高效的基元分配策略。

\paragraph{特征上采样驱动的生成式高斯增密}

前馈模型预测的高斯数量通常不足以表达高频细节（如精细纹理和薄结构）。原始 3DGS 的分裂-克隆增密策略是为逐场景优化设计的，直接迁移到前馈场景中泛化性不佳。Generative Densification（2024.12）~\cite{2412_06234}提出在前馈管线内部通过上采样特征表示并在单次前向传播中生成精细高斯——具体而言，将前馈模型的中间特征表示上采样至更高分辨率，利用网络在大规模训练中学到的先验知识为上采样后的每个位置预测精细高斯属性。这一"生成式"增密策略与原始 3DGS 的迭代式分裂-克隆根本不同：后者在参数空间操作原始高斯属性（缺乏泛化性），前者在特征空间操作上采样后的表示（利用学到的先验增强泛化）。在物体级和场景级重建任务上，Generative Densification 以可比或更小的模型尺寸超越先前最佳方法，在高频细节表达上取得显著改进。

\paragraph{Splat-Voxel-Splat 融合精炼范式}

像素对齐高斯在输入视角重叠时产生大量冗余，且基元位置被约束在相机射线上缺乏三维灵活性，难以自然扩展到动态场景。Fuse-and-Refine（2025.03）~\cite{2503_14698}引入混合 Splat-Voxel 表示解决这些问题。其管线为：首先从初始像素对齐高斯基元中，通过可微分聚合将局部特征提取到粗到细的体素层次结构中——低分辨率体素捕获全局结构，高分辨率体素保留局部细节；然后用稀疏体素 Transformer 在体素空间中处理这些特征，建立跨视角的三维一致性；最后从精炼后的体素特征生成新的高斯基元，替代原始的像素对齐基元。这一"Splat$\to$Voxel$\to$Splat"的融合-精炼范式有效减少了冗余，自然支持任意数量输入（包括时序帧）的融合，在静态和流式动态场景重建中均达到最佳性能，以 15 FPS 和 350ms 延迟在单张 H100 GPU 上运行。

\paragraph{统一三维体量跨帧联合预测}

现有前馈像素对齐 3DGS 方法在城市场景中常出现多视角不一致和内容重复——因为不同视角独立预测的高斯在重叠区域难以对齐。EVolSplat（2025.03）~\cite{2503_20168}采用根本不同的策略：在统一三维体量（unified volume）中通过三维卷积网络跨帧联合预测高斯基元，而非逐视角独立预测。具体流程为：首先用噪声深度预测初始化三维高斯位置，然后在三维空间中通过三维卷积精炼几何属性（位置、尺度、旋转），同时基于二维纹理预测颜色——三维卷积天然建立跨帧的空间一致性。此外，EVolSplat 设计灵活的半球背景模型（hemisphere background model）处理远景和天空——远距离内容被参数化为以相机为中心的半球面上的高斯，避免与近景基元竞争。在 KITTI-360 和 Waymo 数据集上，EVolSplat 实现了前馈 3DGS 和 NeRF 方法中的最佳质量。

\paragraph{卫星-地面正射校正特征对齐融合}

大规模户外场景的完整重建需要良好的相机覆盖，但地面采集耗时且难以扩展。Cross-View Splatter（2026.05）~\cite{2605_19656}利用卫星影像提供全局几何先验——通过公开 API 轻松获取的卫星图像与 GPS 标记的地面照片在统一三维坐标系中融合，预测像素对齐的高斯 splat。关键技术挑战在于卫星视角（俯视正射投影）与地面照片（透视投影）之间存在巨大的视角和分辨率差异。Cross-View Splatter 先对卫星视角进行正射校正（orthorectification），然后通过对齐地面和鸟瞰特征表示，使模型能同时利用两种视角的互补信息——卫星提供场景的全局布局和屋顶结构，地面照片提供立面细节和近距离纹理。在专门构建的地理参照基准上，Cross-View Splatter 相比仅使用地面图像的方法显著改善了场景覆盖和新视角合成质量。

\paragraph{观测锚定与生成补全的分层重建}

从稀疏航空视角重建大规模城市场景面临强烈的证据不平衡（evidence imbalance）：屋顶和开放区域被重复观测，而立面、远处建筑和遮挡结构获得很少的多视角支撑。直接回归确定性表示常导致鬼影、熔化立面和拉伸纹理。AnyCity（2026.05）~\cite{2605_19949}提出观测锚定的生成重建框架，将场景内容分为两层处理：(1) 观测支撑层——先预测观测支撑的几何潜在表示（observation-supported geometry latent），锚定有可靠多视角证据的结构；(2) 生成补全层——用脚手架条件化的航空补全 token（scaffold-conditioned aerial completion tokens）预测弱约束内容的门控残差更新，仅在观测不足的区域激活生成先验。训练中采用稠密到稀疏蒸馏（dense-to-sparse distillation）从稠密视角重建迁移结构线索，同时航空适配的视频扩散先验通过门控 token 条件化提供细粒度的城市外观线索。观测保持目标（observation-preserving objectives）确保精炼后的表示与输入支撑的几何一致。在合成、航空域、UAV 纹理和真实世界场景上，AnyCity 一致优于前馈基线。

\paragraph{解耦点图与外观的立体视觉重建}

现有前馈方法通常纠缠高斯几何和外观的预测——两者共享骨干特征，高度依赖数据驱动的先验，导致回归速度慢且资源需求大。Stereo-GS（2025.07）~\cite{2507_14921}提出解耦框架：立体视觉骨干从局部图像对提取特征后，通过全局注意力块融合跨视角信息，然后由专用的点预测头和独立的高斯预测头分别生成几何和外观。点预测头生成多视角点图（point-map）表示几何结构，高斯预测头输出外观特征——两者组合为 GS-map 表示三维场景。精炼网络进一步增强 GS-map 的质量。这一解耦设计不依赖相机参数，实现无位姿的端到端重建，在降低资源需求的同时保持高质量输出。

\paragraph{遮挡感知的外推视角合成}

标准高斯 splatting 方法假设训练视角充分覆盖目标视角，仅在内插视角上表现良好。EVPGS（2025.03）~\cite{2503_21816}解决更具挑战性的外推视角合成（Extrapolated View Synthesis, EVS）任务——从有限覆盖的训练视角泛化到训练覆盖之外的新视角。其粗到细视角增强框架分为两阶段：粗阶段在外观和几何两个层面引入正则化，减少视角覆盖不足导致的渲染伪影；细阶段生成可靠的视角先验（view prior）提供进一步训练引导，并引入遮挡感知（occlusion awareness）确保先验质量。为全面评估 EVS 任务，作者收集了专门的外推场景数据集 Merchandise3D。在三个数据集上，EVPGS 在外推视角的合成质量上显著优于先前方法。

\paragraph{Z-order 曲线组织的空间连贯高斯}

前馈方法预测的非结构化高斯集合缺乏空间组织，Transformer 在处理这种无序序列时效率低下。Z-Order Transformer（2026.05）~\cite{2605_13465}利用 Z-order 曲线（一种空间填充曲线）将非结构化高斯集合组织为空间连贯的有序序列——空间上相邻的高斯在序列中也相邻，使 Transformer 的局部注意力能高效捕获空间关系。在此基础上，通过稀疏注意力机制自适应抑制冗余高斯：空间上高度重叠的高斯被合并或抑制，而位于结构关键位置（如边缘、角点）的高斯被保留。整个模型在单次前向传播中完成高斯压缩和属性预测，以更少的高斯基元实现快速且高质量的视角合成。

\paragraph{预测-优化-蒸馏的分摊式 2DGS}

2D 高斯 splatting 的迭代优化虽能产生高质量重建，但每张图像需要耗时的逐图像优化。AIR（2026.05）~\cite{2605_20820}提出自监督前馈框架，将迭代高斯拟合分摊到单次网络前向传播中——训练时学习迭代优化的行为，推理时网络直接预测优化后的结果。AIR 采用阶段式残差架构（stage-wise residual architecture）逐步预测额外的高斯基元：每个阶段从上一阶段的重建残差中预测新基元，仅补充未充分重建的区域。显式的阶段控制机制（Stage Control）确保新基元仅在需要的位置被激活。预测-优化-蒸馏（Predict-Optimize-Distill）训练策略稳定多阶段预测：先用短视野优化生成高斯增量作为教师信号，再将优化结果蒸馏回前馈预测器，逐步使预测器逼近优化器的行为。最终模型配备图像自适应量化器（image-adaptive quantizer）实现紧凑的高斯存储，在 Kodak 和 DIV2K 上超越代表性高斯基线方法，编码时间仅 160--300ms。

张 et al.（2025.07）~\cite{2507_14501}的综述论文系统回顾了前馈三维重建和视角合成的最新进展，涵盖点云、3DGS、NeRF 等多种表示架构。

\begin{takeaway}
多视角前馈重建通过跨视角注意力机制和像素对齐高斯预测实现单次前向传播的 3DGS 生成，核心权衡在于 Transformer 效率与重建质量的平衡，Z-order 空间组织和预测-优化-蒸馏策略是提升效率的关键创新。
\end{takeaway}

%=====================================================================

\subsection{近期方法脉络：稳定前馈表示、紧凑前馈高斯预测、前馈高斯预测等}

这一组近期工作延续 feed-forward Gaussian splatting、feed-forward 3DGS、3DGS，主要把改动集中在增量改进、扩展、适配等环节。下面只保留相对既有路线的变化，便于和前文方法直接对照。

\paragraph{Learning Stable Canonical Worlds for Novel View Synthesis and Beyond}
Learning Stable Canonical Worlds for Novel View Synthesis and Beyond（2026.06）~\cite{2606_23027} 可作为“稳定前馈表示”中的新锚点候选。相较于 feed-forward Gaussian splatting，主要改变是：CanonicalGS：用深度、语义和不确定性估计抽取view-centric evidence，再在canonical latent world中不确定性感知融合。3DGS 的显式高斯表示和可微渲染能力保持为基础。
其设计动机是：从view-dependent prediction转向scene-centric representation，减少多输入视图下的冗余和噪声积累。
具体而言，方法先从各视角提取深度、语义特征和不确定性，再在canonical latent world中融合，并解码为稳定高斯场。
这一锚点判断来自当前分类中的机制关系：候选锚点：路线变化，且相关性 0.95；后续方法可围绕 用canonical latent world融合多视图证据 继续比较和扩展。
适用边界是：canonical融合依赖深度/语义/不确定性估计质量；复杂动态场景仍难。

\paragraph{ZipSplat}
ZipSplat（2026.06）~\cite{2606_05102} 可作为“紧凑前馈高斯预测”中的新锚点候选。相较于 feed-forward 3DGS，主要改变是：ZipSplat：token-based feed-forward模型，用多视图视觉token聚类生成紧凑scene tokens并解码为高斯组。3DGS 的显式高斯表示和可微渲染能力保持为基础。
其设计动机是：提出解耦Gaussian placement和像素网格的前馈模型，通过推理时聚类控制质量-速度折衷，减少平坦区域过量高斯。
具体而言，多视图backbone提取dense tokens，经k-means聚合为紧凑scene tokens，再用cross/self-attention细化并由MLP解码成任意3D位置的高斯组。
这一锚点判断来自当前分类中的机制关系：候选锚点：路线变化，且相关性 0.93；后续方法可围绕 用token聚类生成紧凑高斯集合 继续比较和扩展。
适用边界是：聚类质量影响高斯布局；复杂细节可能需要较高token预算；推理时聚类增加额外步骤。

\paragraph{\$\text\{VG\}\textasciicircum{}2\$GT}
\$\text\{VG\}\textasciicircum{}2\$GT（2026.06）~\cite{2606_01573} 延续 feed-forward 3DGS 在“前馈高斯预测”上的基本设定，主要改变是：VG2GT：Voxel-Gaussian Splatting Visual Geometry Grounded Transformer，冻结VFM并通过多尺度可微体素模块回归高斯参数。3DGS 的显式高斯表示和可微渲染能力保持为基础。
其设计动机是：提出体素-高斯视觉几何接地Transformer，将VFM特征、可微体素表示和高斯参数回归结合，并用随机实体体渲染监督深度。
具体而言，方法从图像特征构建多尺度体素几何表示，再拆分并回归高斯位置、尺度、颜色等参数；冻结VFM减少训练成本并提高泛化。
适用边界是：体素模块带来内存开销；对大规模场景和极稀疏输入的泛化仍需验证。

\paragraph{GADA}
GADA（2026.07）~\cite{2607_00595} 延续 3DGS 在“图像驱动Gaussian Splatting”上的基本设定，主要改变是：GADA：Geometry-Aware Deformable Aggregation，用可变形偏移迭代校正warped image错位，并进行几何感知特征聚合。3DGS 的显式高斯表示和可微渲染能力保持为基础。
其设计动机是：提出GADA模块，用可学习可变形偏移主动对齐由不确定几何造成的局部错位，将原本错开的源视角线索重新聚合到目标视角，从而改善warping-based Gaussian Splatting在细结构和高频区域的渲染质量。
具体而言，给定非结构化源图像和不确定几何，方法先渲染基础颜色与几何引导，再通过迭代refinement模块预测deformable offsets，对warped source features进行局部重采样和聚合；与标准像素级残差校正不同，GADA假设有用线索在轻微位移范围内仍局部保留，因此用几何感知偏移恢复被错位破坏的细节，并结合Gaussian Splatting渲染管线输出高保真新视角。
适用边界是：依赖初始几何质量和源视角覆盖；可变形聚合增加推理开销；大幅遮挡或几何严重错误时局部偏移可能不足。

\paragraph{FreeStreamGS}
FreeStreamGS（2026.06）~\cite{2606_03254} 延续 feed-forward 3DGS 在“在线前馈重建”上的基本设定，主要改变是：FreeStreamGS：在线前馈框架，包含解耦内参恢复头和流式位姿-几何一致机制。3DGS 的显式高斯表示和可微渲染能力保持为基础。
其设计动机是：提出可在线处理streaming unposed inputs的前馈3DGS框架，通过解耦内参恢复和长期一致性机制减少相机/几何误差累积。
具体而言，方法从连续输入中估计内参、位姿和高斯表示，并针对NVS对尺度、姿态和高斯参数一致性的敏感性设计解耦恢复模块。
适用边界是：仍依赖在线几何估计质量；长时间动态场景和闭环一致性可能是挑战。

\paragraph{StructSplat}
StructSplat（2026.06）~\cite{2606_28321} 可作为“Uncalibrated sparse-view Gaussian reconstruction”中的新锚点候选。相较于 generalizable 3DGS，主要改变是：结构化表示、像素对齐特征注入和语义感知几何建模。3DGS 的显式高斯表示和可微渲染能力保持为基础。
其设计动机是：提出组织 geometry、semantic 和 texture cues 的结构化表示，降低统一 backbone 中几何与外观纠缠。
具体而言，通过 pixel-aligned feature injection 改善纹理建模，并引入 semantic-aware 几何组件，在未知相机参数输入下前馈预测可渲染 3D Gaussians。
这一锚点判断来自当前分类中的机制关系：候选锚点：扩展，且相关性 0.95；后续方法可围绕 将 3D Gaussian 重建推进到无标定稀疏视图前馈场景。 继续比较和扩展。
适用边界是：无标定设定难度高，重建质量可能受相机估计、语义先验和输入视角覆盖限制。

\paragraph{Wild3R}
Wild3R（2026.06）~\cite{2606_11894} 延续 feed-forward 3DGS 在“野外稀疏照片集合”上的基本设定，主要改变是：Wild3R：基于WildCity大规模数据集学习外观一致性条件的feed-forward 3DGS。3DGS 的显式高斯表示和可微渲染能力保持为基础。
其设计动机是：提出WildCity数据集和相应前馈模型，使3DGS不再局限于干净、受控、多视角训练集。
具体而言，模型利用多视角、多光照和瞬态变化数据学习跨视角外观一致性，并从稀疏照片集合直接预测高斯场。
适用边界是：数据集构建成本高；极端天气/遮挡和长时间变化仍难处理。

\paragraph{L2D2-GS}
L2D2-GS（2026.06）~\cite{2606_29374} 延续 generalizable 3DGS 在“Learned densification for dynamic Gaussian scenes”上的基本设定，主要改变是：Learning to Densify：把通用重建改写为鲁棒迭代式 Gaussian densification。3DGS 的显式高斯表示和可微渲染能力保持为基础。
其设计动机是：将 feedforward Gaussian reconstruction 从一次性参数回归转为可学习的多步 densification，提高动态城市场景的可扩展性。
具体而言，在多视图动态输入下，逐步分配和细化 Gaussian primitives，使密集观测融合更稳健，兼顾高保真重建和比逐场优化更快的推理。
适用边界是：迭代致密化仍可能增加推理步骤；动态城市长尾对象和遮挡场景需要更多验证。

\paragraph{FastPano3D}
FastPano3D（2026.06）~\cite{2606_30352} 可作为“Panoramic feed-forward Gaussian reconstruction”中的新锚点候选。相较于 generalizable 3DGS，主要改变是：端到端从 equirectangular panorama 直接生成可渲染 3D Gaussian 表示。3DGS 的显式高斯表示和可微渲染能力保持为基础。
其设计动机是：针对 perspective 方法不适配 panorama 的问题，设计面向 equirectangular 图像的快速 Gaussian 生成框架。
具体而言，模型直接从单张全景输入预测可渲染 Gaussian scene representation，处理全景投影带来的空间非均匀特征分布。
这一锚点判断来自当前分类中的机制关系：候选锚点：扩展，且相关性 0.93；后续方法可围绕 将前馈 Gaussian 重建扩展到单张全景图输入。 继续比较和扩展。
适用边界是：单图重建存在未观测区域幻觉；全景图质量和室内先验对结果影响较大。

\paragraph{MooMIns -- Monocular 3D Reconstruction and Object Pose Estimation from Multiple Instances}
MooMIns -- Monocular 3D Reconstruction and Object Pose Estimation from Multiple Instances（2026.06）~\cite{2606_14389} 可作为“单图与对象级重建”中的新锚点候选。相较于 Gaussian Splatting，主要改变是：MooMIns：反转Gaussian splatting公式，单相机下渲染多个物体实例，结合SAM3 mask和修改SfM初始化。3DGS 的显式高斯表示和可微渲染能力保持为基础。
其设计动机是：提出工业bin场景中的多实例单图GS重建范式：不是多相机看一个场景，而是单相机看同一物体多个姿态实例。
具体而言，方法用SAM3分割实例mask，修改SfM初始化物体形状和实例位姿，再优化共享Gaussian物体模型和各实例6D姿态。
这一锚点判断来自当前分类中的机制关系：候选锚点：适配，且相关性 0.85；后续方法可围绕 反转多视图渲染设置，用多个实例提供隐含视角 继续比较和扩展。
适用边界是：需要同一对象多个可见实例；实例分割错误和遮挡会影响隐式多视角约束。
\section{无标定与自由视角}
\label{sec:13_2}

前馈 3DGS 的一个关键假设是输入图像具有已知的相机位姿，通常由运动恢复结构（Structure-from-Motion, SfM）计算获得。然而在实际场景中，获取精确位姿本身就是一个耗时而脆弱的过程，在弱纹理或低重叠区域尤其容易失败。本节综述的方法去除了这一假设，从无标定图像同时估计相机参数和三维高斯表示，使前馈重建适用于随意拍摄的图像集合。

\subsection{无位姿前馈重建}
\label{ss:13_2a}

本节讨论基于坐标规范化的无位姿前馈重建方法，通过将所有视角统一到规范坐标空间中消除对已知位姿的依赖。

\subsubsection{基于坐标规范化的无位姿方法}

本路线介绍通过坐标规范化消除位姿依赖的前馈重建方法。此前的前馈 3DGS 方法（如 pixelSplat、MVSplat）均遵循"局部到全局"（local-to-global）范式：先在各帧的局部相机坐标系中预测像素级高斯基元，再利用已知的相机外参 $[R^v \mid t^v]$ 将其变换到世界坐标系进行融合。这一范式存在两个固有缺陷：(1) 高斯变换依赖精确位姿，在稀疏视角和低重叠场景中尤为脆弱；(2) 变换-融合策略难以将来自不同视角的高斯整合为连贯的全局表示，在分布外数据上尤其容易失败。

NoPoSplat（2024.10）~\cite{2410_24207}对此提出了根本性的范式转变——不再"局部预测再变换"，而是让所有视角直接在同一个\textbf{规范高斯空间}（canonical Gaussian space）中预测高斯。具体做法是将第一个输入视角的相机坐标系作为全局参考帧，所有视角的高斯基元直接预测到这一统一坐标系中，完全绕过了局部到全局的坐标变换。形式化地，网络学习映射 $f_\theta: \{(I^v, K^v)\}_{v=1}^V \rightarrow \{(\mu_j^{v \rightarrow 1}, \alpha_j, r_j, s_j, c_j)\}$，其中上标 $v \rightarrow 1$ 表示第 $v$ 个视角中像素 $p_j$ 对应的高斯参数直接表达在第 1 帧的相机坐标系下，而非第 $v$ 帧的局部坐标系。第一个视角的位姿固定为 $[I \mid 0]$（单位旋转、零平移）。直觉上，这就像让所有画家面对同一幅画布各自画自己看到的部分，而不是每个人在自己的画布上画完再拼到一起。

在架构层面，NoPoSplat 采用纯 Vision Transformer（ViT）编码器-解码器，不注入任何几何先验（如对极约束或代价体）。各视角图像经 patch 嵌入后与内参 token 拼接，分别送入共享权重的 ViT 编码器；解码器通过跨注意力层实现多视角信息交互。参数预测采用基于 DPT 架构的双头设计：(1) \textbf{位置头}仅利用 Transformer 解码器特征预测高斯中心位置；(2) \textbf{属性头}在解码器特征之外额外引入 RGB 图像作为输入（RGB shortcut），确保纹理细节信息直接流向高斯基元，弥补 ViT 特征因 16 倍下采样而丢失的高频结构信息。

为解决无位姿设置下的\textbf{尺度歧义}问题，NoPoSplat 系统比较了三种相机内参编码策略：加法嵌入（将内参向量广播并加到图像特征上）、拼接 token（将内参编码为独立 token 与图像 token 拼接）、像素级球谐射线嵌入（将逐像素相机射线方向 $K^{-1}p_j$ 转换为球谐特征并与 RGB 拼接，不依赖外参）。实验表明内参 token 策略性能最优且设计优雅，因此作为默认方案。

NoPoSplat 的设计动机是多方面的：(1) 消除对位姿的依赖，将重建与标定解耦——这也是 NoPoSplat（No Pose Splatting）名称的由来；(2) 网络直接在规范空间中学习融合不同视角的信息，避免 transform-then-fuse 步骤带来的融合不一致；(3) 规范空间中的高斯一致性\textbf{反向解锁}了位姿估计应用——通过对齐不同视角预测的高斯，可直接估计输入视角间的相对位姿，无需深度真值或显式匹配损失即大幅超越先前最佳方法。训练采用端到端监督，损失函数为 $\mathcal{L} = \mathcal{L}_\text{MSE} + 0.05 \cdot \mathcal{L}_\text{LPIPS}$。

有趣的是，这一不含几何先验的纯 ViT 架构在低重叠视角场景下反而优于注入对极约束或代价体的专用骨干网络。原因在于，对极约束和代价体通常需要输入视角间有充分重叠才能有效——低重叠场景中这些几何先验反而成为负担。

\paragraph{几何置信度引导的三阶段精炼}

PF3plat（2024.10）~\cite{2410_22128}在 NoPoSplat 的坐标规范化框架上，针对稀疏视角、噪声位姿和低重叠等非理想条件下的训练不稳定性问题，提出了三阶段精炼策略。第一阶段，利用预训练单目深度估计和视觉对应模型实现跨视角高斯的粗对齐——为不同视角预测的高斯提供初始的空间一致性。第二阶段，引入轻量可学习模块从粗对齐出发精炼深度和位姿估计，提高三维重建和新视角合成质量。第三阶段，利用精炼后的估计计算\textbf{几何置信度分数}——评估每个高斯中心的可靠性——并据此条件化高斯参数预测：对低置信度区域的高斯施加更强的正则化约束（抑制其在训练中的梯度更新幅度），对高置信度区域则允许更灵活的参数调整。这一置信度感知策略有效区分了几何可靠与不可靠的区域，使模型在弱几何约束下仍能产生鲁棒的重建。在大规模真实世界数据集的广泛评估中，PF3plat 在所有基准上均达到最佳性能。

\paragraph{匹配感知位姿与深度精炼的自监督重建}

SelfSplat（2024.11）~\cite{2411_17190}实现了无位姿且无三维先验（无需预训练深度模型或对应模型）的可泛化重建——这一设置本质上是不适定的，因为缺乏真值监督、学到的几何信息和测试时微调。SelfSplat 的核心策略是将显式三维表示与自监督深度和位姿估计技术有效整合，使位姿精度和三维重建质量互相促进。具体而言，SelfSplat 设计匹配感知的位姿估计网络（matching-aware pose estimation network）——利用跨视角特征匹配的置信度加权位姿估计，避免低质量匹配引入噪声；同时引入深度精炼模块（depth refinement module）在初始深度预测的基础上，利用跨视角光度一致性约束增强几何精度。在 RealEstate10K、ACID 和 DL3DV 上，SelfSplat 在外观和几何质量上均超越先前最佳方法，并展现出强大的跨数据集泛化能力。

\paragraph{位姿感知规范体与锚点对齐预测}

SHARE（2025.05）~\cite{2505_22978}从另一个角度解决位姿不准确的问题——不依赖显式三维变换，而是提出联合形状和相机射线估计（shape-ray estimation）。其核心设计是构建位姿感知的规范体表示（pose-aware canonical volume representation），在规范坐标系中无缝整合多视角信息：每个空间位置的体素特征由所有输入视角按其估计的射线-体素交点加权聚合，位姿不确定性通过射线权重隐式建模——位姿越不确定的视角对体素特征的贡献越小。这一表示天然减少了不准确位姿引起的错位，因为体素特征是连续且视角加权的，而非依赖离散的坐标变换。此外，SHARE 引入锚点对齐高斯预测（anchor-aligned Gaussian prediction），在粗锚点周围精炼局部几何，允许更精确的高斯放置。在多种真实世界数据集上，SHARE 展现了鲁棒的无位姿泛化高斯 splatting 性能。

\subsection{基于基础模型的无标定方法}

通用前馈 3DGS 方法（如 pixelSplat、MVSplat）依赖已知相机内参计算相机射线，并据此将像素提升到三维空间。然而在无标定设置下，相机射线未知，这些方法无法直接适用。MASt3R 作为三维几何重建基础模型，已经能够从无标定图像对预测像素对齐的三维点云和置信度，但点云输出尚不足以支持高质量的新视角合成。

Splatt3R（2024.08）~\cite{2408_13912}走了一条与前述方法不同的路线——不从头训练前馈 3DGS 网络，而是直接在三维几何基础模型 MASt3R 的输出上追加高斯属性。其\textbf{核心洞察}在于：MASt3R 预测的像素对齐三维点与通用 3DGS 方法预测的像素对齐高斯基元在结构上高度相似——二者都基于跨注意力 Transformer 架构提取多视角信息，并为每个像素预测三维位置。因此，只需在 MASt3R 的点云输出上追加高斯基元所需的额外属性（协方差、不透明度、颜色），即可将其转化为完整的三维高斯表示，无需相机标定。

在具体实现上，Splatt3R 在冻结的 MASt3R 模型上新增第三条预测头——\textbf{Gaussian head}，与原有的点云预测头和特征匹配头并行运行。该头同样基于 DPT 架构，为每个像素预测：(1) 协方差，参数化为四元数旋转 $q \in \mathbb{R}^4$ 和缩放 $s \in \mathbb{R}^3$；(2) 球谐颜色参数 $S \in \mathbb{R}^{3 \times d}$；(3) 不透明度 $\alpha \in \mathbb{R}$；(4) 位置偏移 $\delta \in \mathbb{R}^3$。高斯均值参数化为 $\mu = x + \delta$，其中 $x$ 是 MASt3R 预测的三维点。所有参数均在第一帧的相机坐标系中表示，避免使用真值变换在不同参考帧间转换参数。各参数类型使用对应的激活函数：四元数归一化、缩放和偏移指数激活、不透明度 sigmoid 激活。为辅助高频颜色学习，模型预测像素颜色与高斯颜色之间的残差。最终的三维高斯集合是两幅输入图像预测的高斯的并集。

关于训练策略，Splatt3R 的设计体现了精心的考量与平衡。MASt3R 的点云预测为每个像素提供了从损失函数到正确三维位置的\textbf{直接梯度路径}，其具体形式为$L_\text{pts} = \sum_{v} \sum_i C_i^v \left(\frac{1}{z}\|X_i^v - \hat{X}_i^v\| - \lambda \log(C_i^v)\right)$，这使得高斯训练不会陷入局部最优（原始 3DGS 优化的核心痛点）。训练采用两阶段方案：第一阶段使用 MASt3R 预训练的点云几何损失提供稳定的三维初始化；第二阶段冻结 MASt3R 的其余参数，仅训练 Gaussian head，使用掩码重建损失优化新视角合成质量：

\begin{equation}
\mathcal{L} = \lambda_\text{MSE} \mathcal{L}_\text{MSE}(M \hat{I}, MI) + \lambda_\text{LPIPS} \mathcal{L}_\text{LPIPS}(M \hat{I}, MI)
\end{equation}

其中 $M$ 为\textbf{损失掩码}——通过将目标图像的像素反投影并重投影到上下文图像中，检查渲染深度是否与真值深度匹配来识别可见像素，仅对在至少一幅上下文图像中可见的区域计算渲染损失。这一掩码策略使 Splatt3R 能够训练外推到上下文视角之外的更远视点，而非仅限于视角间插值。

\paragraph{空间记忆网络的增量高斯场重建}

PreF3R（2024.11）~\cite{2411_16877}利用 DUSt3R 的成对三维结构重建能力，通过空间记忆网络（spatial memory network）扩展到变长度图像序列的多视角输入。其核心设计是将 DUSt3R 的成对重建能力以增量方式扩展：每当新图像到达时，空间记忆网络将当前帧的局部三维信息与已构建的全局场景表示进行对齐和融合，消除基于优化的全局对齐需求。PreF3R 同时引入稠密高斯参数预测头（dense Gaussian parameter prediction head），在 DUSt3R 的点图预测基础上追加协方差、不透明度和球谐颜色等高斯属性，使模型可监督以光度损失和点图回归损失的组合训练，同时增强光度真实感和结构精度。给定有序图像序列，PreF3R 以 20 FPS 增量重建三维高斯场，支持实时新视角渲染。

\paragraph{两阶段退火的流线型无位姿重建}

FreeSplatter（2024.12）~\cite{2412_09573}采用流线型 Transformer 架构，自注意力块促进多视角图像 token 之间的信息交换，在统一参考帧中解码像素对齐的三维高斯基元。在训练策略上，FreeSplatter 设计了两阶段方案：预训练阶段同时使用深度位置监督损失（将预测高斯位置与深度解投影真值对齐）和像素对齐余弦相似度损失（约束高斯分布在相机射线上），为高斯位置学习提供稳定的三维几何初始化；主训练阶段逐步退火位置监督并过渡到纯渲染损失加像素对齐约束，使模型最终以渲染质量为目标优化。FreeSplatter 发展了两个专用变体——面向物体级和场景级重建——分别在大规模数据集上训练。一个引人注目的结果是：FreeSplatter 不仅在重建质量上超越多个依赖位姿的大规模重建模型（LRMs），还能利用 PnP-RANSAC 从预测的高斯位置图高效估计相机参数，在位姿估计精度上与最先进的无位姿方法 MASt3R 可比甚至更优。此外，FreeSplatter 简化了文本/图像到 3D 的内容创建管线，消除了相机位姿管理的复杂性。

\paragraph{可微体素化与可学习相机 token 的全能重建}

AnySplat（2025.05）~\cite{2505_23716}通过单次前向传播同时预测三维高斯基元和相机内参外参，无需任何位姿标注即可扩展到随意拍摄的多视角数据集。在架构层面，AnySplat 基于 DINOv2 编码器和交替注意力 Transformer（alternating attention Transformer，在自注意力和跨视角注意力之间交替）提取跨视角特征，通过可学习相机 token（learnable camera tokens）和专用解码器预测相机参数（内外参），双头 DPT 解码器分别预测深度图和高斯属性。核心设计是可微分体素化模块——按高斯中心位置将逐像素基元聚类到体素网格中，以置信度加权的 softmax 归一化分配体素内权重并聚合属性，将基元数量从与像素数线性增长降至与场景复杂度相关的体素数，使模型可扩展到数百张输入视角。在大规模零样本评估中，AnySplat 在稀疏和稠密视角场景下匹配依赖位姿的基线方法质量，同时超越现有无位姿方法，大幅降低渲染延迟。

\subsubsection{多任务统一与语义重建}

本路线介绍将几何重建与语义理解统一到前馈 3DGS 框架中的多任务方法。前馈 3DGS 的初始目标仅是新视角合成，但实际应用中场景理解与重建往往需要同步完成。本节介绍 6 个方法：Uni3R、UniSem、UniSplat、FLEG、GSemSplat 和 VicaSplat，分别从联合建模、错误感知训练、双重掩码、紧凑语言嵌入、无标定语义场和视频帧重建等角度推进多任务统一。

\paragraph{开放词汇语义特征场的统一重建}

Uni3R（2025.08）~\cite{2508_03643}利用跨视角 Transformer（Cross-View Transformer）整合任意多视角输入，回归具有开放词汇语义特征场的三维高斯基元。其核心设计是将场景的外观和语义信息统一到同一组三维高斯基元中——每个高斯不仅携带颜色和几何属性，还嵌入高维语义特征，使其渲染出的语义特征可通过任意二维分割器解码为开放词汇分割结果。Uni3R 在单次前馈传播中同时实现三项任务：高保真新视角合成（RealEstate10K 上 25.07 PSNR）、开放词汇三维语义分割（ScanNet 上 55.84 mIoU）和深度预测，在多个基准上建立了新的最佳性能。

\paragraph{错误感知高斯丢弃与混合训练课程}

UniSem（2026.03）~\cite{2603_17519}从两个互补维度改进无位姿语义重建。首先，针对稀疏视角监督下高斯过度完整（over-complete）导致的几何不稳定和深度质量差的问题，提出错误感知高斯丢弃（Error-aware Gaussian Dropout, EGD）——利用渲染误差线索识别并抑制冗余倾向的高斯：渲染误差大的区域说明高斯配置不合理，EGD 以渲染误差为概率丢弃这些区域的部分高斯，迫使模型用更少的基元重建更准确的几何。其次，针对仅依赖二维分割器特征进行语义提升导致的三维语义不完整和泛化性差的问题，提出混合训练课程（Mix-training Curriculum, MTC）——训练初期使用二维分割器提升的语义作为监督信号，随着训练推进逐步将模型自身涌现的三维语义先验混入监督，配合物体级原型对齐（object-level prototype alignment）增强语义一致性和完整性。在 ScanNet 和 Replica 上，UniSem 在深度预测和开放词汇三维分割上均取得显著改进。

\paragraph{双重掩码与位姿重标定的几何-语义学习}

UniSplat（2026.04）~\cite{2604_10573}针对无位姿多视角输入下几何、外观和语义的联合表示学习，提出三个互补组件。(1) 双重掩码策略（dual-masking strategy）增强编码器的几何归纳偏置：同时掩码编码器和解码器的 token，并将解码器掩码定向到几何丰富区域，迫使模型从不完整的视觉线索中推断结构信息。(2) 粗到细高斯 splatting 策略逐步精炼辐射场，减少外观与语义之间的不一致性。(3) 位姿条件化重标定机制（pose-conditioned recalibration）强制几何-语义一致性——利用估计的相机参数将预测的三维点和语义图重投影到图像平面，与对应的 RGB 和语义预测对齐，确保不同预测头之间的跨任务一致性。三个组件共同产生对无位姿稀疏视角输入鲁棒的统一三维表示。

\paragraph{几何-语义双分支的紧凑语言嵌入}

先前的前馈语言嵌入高斯方法受限于固定输入视角数量，且为每个高斯附着语言对齐的语义嵌入，导致语义冗余和存储浪费。FLEG（2025.12）~\cite{2512_17541}提出几何-语义双分支蒸馏框架，支持无相机参数的任意多视角输入。其核心发现是：语义表示的空间稀疏性远高于几何表示——同一物体的大片区域共享相同的语义标签，不需要每个高斯都携带语言嵌入。基于此，FLEG 设计解耦语言嵌入策略（decoupled language embedding），仅用一组稀疏的语义高斯（semantic Gaussians）表示语言信息，而非为每个高斯都附加嵌入——仅需 5\% 的语言嵌入即可维持可比的语义保真度，有效降低存储成本。此外，FLEG 引入基于新视角的蒸馏策略（novel-view-based distillation），防止模型过拟合于输入视角。

\paragraph{双重特征的无标定语义场建模}

GSemSplat（2024.12）~\cite{2412_16932}在 Splatt3R 架构上构建，从稀疏无标定图像对学习与三维高斯关联的开放词汇语义表示。其核心挑战在于如何将二维语义监督有效提升到三维空间。GSemSplat 采用双重特征方法（dual-feature approach）作为二维空间监督——同时利用区域特定语义特征（region-specific features，捕获局部语义边界）和上下文感知语义特征（context-aware features，捕获全局语义上下文）的互补优势，使学习到的三维语义特征在细粒度和全局一致性之间取得平衡。在 ScanNet++ 上，GSemSplat 在开放词汇三维语义任务上超越了传统的逐场景优化方法。

\paragraph{可学习相机 token 的视频帧联合重建}

VicaSplat（2025.03）~\cite{2503_10286}从无位姿视频帧序列同时估计三维高斯和相机位姿。其架构创新在于可学习相机 token 的引入——图像编码器将每帧映射为视觉 token 序列，所有视角的视觉 token 与额外插入的可学习相机 token 拼接后，在定制的 Transformer 解码器中全交互。相机 token 以因果方式聚合不同视角视觉 token 的特征，并进一步逐帧调制视觉特征以注入视角依赖信息——这一设计使相机 token 成为编码全局位姿关系的瓶颈，视觉 token 则保持对局部场景内容的关注。三维高斯 splat 和相机位姿参数通过各自的预测头从增强后的 token 中估计。VicaSplat 不仅在多视角输入上超越基线方法，在跨数据集泛化（如零样本迁移至 ScanNet）上也展现出卓越性能。

\subsubsection{鲁棒性与特殊设置}

本路线介绍在噪声位姿、干扰物体、变化光照等真实世界挑战下保持鲁棒性的方法。真实世界的前馈重建面临噪声位姿、干扰物体、变化光照、非标准视角、全景输入等多种挑战。本节介绍 13 个应对上述鲁棒性问题的方法，涵盖几何一致性约束（G3Splat）、token 对齐预测（TokenSplat）、双专家分解（2Xplat）、基础模型几何先验适配（AirSplat）、任意输入统一模型（YoNoSplat）、长期在线重建（SaLon3R）、无约束光照（WildSplatter）、非有利视角适配（UFV-Splatter）、干扰物感知（DGGS）、全景图像适配（OmniSplat）、高效渲染（ERUPT）、无约束高斯 Surfel（SmileSplat）以及数据驱动的可扩展性分析（Wang et al.）等方向。

\paragraph{深度平滑与法向一致的几何自监督}

前馈无位姿方法主要依赖视角合成损失训练，但纯光度监督不足以恢复几何有意义的高斯——网络可以找到"欺骗"渲染损失的高斯配置（如半透明浮游体），在输入视角上渲染正确但在几何上毫无意义。G3Splat（2025.12）~\cite{2512_17547}系统分析并解决了自监督下学习三维高斯 splat 的歧义性问题，强制引入几何先验（如深度平滑性、法向量一致性）获得几何一致的三维场景表示。在 RealEstate10K 上训练后，G3Splat 在几何一致性重建、相对位姿估计和新视角合成三个任务上均达到最佳水平，并在 ScanNet 上展现出强大的零样本泛化能力。

\paragraph{Token 对齐与非对称双流解码器}

TokenSplat（2026.03）~\cite{2603_00697}引入 token 对齐的高斯预测模块（Token-aligned Gaussian Prediction），在特征空间直接对齐跨视角语义对应信息——通过粗略 token 位置和融合置信度引导，聚合多尺度上下文特征实现长程跨视角推理，减少重叠区域高斯的冗余。为增强位姿鲁棒性并解耦视角线索与场景语义，TokenSplat 采用可学习相机 token 和非对称双流解码器（Asymmetric Dual-Flow Decoder, ADF-Decoder）：相机 token 与图像 token 之间的通信被约束为单向的（方向性约束），维持前馈架构中相机估计与场景重建的清晰因子化，无需迭代精炼即可实现连贯重建和稳定位姿估计。

\paragraph{几何-外观双专家分解框架}

当前无位姿前馈 3DGS 普遍采用统一的单体架构，在单一网络中联合估计相机位姿和合成三维高斯。2Xplat（2026.03）~\cite{2603_21064}挑战了这一"全合一"范式，提出双专家框架——几何专家（专用位姿估计网络）预测相机位姿，外观专家（专用高斯生成网络）根据显式传递的位姿合成三维高斯。两个专家各司其职：几何专家专注于从图像对中恢复相对位姿，外观专家专注于从已知位姿的多视角图像预测高质量高斯。这一概念上简单但在先前工作中被忽视的设计极为有效：在不到 5000 次训练迭代内即大幅超越先前无位姿方法，性能与依赖位姿的最佳方法可比。这一结果表明，复杂的三维几何估计和外观合成任务可能更适合模块化设计，而非单一网络的端到端联合学习。

\paragraph{自洽位姿对齐与评分不透明度匹配}

三维视觉基础模型（3DVFMs，如 DUSt3R/MASt3R）在视觉几何估计上展现了强大的零样本能力，但直接将其用于高质量新视角合成仍然困难——基础模型预测的几何虽鲁棒但精度有限，且缺乏外观建模能力。AirSplat（2026.03）~\cite{2603_25129}提出训练框架将三维视觉基础模型的鲁棒几何先验适配到高保真无位姿新视角合成。其两项关键技术贡献为：(1) 自洽位姿对齐（Self-Consistent Pose Alignment, SCPA）——训练时的反馈回路，确保像素对齐监督能解决位姿-几何不一致性：先从当前预测渲染监督信号，检查渲染与真值的一致性，将不一致性反馈到位姿估计中修正；(2) 基于评分的不透明度匹配（Rating-based Opacity Matching, ROM）——利用稀疏视角 NVS 教师模型的局部三维几何一致性知识，过滤退化基元（opacity rating 低的高斯被抑制）。在大规模基准上，AirSplat 显著超越最先进的无位姿 NVS 方法。

\paragraph{混合训练与距离归一化的任意输入模型}

YoNoSplat（2025.11）~\cite{2511_07321}支持任意数量图像的输入，无论是否提供位姿和标定——模型预测每个视角的局部高斯和相机位姿，通过预测或提供的位姿聚合为全局表示。联合学习高斯和相机参数的核心困难在于两个任务的梯度纠缠。YoNoSplat 引入混合训练策略（mixing training strategy）缓解这一问题：初始阶段使用真值位姿聚合局部高斯（提供稳定的几何监督），逐步过渡到预测位姿与真值位姿的混合（防止训练不稳定和曝光偏差）。同时提出成对相机距离归一化（pairwise camera-distance normalization）解决尺度歧义——将相邻视角间的平移距离归一化，消除绝对尺度的不确定性；内参嵌入使模型能处理未标定输入（预测相机内参）。YoNoSplat 在 NVIDIA GH200 GPU 上仅需 2.69 秒即可从 100 张 $280 \times 518$ 分辨率图像完成场景重建。

\paragraph{显著性量化与三维点 Transformer 的长期在线重建}

现有在线前馈方法逐帧预测像素级高斯并简单拼接，在长视频序列中产生大量冗余和几何不一致。SaLon3R（2025.10）~\cite{2510_15072}首次实现 50 视角以上以超过 10 FPS 重建的在线前馈方法，冗余消除率达 50\%--90\%。其核心设计是紧凑锚点基元（compact anchor primitives）通过可微分显著性感知高斯量化（differentiable saliency-aware Gaussian quantization）消除冗余——先利用三维重建骨干预测密集像素级高斯和编码区域几何复杂度的显著性图，然后优先保留高复杂度区域的高斯，将冗余高斯压缩为紧凑锚点。三维点 Transformer（3D Point Transformer）从训练数据中学习空间结构先验，精炼锚点属性和显著性，支持区域自适应的高斯解码以保持几何保真度。无需已知相机参数或测试时优化，SaLon3R 在单次前馈中有效解决长序列重建的伪影和冗余问题。

\paragraph{联合外观嵌入的无约束变化光照}

WildSplatter（2026.04）~\cite{2604_21182}处理无约束照片集合中两个根本性挑战：未知相机参数和变化光照条件。其核心设计是联合学习三维高斯和外观嵌入（appearance embeddings）——每个高斯除几何和颜色属性外，还携带一个以输入图像为条件的外观嵌入，支持在不同光照条件下灵活调制高斯颜色。具体而言，外观嵌入与空间位置联合条件化颜色预测：同一三维位置在不同光照条件下通过不同的外观嵌入产生不同颜色，而解耦的外观嵌入确保光照变化不影响几何表示的一致性。WildSplatter 从稀疏输入视角在不到 1 秒内完成三维高斯重建，同时支持在多样光照条件下进行外观控制，在变化光照的真实世界数据集上超越现有无位姿 3DGS 方法。

\paragraph{LoRA 适配的非有利视角重建}

前馈模型通常假设输入来自有利视角（相机朝向场景中心），但实际应用中视角可能朝向任意方向。UFV-Splatter（2025.07）~\cite{2507_22342}提出适配框架，通过低秩适配（LoRA）层增强预训练无位姿模型处理非有利视角的能力——将输入图像重新居中（recenter）后送入预训练模型，LoRA 层以极少参数微调骨干适应新的视角分布。此外，高斯适配器模块（Gaussian adapter module）增强从重新居中输入导出的高斯的几何一致性，配合高斯对齐方法（Gaussian alignment method）为训练渲染准确的目标视角。整个框架仅需有利视角图像的数据集即可训练，无需收集非有利视角数据。

\paragraph{场景无关干扰物掩码与两阶段推理}

DGGS（2024.11）~\cite{2411_17605}首次解决干扰物（distractors，如行人、移动车辆等瞬态物体）存在下的可泛化 3DGS 重建问题。干扰物在跨场景泛化训练中引入三维不一致性和训练不稳定性。DGGS 提出场景无关的参考掩码预测和精炼模块（scene-agnostic reference-based mask prediction and refinement）——在训练阶段从参考视角预测干扰物掩码并从训练信号中剔除干扰物区域，有效消除干扰物对训练稳定性的影响。推理时，DGGS 通过两阶段推理框架：首先对参考视角评分和重新选择以获得最佳参考集，然后通过干扰物剪枝机制（distractor pruning）进一步移除残余的干扰物三维高斯基元。值得注意的是，DGGS 的可泛化掩码预测甚至在准确度上优于现有逐场景训练的专用干扰物检测方法。

\paragraph{捆绑调整的高斯 Surfel 联合优化}

SmileSplat（2024.11）~\cite{2411_18072}从无约束稀疏多视角图像预测像素对齐的高斯 surfel（二维高斯面元）。相比三维高斯，高斯 surfel 以更少的自由度表示表面元素，天然具备更好的多视角一致性。SmileSplat 设计多头高斯回归解码器预测 surfel 参数，并利用高质量法线先验增强 surfel 法向量的精度。在获得初始 surfel 后，通过捆绑调整高斯 splatting 模块（Bundle-Adjusting Gaussian Splatting module）联合优化高斯和相机参数（外参和内参），在端到端框架中同时提升几何精度和新视角渲染质量。

\paragraph{阴阳网格分解的全景图像适配}

现有前馈 3DGS 模型仅针对透视图像设计，全景图像的独特光学特性（如极不均匀的空间分辨率）导致特征编码器难以正确理解图像上下文，预测的高斯在空间中分布不均匀。OmniSplat（2024.12）~\cite{2412_16604}采用阴阳网格（Yin-Yang grid）分解将全景图像分解为多个准均匀子图像——阴阳网格的准均匀特性使分解后的子图像近似为透视图像，可直接利用已有前馈网络的强先验知识，无需额外训练。OmniSplat 在重建精度上超越在透视图像上训练的前馈网络，同时支持全景图像的编辑操作。

\paragraph{图块级查询与学习隐位姿的高效渲染}

ERUPT（2025.03）~\cite{2503_24374}引入基于图块的查询（patch-based querying）替代像素查询，大幅降低渲染目标视角所需的计算量——图块级查询将计算复杂度从像素级降低为图块级，使模型在训练和推理中均极为高效（在商用硬件上达到 600 FPS 渲染）。模型设计使用学习的潜在相机位姿（learned latent camera pose），允许在标注稀疏或不准确的位姿数据集上以无位姿目标训练。ERUPT 仅需 5 张无位姿输入图像即可高效渲染新视角，标注数据需求减少约 95\%，计算需求降低一个数量级。作者还引入了基于 Mapillary 街景图像构建的 MSVS-1M 基准数据集用于潜在视角合成评估。

\paragraph{数据驱动可扩展性分析}

Wang et al.（2025.06）~\cite{2506_09885}对前馈新视角合成方法进行了系统性分析，揭示了一个关键趋势：越少依赖显式三维知识的方法（如手工设计的 3D 表示和 SfM 位姿标注），在大规模训练数据上的性能加速越快，最终超越三维知识驱动的方法。作者将这一现象总结为"the less you depend, the more you learn"（越少依赖，学得越多）。基于此发现，Wang et al. 设计了完全去除显式场景结构和位姿标注依赖的前馈 NVS 框架，从大量二维图像中直接学习隐式三维感知——在训练和推理中均无需任何位姿信息。在大规模数据上训练后，该模型甚至在 NVS 性能上超越了依赖位姿训练数据的方法，验证了数据驱动范式的有效性。

\begin{takeaway}
无标定前馈重建的核心洞见是"越少依赖，学得越多"——去除对显式三维知识和位姿标注的依赖后，模型能更充分利用大规模训练数据，数据驱动范式正在取代传统几何管道。
\end{takeaway}

%=====================================================================
\section{单图到 3D 前馈方法}
\label{sec:13_3}

从单张图像重建三维场景是前馈重建中最具挑战性的设置。与多视角输入不同，单目设置面临固有的歧义性：深度不可知、背面信息缺失、尺度不确定。本节综述的方法通过生成先验、几何约束和蒸馏策略来补偿单目信息不足，实现从单张图像到三维高斯表示的前馈映射。

\subsection{单图前馈重建}
\label{ss:13_3a}

本节讨论两条技术路线：（1）对象级单图重建；（2）场景级单图重建。

\subsubsection{对象级单图重建}

本路线介绍从单张图像预测物体级三维高斯基元的方法。对象级单图重建聚焦于从单张图像预测物体级三维高斯基元，与 §\ref{sec:13_1}~节中的多视角前馈方法不同，这里面临的核心困难是深度信息完全缺失，背面几何和外观必须从有限的可见信息中推断。以下三个方法分别代表了 Mamba 序列建模、等变性约束和法线引导三条技术路线。

\paragraph{Mamba 线性序列建模与径向掩码约束}

Gamba（2024.03）~\cite{2403_18795}是首个将 Mamba 状态空间模型应用于单目 3DGS 重建的端到端前馈框架。与 §\ref{sec:13_1}~节中基于 Transformer 的 GS-LRM（二次复杂度）不同，Gamba 利用 Mamba 的线性序列复杂度直接预测 $L{=}16384$ 个三维高斯，避免了先前 LRM 系列（如 TriplaneLRM）所需的两阶段训练——第一阶段预测 4096 个高斯，第二阶段超分辨至 16384 个。

\textbf{架构设计。}Gamba 的核心是 GambaFormer，将参考图像通过 DINOv2 ViT 编码为条件 token $X \in \mathbb{R}^{K \times C}$，相机参数通过 MLP 映射为高维嵌入 $T \in \mathbb{R}^C$。关键创新在于 3DGS token 的构造：将参考图像与 Pl\"ucker 射线嵌入拼接为 9 通道输入，经 $p \times p$ 非重叠卷积生成特征图，再按四种预定义扫描顺序展平为长度 $L = 4hw$ 的一维序列，与可学习高斯嵌入相加后送入 Gamba 块。每个 Gamba 块在标准 Mamba 块基础上增加跨模态条件机制：将投影后的相机嵌入和图像 token \emph{前置}（prepend）到高斯序列头部，经 Mamba 处理后 \emph{移除}（remove），实现条件生成与线性复杂度的兼容。

\textbf{位置离散化。}高斯中心位置的预测采用离散化策略而非连续回归：将 $[-0.5, 0.5]^3$ 的坐标空间沿每个轴离散化为 21 个均匀点 $c_j$，通过线性投影和 softmax 获得逐坐标的概率分布，最终位置为各离散坐标的加权平均。这一设计比直接回归更稳定，避免了位置预测的退化。

\textbf{径向掩码约束。}与先前单目重建方法依赖三维点云监督防止模型坍塌不同，Gamba 提出纯图像空间的径向掩码约束：从图像中心向掩码轮廓发射射线构建距离场，对投影中心落在轮廓之外的高斯施加 MSE 约束损失 $\mathcal{L}_{\text{rdist}}$，从而将高斯约束在物体轮廓内。总训练损失为多视角渲染损失（MSE + LPIPS + alpha 掩码 MSE）与径向约束的加权和，配合渐进式训练策略（视角数从 2 逐步增至 6）。在 Objaverse 约 20K 物体的训练集上，Gamba 在 0.05 秒内完成重建，比 SDS 优化方法快约 1000 倍。

\paragraph{SO(2) 等变约束的规范立方体偏离预测}

SO(2)-等变高斯雕塑网络（GSNs）（2024.09）~\cite{2409_07245}从机器人操作的实际需求出发，将 SO(2) 旋转等变性引入单目 3DGS 重建——物体绕相机主轴旋转后，预测的高斯表示应相应旋转。这一性质对机器人抓取至关重要：物体在相机视野中的平面内旋转不应影响其三维几何预测。

\textbf{从规范立方体偏离的预测策略。}与 Gamba 从可学习 token 直接预测高斯参数不同，GSNs 在规范立方体上均匀放置 $N$ 个高斯（初始协方差 $\Sigma_0 = 0.0075\mathbf{I}$），网络仅预测偏离规范状态的增量——位置偏移 $\delta\mu_i$ 和尺度增量 $s_i$，而旋转 $r_i$、颜色 $c_i$ 和不透明度 $\alpha_i$ 直接预测。这种"从规范偏离"的设计与 GS-LRM 的射线距离参数化思路形成对比：GS-LRM 将高斯约束在射线上，GSNs 将高斯约束在规范立方体附近。

\textbf{共享编码器加并行解码器。}图像经 ResNet 骨干（取倒数第二层的 $2048 \times 7 \times 7$ 张量以保留更多空间细节）编码后，由五个独立的双层 MLP 并行解码五组高斯参数（位置偏移、尺度、旋转、颜色、不透明度），分别施加 ReLU+标量偏移（保证协方差正定）、归一化（单位四元数）和 Sigmoid（颜色和不透明度）等激活约束。

\textbf{扩展倒角距离与等变损失。}为训练等变性，输入图像绕主轴旋转随机角度 $\theta$ 后分别通过网络，预测的两组高斯之间计算扩展倒角距离（Extended Chamfer Distance, ECD）——在标准倒角距离的最近邻匹配基础上，对匹配点对额外施加颜色、尺度、旋转和不透明度的逐元素 $L_2$ 损失。总训练损失为 $L_{\text{image}} = L_2 + L_{\text{LPIPS}} + L_{\text{D-SSIM}}$（多视角渲染损失）与 $L_{\text{rot}} = \text{ECD}$（等变约束）之和。GSNs 在 Chairs 和 Cars 数据集上实现超过 150 FPS 的推理速度，并已验证可直接集成到机器人抓取管线中。

\paragraph{法线引导旋转确定与多尺度三平面外观注入}

GaussianPainter（2024.12）~\cite{2412_17715}提出了与前两者不同的输入假设——给定三维点云和参考图像，前馈地将点云"绘制"为高质量三维高斯。其核心洞察是高斯参数空间的\emph{非唯一性问题}：旋转、各向异性缩放和球谐系数的不同组合可以产生几乎相同的渲染结果，这使前馈网络的优化空间极度不规则，难以收敛。

\textbf{法线引导旋转确定。}GaussianPainter 的关键设计是分两步解决非唯一性。第一步，用 3D UNet（3D 稀疏卷积，批归一化替换为层归一化以适应点云稀疏性）将输入点云编码为 $200^3$ 占据网格的特征。第二步，法线解码器在\emph{各向同性}设置下（仅预测法线方向、各向同性尺度和不透明度）渲染法线图，在二维图像空间用法线图真值的 $L_1$ + SSIM 损失监督——这一设计巧妙地绕过了三维空间中偏离表面的点法线难以定义的困难。预测的法线 $\hat{n}$ 通过 Rodrigues 旋转公式直接确定高斯旋转 $R$：以 $\hat{n}$ 与世界坐标系 $z$ 轴的叉积为旋转轴、夹角为旋转角，将规范坐标系旋转到法线方向对齐的姿态。由此，高斯旋转从自由变量变为法线的确定性函数，大幅压缩了优化空间。

\textbf{多尺度三平面外观注入。}参考图像经 DINOv2 编码后与 3D 占据特征做交叉注意力获得 3D 增强特征，再构建多尺度三平面表示（每层尺度空间分辨率加倍）。各占据位置通过三平面双线性插值获取外观特征，注入到外观解码器的多个层级中，使高斯在几何确定后获得与参考图像一致的外观。在 ShapeNet 和 Objaverse 上的实验表明，法线引导将 PSNR 提升约 1.5 dB，多尺度三平面进一步贡献约 0.8 dB，两者缺一不可。

\subsubsection{场景级单图重建}

本路线介绍从单张图像重建完整三维场景的方法。场景级单图重建面临比对象级更大的挑战：场景的深度范围远大于物体、背面信息涵盖大范围背景、尺度歧义更为严重。以下五个方法分别从语言先验、教师-学生蒸馏、视频扩散潜在空间、几何引导生成和全景分解五个角度应对这些挑战。

\paragraph{VLM 上下文先验与空间先验辅助}

CATSplat（2024.12）~\cite{2412_12906}的核心观察是：单张图像无法提供足够的场景理解信息，但视觉语言模型（VLM）可以从同一图像中提取丰富的语义和上下文。CATSplat 引入两类互补先验来弥补单目信息不足。

\textbf{上下文先验。}输入图像经预训练 VLM 生成描述文本，文本经 CLIP 文本编码器映射为嵌入向量。在 3DGS 预测管线中，通过交叉注意力将场景特定的上下文细节融入文本嵌入——使网络不仅看到"像素长什么样"，还理解"这是什么场景、有什么特征"。这一设计与 §\ref{sec:13_1}~节中纯视觉特征的方法（如 DepthSplat、TranSplat）形成互补：那些方法依赖多视角几何约束，CATSplat 在仅有单视角时用语义约束替代。

\textbf{空间先验。}从输入图像预测的三维点特征提供几何引导：将点特征的 3D 空间信息注入高斯参数预测头，增强对场景几何布局的理解。在 RealEstate10K 上的实验表明，两类先验的联合使用使 CATSplat 在单目设置下超越纯视觉基线；跨数据集迁移至 NYUv2（室内）、ACID（自然场景）和 KITTI（驾驶场景）时，VLM 的语义泛化能力使 CATSplat 保持鲁棒性能。

\paragraph{多视角教师蒸馏与外推器网络}

studentSplat（2026.01）~\cite{2601_11772}解决的核心问题是\emph{尺度歧义}——从单张图像无法确定场景的绝对尺度，而直接用新视角监督训练单目模型会导致深度预测在任意尺度间摇摆。其设计灵感来自 §\ref{sec:13_1}~节中多视角方法（如 DepthSplat）已具备可靠的几何推理能力，但单目方法无法直接利用这些能力。

\textbf{教师-学生架构。}studentSplat 包含三个组件：(1) 多视角教师模型——以 DINOv2 为骨干的多视角 3DGS 网络，从多张输入图像预测可靠深度图；(2) 单目学生模型——以同一 DINOv2 骨干但仅接收单张图像，骨干分支沿 $z$ 轴定位高斯中心，细化分支用 CNN 特征精炼预测并解码其他高斯参数；(3) 外推器网络——预测缺失的场景上下文（输入视野之外的区域）。训练期间，教师模型为学生提供逐像素几何监督（深度损失），解决尺度歧义；外推器补充被遮挡区域的上下文信息。

\textbf{渐进式微调。}studentSplat 还展示了一种实用策略：先用单目模型从一张输入图像和虚拟小偏移位姿生成若干新视角，再将这些新视角作为多视角教师模型的输入进行进一步精炼，形成"单目预测$\to$多视角精炼"的级联管线。在 RealEstate10K 和 ACID 的外推设置下，studentSplat 首次使单目方法的场景级新视角重建质量接近多视角方法。

\paragraph{视频扩散潜在空间的大规模重建}

Wonderland（2024.12）~\cite{2412_12091}提出了一条与前述方法根本不同的路线——不直接从单张图像预测 3DGS，而是先用视频扩散模型生成多视角视频，再从视频的潜在表示中重建三维场景。这一思路的关键优势是：视频扩散模型在大规模视频数据上学习了丰富的场景先验（几何连续性、外观一致性、合理的遮挡补全），这些先验自然编码在视频潜在空间中。

\textbf{双分支相机引导视频生成。}Wonderland 以 CogVideoX-5B-I2V 为基础模型（生成 $480 \times 720$ 的 49 帧视频），引入双分支相机条件机制：主分支保留原模型的视频生成能力，ControlNet 分支（用前 21 个 Transformer 块初始化）注入相机轨迹条件。双分支设计确保精确的相机位姿控制，同时保持视频的静态一致性——背景不随相机移动产生时序抖动。生成的视频经 3DVAE 压缩为紧凑潜在表示（时间压缩比 $r_t{=}4$，空间压缩比 $r_s{=}8$，潜在维度 $13 \times 60 \times 90$）。

\textbf{潜在空间大规模重建模型（LaLRM）。}与传统 LRM 在图像像素空间操作不同，LaLRM 直接在视频潜在空间进行前馈三维重建。这一设计利用了视频潜在空间的紧凑性和三维感知性——潜在表示已将多视角信息编码为一致的结构。LaLRM 通过渐进训练策略（先在合成数据训练，再用自生成的野外数据微调）逐步提升重建质量。在 RealEstate10K 和 ACID 上的实验表明，Wonderland 从单张图像生成的三维场景在覆盖范围和视觉质量上均显著优于直接前馈方法。

\paragraph{平面感知度量深度引导的生成增强}

G4Splat（2025.10）~\cite{2510_12099}的核心论点是：准确几何是有效利用生成模型增强三维场景重建的\emph{前提}而非副产品——如果几何本身就不准确，生成模型产生的多视角图像会在错误的 3D 结构上"糊"出看似合理但几何上不一致的结果。

\textbf{平面感知几何建模。}G4Splat 基于一个强先验：室内外场景中存在大量平面结构（墙壁、地板、桌面）。通过多视角观测检测和合并平面掩码，为每个平面拟合三维平面方程，从中导出具有\emph{度量尺度}的深度图作为监督信号。这一设计与 §\ref{sec:13_1}~节中单目深度先验（如 DepthSplat 使用的 Depth Anything）的关键差异在于：单目深度仅有相对尺度，平面拟合提供绝对度量尺度，从根本上消除了尺度歧义。

\textbf{几何引导贯穿生成管线。}度量深度图在三处引导生成管线：(1) 改善可见性掩码估计——区分已观测和未观测区域；(2) 引导新视角选择——确保生成模型补全的区域在几何上连贯；(3) 增强视频扩散模型的多视角一致性——用几何约束过滤不一致的生成结果。基于 MAtCha 的基线模型在 Replica 和 ScanNet++ 上评估，G4Splat 在可见和未观测区域的几何精度与外观质量上均一致超越先前方法，并天然支持单目输入和未知位姿视频。

\paragraph{立方体面分解与逆深度拉普拉斯融合}

Genie Sim PanoRecon（2026.04）~\cite{2604_07105}面向机器人仿真中从单张全景图快速生成高保真 3D 场景的需求，提出全景分解-融合管线，在秒级时间内生成可用于仿真的三维场景。

\textbf{立方体面分解与并行处理。}将全景输入分解为六个不重叠的立方体面（cube-map faces），分别送入前馈 3DGS 网络（基于 SHARP）独立重建，再无缝拼接。这一设计利用了立方体面之间的独立性实现并行化，但引入了跨面深度不一致的挑战。

\textbf{深度感知融合。}为解决跨面深度尺度和接缝问题，Genie Sim PanoRecon 提出两阶段深度融合策略：先用 DA360 估计全景粗深度作为全局一致性参考，再通过逆深度拉普拉斯融合（inverse-depth Laplacian fusion）将 DA360 的全局一致性与 SHARP 的局部细节融合。此外，免训练深度注入模块（training-free depth injection）将融合后的深度图直接引导单目前馈网络的高斯位置预测，无需重新训练。整个系统在秒级内从单张全景图重建出照片级真实感的 3D 场景，已集成到 Genie Sim 机器人仿真平台中用于合成数据生成和仿真训练。

\begin{takeaway}
单图到 3D 前馈方法通过生成先验和几何蒸馏弥补单目信息不足的固有歧义，立方体面分解与逆深度拉普拉斯融合等策略使全景图到 3D 场景的秒级重建成为可能。
\end{takeaway}

%=====================================================================


\section{大规模重建模型}
\label{sec:13_4}

前馈 3DGS 的早期方法通常限于少量输入视角（2--4 张）和小范围场景重建。大规模重建模型（Large Reconstruction Models, LRMs）旨在将前馈重建扩展到数十甚至上百张输入视角，实现 360\textdegree{}场景级的全覆盖重建。核心挑战在于如何处理长序列输入带来的巨大 token 数量，同时保持重建质量和计算效率。

\subsection{逐像素高斯预测与长序列扩展}
\label{ss:13_4a}

本节讨论逐像素高斯预测与长序列扩展的大规模重建方法，将前馈重建扩展到数十乃至上百张输入视角。

\subsubsection{序列到序列映射的大规模重建}

本路线介绍将大规模重建建模为序列到序列映射的方法。

\paragraph{逐像素高斯预测的开创性工作}

GS-LRM（2024.04）~\cite{2404_19702}作为大规模重建模型的开创性工作，旨在解决先前 LRM 系列方法（LRM、Instant3D、DMV3D、PF-LRM 等）依赖固定分辨率 triplane 表示的根本局限——triplane 无法随输入分辨率自然扩展，难以应对大尺度变化和复杂场景。

\textbf{核心改动。}GS-LRM 用逐像素高斯预测取代 triplane，将前馈三维重建重新建模为序列到序列映射问题：多视角图像 token 序列经 Transformer 处理后，直接从每个输出 token 解码对应 patch 内的高斯参数。

\textbf{实现原理。}GS-LRM 的管线分为三个阶段——直觉上，它将三维重建转化为一个标准的"输入图像序列→输出高斯序列"映射，不需要手工设计几何推理模块。

（1）\textit{分块化}：对于 $N$ 张输入图像及其已知相机参数，首先计算每张图像的 \textbf{Pl\"ucker 射线坐标}——一种用方向向量 $\mathbf{d}$ 和矩向量 $\mathbf{o} \times \mathbf{d}$ 组成的六维参数化，编码每个像素对应的相机射线几何——作为位姿条件。由于 Pl\"ucker 坐标随像素和视角自然变化，无需额外的位置编码或视角嵌入。RGB 图像与 Pl\"ucker 坐标在通道维度拼接（9 通道），按 patch 大小 $p=8$ 切分，每个 patch 经线性层映射为 $d=1024$ 维 token。

（2）\textit{Transformer 处理}：所有 $N$ 张图像的 token 拼接为一个长序列，通过 $L=24$ 层 Transformer 块（Pre-LayerNorm + 16 头自注意力 + 两层 MLP，MLP 隐藏维度 4096，GeLU 激活，带残差连接）。采用 FlashAttention-v2 加速训练。

（3）\textit{高斯解码与参数化}：每个输出 token 经单个线性层解码为 patch 内的逐像素高斯参数（$q=12$ 维：3 通道 RGB、3 通道缩放、4 通道旋转四元数、1 通道不透明度、1 通道射线距离）。高斯中心位置由射线距离和相机参数计算：$\mathbf{x} = \mathbf{o}_{\mathrm{ray}} + t \cdot \mathbf{d}_{\mathrm{ray}}$。参数化设计体现工程考量——缩放初始化为 $\min(\exp(G_{\text{scale}} - 2.3),\, 0.3)$，其中 $-2.3$ 使初始值约为 $0.1$，$0.3$ 的上界防止产生"线状退化"高斯（过大的高斯会将梯度分散到多个像素，损害训练稳定性）；不透明度初始化为 $\sigma(G_{\text{opacity}} - 2.0)$，同样使初始值约为 $0.1$；旋转经 $L_2$ 归一化为单位四元数。最终输出为所有 $N$ 个视角的高斯并集，共计 $N \times H \times W$ 个高斯。

\textbf{设计动机。}逐像素高斯预测的设计有两个关键考量：其一，高斯数量与输入分辨率成正比，能够自适应地表达高频细节和大规模场景；其二，将三维重建简化为标准的序列到序列问题，无需设计复杂的几何推理模块。

\textbf{训练目标。}训练损失为 MSE 损失与感知损失的加权组合：

\begin{equation}
\mathcal{L} = \frac{1}{M} \sum_{i'=1}^{M} \left( \mathcal{L}_{\mathrm{MSE}}(\hat{I}_{i'}, I_{i'}) + \lambda \cdot \mathcal{L}_{\mathrm{Perceptual}}(\hat{I}_{i'}, I_{i'}) \right),
\end{equation}

其中 $M$ 为监督视角数，$\lambda=0.5$ 为感知损失权重——感知损失基于 VGG-19 特征而非 LPIPS，实验发现前者在大规模训练中更为稳定。

\textbf{训练与推理。}模型分两阶段训练：先在 $256\times 256$ 分辨率预训练 80K 步，再在 $512\times 512$ 微调 20K 步。物体级模型使用 Objaverse 730K 物体（4 输入视角 + 4 监督视角），场景级模型使用 RealEstate10K 80K 视频（2 输入视角 + 6 监督视角）。训练使用 64 张 A100 GPU，AdamW 优化器（$\beta_1=0.9$, $\beta_2=0.95$），余弦学习率衰减（峰值 $4\times 10^{-4}$）加 2000 步线性预热。推理时，模型含 300M 参数，在单张 A100 上仅需 0.23 秒即可从 2--4 张位姿图像预测高质量三维高斯。在 Objaverse 和 RealEstate10K 上分别训练后，物体级模型在 GSO 基准上超越 Instant3D 达 3.98 dB PSNR，场景级模型在 RealEstate10K 上超越 pixelSplat 达 2.2 dB PSNR。

\paragraph{长序列扩展}

Long-LRM（2024.10）~\cite{2410_12781}从根本上改变了 GS-LRM 的序列处理架构，用状态空间模型（State Space Model, SSM）的线性复杂度替代纯 Transformer 的二次复杂度，使前馈高斯重建从 2--4 视角扩展到 32 视角、分辨率从 $256\times256$ 提升至 $960\times540$——在单张 A100 GPU 上仅需约 1 秒即可完成。

\textbf{Baseline 瓶颈。}GS-LRM 的纯 Transformer 架构在扩展到 32 张 $960\times540$ 输入图像时遭遇根本性障碍：以 patch 大小 8 分块后，token 序列长度达到约 25 万——甚至超过 LLama3 的 128K 上下文长度。Transformer 自注意力的 $O(L^2)$ 复杂度意味着，在 32 视角 512 分辨率下，单次迭代（batch size 1）的训练时间即达 50.5 秒，内存和时间均不可承受。

\textbf{核心改动。}Long-LRM 提出三项关键设计：（1）用 Mamba2（一种高效的状态空间模型）块与 Transformer 块的 7:1 混合架构取代纯 Transformer；（2）在网络中间注入轻量级 token 合并模块，将序列长度压缩至 $1/4$；（3）通过 opacity 正则化配合训练/推理时高斯剪枝，将 1700 万逐像素高斯压缩至有效集合。

\textbf{实现原理：Mamba2-Transformer 混合架构。}理解混合架构的关键在于状态空间模型（SSM）的核心思想：与 Transformer 为每对 token 显式计算注意力权重不同，SSM 将整个历史信息压缩到一个固定维度的隐状态向量中，新 token 只需与当前隐状态交互即可获得上下文信息——因此计算量与序列长度成正比（$O(L)$），而非二次增长。具体而言，Mamba2 对输入 token $x_t$ 按递推公式处理：隐状态更新 $h_t = \mathbf{A}h_{t-1} + \mathbf{B}x_t$，输出 $y_t = \mathbf{C}h_t$，其中 $\mathbf{A}, \mathbf{B}, \mathbf{C}$ 由当前输入动态计算（而非固定参数），使模型能够根据输入内容自适应调整信息传递方式。Mamba2 在原始 Mamba 基础上进一步将状态矩阵 $\mathbf{A}$ 约束为标量乘单位矩阵（$a \cdot \mathbf{I}$），支持高效的块矩阵乘法和更大的状态维度（Long-LRM 使用 256 维状态、2 倍扩展率和 64 维注意力头），在多项语言任务上达到与 Transformer 相当的性能。

由于 Mamba2 原始的单向扫描（从左到右依次处理）是为语言任务设计的——文本天然具有因果方向性——而图像 token 没有固定的读取顺序，Long-LRM 改为双向扫描：对拼接的 token 序列分别进行前向和后向扫描，将两次输出求和后再经线性层投影。模型共 24 个块，每 7 个 Mamba2 块后跟 1 个全局自注意力 Transformer 块（使用 FlashAttentionV2 实现以优化 GPU IO），该模式重复 3 次，token 初始维度为 256。

\textbf{实现原理：Token 合并模块。}即使线性复杂度的 Mamba2 在最高分辨率下也会耗尽内存——25 万 token 的序列对任何架构都是严峻挑战。Long-LRM 受传统多级 CNN 编码器启发，在第 9 个块（即第二个混合块组的起始处）注入 token 合并模块：先将 token 序列从 $L \times D$ 重塑回空间形式 $N \times H' \times W' \times D$（$H'=H/p$，$p$ 为 patch 大小），施加 $2\times 2$ stride-2 的二维卷积，输出形状为 $N \times \frac{H'}{2} \times \frac{W'}{2} \times D'$，再展平为 $\frac{L}{4} \times D'$。合并后 patch 的有效大小从 8 扩展到 16，token 维度从 256 扩展到 1024——维度扩展补偿了空间分辨率的损失。这一设计将序列长度压缩至原来的 $1/4$，在保留空间结构的前提下使超高分辨率训练成为可能。

\textbf{实现原理：高斯剪枝策略。}32 张 $960\times540$ 图像产生约 1700 万个逐像素高斯，跨视角的重叠带来巨大冗余。直接剪除低 opacity 高斯（即使仅 5\% opacity 的那些）会导致场景中出现不完整的半透明物体——模型倾向于将一个不透明表面的颜色分散到多个半透明高斯中。Long-LRM 的策略分两步：首先，对所有高斯施加 opacity L1 正则化鼓励稀疏性：

\begin{equation}
\mathcal{L}_{\mathrm{opacity}} = \frac{1}{N} \sum_{i=1}^{N} |o_i|,
\end{equation}

其中 $o_i \in [0, 1]$ 为第 $i$ 个高斯的不透明度。L1 范数天然鼓励被约束项趋向稀疏，将有效高斯（opacity $>$ 0.001）占比从 99\% 压缩至约 40\%。然后剪除 opacity 低于 0.001 的高斯——实验证实该阈值下的剪枝不产生可感知的质量损失。训练时使用固定百分比剪枝（保留 top 40\% 加随机采样的 10\%）以保证恒定内存消耗，推理时使用阈值剪枝（保留 top 50\%）。

\textbf{实现原理：深度正则化。}大规模训练中一个典型的失败模式是产生"浮游体"（floaters）——不附着于真实 3D 表面的高斯，它们悬浮在相机前方充当半透镜片，试图通过视角相关的颜色混合来拟合训练图像，而非学习正确的表面位置和球谐系数。Long-LRM 引入 DepthAnything 预测的视差图作为软监督，采用尺度不变的深度损失：

\begin{equation}
\mathcal{L}_{\mathrm{depth}} = \frac{1}{M} \sum_{i=1}^{M} \mathcal{L}_{\mathrm{L1}}(\hat{D}^{\mathrm{DA}}_i, \hat{D}^{\mathrm{pred}}_i),
\end{equation}

其中视差图通过减去中位数 $t(d_i)$ 再除以平均绝对偏差 $\frac{1}{HW}|d_i - t(d_i)|$ 进行归一化，消除尺度和偏移的不确定性。总训练损失为：

\begin{equation}
\mathcal{L} = \mathcal{L}_{\mathrm{render}} + \lambda_{\mathrm{depth}} \cdot \mathcal{L}_{\mathrm{depth}} + \lambda_{\mathrm{opacity}} \cdot \mathcal{L}_{\mathrm{opacity}},
\end{equation}

其中渲染损失 $\mathcal{L}_{\mathrm{render}}$ 为 MSE 损失与 VGG-19 感知损失的加权组合（权重 0.5），$\lambda_{\mathrm{depth}}=0.1$，$\lambda_{\mathrm{opacity}}=0.01$。深度正则化引导浮游体高斯向真实表面移动，使模型转而通过球谐系数来表达视角相关的颜色变化，有效稳定大规模训练的收敛。

\textbf{实现原理：课程学习训练策略。}直接在 $960\times540$ 高分辨率上训练极其低效，Long-LRM 采用三阶段由低到高的课程学习策略。第一阶段：图像短边缩放至 256 后中心裁剪为正方形，随机采样 32 张输入和 8 张目标视角，以学习率 $4\times10^{-4}$ 训练 60K 步（batch size 256）。第二阶段：缩放至 $512\times512$，学习率降至 $4\times10^{-5}$，训练 10K 步。第三阶段：缩放至 $960\times540$（不做正方形裁剪），视角采样范围扩展到整个视频序列（DL3DV 约 200--300 帧），训练 10K 步并启用高斯剪枝。

\textbf{设计动机。}每一项设计选择都有明确的实证支撑。混合架构的选取基于关键观察：纯 Mamba2 虽然线性复杂度适合长序列，但完全移除自注意力会严重损失全局建模能力——消融实验显示，纯 Mamba2 在 4 视角 256 分辨率下比混合模型低 1.8dB PSNR，在 32 视角 512 分辨率下差距扩大至 2.5dB，这源于 SSM 基于状态的递推设计难以捕捉长程依赖。7:1 的 Mamba2-Transformer 比例在质量和效率间取得最优平衡；将 Transformer 块均匀分布在 Mamba2 块中（而非集中在开头或末尾）使训练更稳定。Token 合并模块在混合架构基础上进一步将超高分辨率训练时间压缩至 $2/3$，且初始维度 256→合并后 1024 的维度扩张策略额外节省了 $1/3$ 的训练时间和内存——这是唯一能在 32 视角 $960\times540$ 设置下不 OOM 的方案。Opacity L1 正则化的设计则是为了解决逐像素预测带来的固有冗余——模型倾向于将不透明颜色分散到多个半透明高斯中，简单阈值剪枝会导致不完整重建。

其后续工作 Long-LRM++（2025.12）~\cite{2512_10267}在 Long-LRM 的基础上将逐像素颜色高斯替换为半显式的"特征高斯"表示，在保持前馈效率的同时显著提升细节重建质量。核心改动在于：每个图像 token 预测 $K$ 个自由移动的特征高斯（而非射线对齐的颜色高斯），每个高斯携带 $F$ 维特征向量而非 SH 颜色系数。渲染时，特征高斯通过 splatting 将特征渲染到目标视角的特征画布上，再与目标相机的 Pl\"ucker 射线坐标拼接后送入轻量级 Transformer 解码器（平移不变局部注意力捕获高频细节 + 全局注意力建模长程依赖），最终线性层输出 RGB 颜色。这一设计的核心优势在于：(1) 特征高斯的自由移动特性提供比像素对齐更强的鲁棒性和表达力——无需精确对齐像素即可编码场景信息；(2) 解码器的隐式表达能力弥补了显式高斯在细粒度纹理上的不足。此外，Long-LRM++ 引入多空间分割合并机制：将高斯分为多组独立渲染解码，再通过自适应加权融合各组结果，有效处理镜面反射等非物理几何场景。实验表明，Long-LRM++ 在 DL3DV 上匹配隐式方法 LaCT 的渲染质量，仅需 Long-LRM 1/4 的高斯数量（约 120K vs.\ 480K），实现实时 14 FPS 渲染，并可扩展到 64 张输入视角。超参数分析显示，每 token 预测的高斯数 $K$ 和特征维度 $F$ 对质量影响显著——增大 $K$ 或 $F$ 允许模型存储更多场景信息，但计算开销线性增加。

\subsection{迭代精炼与可扩展架构}
\label{ss:13_4b}

本节讨论通过解耦表示与迭代精炼机制提升大规模重建质量的方法。

\subsubsection{解耦表示与迭代精炼架构}

本路线介绍通过解耦表示设计与迭代精炼机制提升大规模重建质量的方法。

\paragraph{迭代精炼与测试时训练}

Long-LRM++ 的半显式设计通过解码器增强表达能力，但整体仍是单次前馈推断。以下两个方法从不同角度引入迭代/递推机制进一步提升大规模重建质量。

iLRM（2025.07）~\cite{2507_23277}提出与 GS-LRM/Long-LRM 的像素对齐范式根本不同的架构——将三维高斯表示与输入图像\emph{解耦}，通过迭代精炼逐步构建场景。其三个核心设计原则分别对应大规模重建的三个瓶颈：

\textbf{（1）解耦表示。}iLRM 不在输入图像分辨率上直接预测高斯参数，而是初始化低分辨率的\emph{视角嵌入}（viewpoint embeddings）$V_i^{(0)} \in \mathbb{R}^{H^v W^v / p^2 \times d}$（$H^v, W^v$ 远小于输入分辨率 $H, W$）作为三维场景的基础表示，再通过线性层逐 token 解码为高斯参数。这一设计使高斯数量与输入图像分辨率无关，计算量随视角数线性而非二次增长。

\textbf{（2）两阶段注意力。}每个迭代更新块由交叉注意力 + 自注意力组成。交叉注意力层以视角嵌入为查询、对应视角的图像 token 为键值，每步注入视觉信息——与 GS-LRM 的全局自注意力（所有视角 token 混合计算）不同，iLRM 的交叉注意力是逐视角的，计算复杂度与视角数成正比。自注意力层则让所有视角嵌入交互，建立全局一致性。以 16 张 $256 \times 256$ 输入、patch 大小 8 为例，相对计算量比 GS-LRM 的全局注意力降低至约 8\%（1:0.08）。

\textbf{（3）Token 提升与 Mini-batch 交叉注意力。}低分辨率视角嵌入难以充分利用高分辨率图像 token 的信息。iLRM 引入 token 提升策略：每个 LR 视角 token 经线性层扩展维度 $k$ 倍后重塑为 $k$ 个细粒度查询 token（$k{=}2$），与 HR 图像 token 做交叉注意力后再合并回 LR 维度。此外，受优化中 mini-batch 梯度下降的启发，iLRM 设计了结构化的 token 采样策略，在交叉注意力中仅采样图像 token 和视角 token 的子集，进一步降低计算开销。

iLRM 在 RealEstate10K 和 DL3DV 上训练，支持从 2 到 16+ 视角的灵活输入，在宽基线和低重叠场景中表现尤为出色——这是因为解耦表示不依赖逐视角的像素对齐，对视角间重叠度更鲁棒。

tttLRM（2026.02）~\cite{2602_20160}从另一个维度突破输入规模限制——引入测试时训练（Test-Time Training, TTT）层，将多个图像观测高效压缩到 TTT 层的\emph{快速权重}中，形成潜在空间中的隐式三维表示。其核心思想与 Long-LRM 的 Mamba2 状态空间模型有本质区别：Mamba2 将历史信息压缩到固定维度的隐状态中（递推式更新），而 tttLRM 的 LaCT（Large Chunk Test-Time Training）层将输入 token 的键-值对视为训练数据，用 Muon 优化器在线更新快速权重矩阵 $W$（梯度下降式更新），能够一次处理多达百万 token 的大块输入。

\textbf{前馈重建模式。}模型由 24 个 LaCT 块组成（隐藏维度 768），输入图像与 Pl\"ucker 射线嵌入拼接后分块化为 token。LaCT 块包含窗口注意力（捕获视角内局部关系）和快速权重更新/查询操作（均为线性复杂度）。虚拟视角 token 查询快速权重获取更新后的表示，再经线性解码器预测逐像素高斯参数。训练损失为 MSE + VGG-19 感知损失，配合深度正则化和 opacity 正则化（与 Long-LRM 类似）。

\textbf{自回归重建模式。}tttLRM 的独特能力在于支持流式输入——新视角以 mini-batch 方式逐批到达，快速权重增量更新，每步立即预测当前的三维高斯。这一模式模拟了机器人在环境中逐步探索并实时构建三维地图的场景。实验表明，tttLRM 在 DL3DV 场景级重建中超越 Long-LRM 的前馈质量，且不经后优化时的 PSNR 已高于 Long-LRM 经 3 步后优化的结果。分布式序列并行策略允许多 GPU 分片处理长序列，进一步扩展输入规模。

\paragraph{MVP: 双层次注意力的大规模可扩展架构}

从 GS-LRM 的 2--4 视角到 Long-LRM 的 32 视角再到 tttLRM 的流式处理，大规模重建的输入规模持续增长。以下四个方法从可扩展架构、紧凑表示和稠密视角推理三个角度进一步推进。

Multi-view Pyramid Transformer（MVP）（2025.12）~\cite{2512_07806}提出与 GS-LRM 的纯全局注意力和 Long-LRM 的 Mamba2-Transformer 混合架构不同的第三条路线——双层次注意力架构，在一次性前馈推断中处理 100+ 张输入图像。

\textbf{视角间层次（Inter-view Hierarchy）。}MVP 将多视角交互分为三个阶段：Stage 1 仅帧内自注意力提取空间特征；Stage 2 引入分组注意力——将 $N$ 个视角按时序邻近度分为 $N/M$ 组（每组 $M$ 个视角），组内执行全局自注意力，兼顾局部一致性和计算效率；Stage 3 将所有视角合并为单一组执行全局自注意力，建立场景级全局一致性。这一局部$\to$分组$\to$全局的渐进式交互使注意力计算量大幅降低——越早期的阶段 token 数越多但交互范围越受限。

\textbf{视角内层次（Intra-view Hierarchy）。}空间维度上，每个阶段之间通过卷积层将空间分辨率减半（$h/2, w/2$），token 维度加倍，形成从细到粗的特征金字塔。这一设计与 §\ref{sec:13_1}~节中 DepthSplat 的层次化匹配架构思路类似，但 MVP 的金字塔用于\emph{多阶段注意力}而非深度估计。

\textbf{金字塔特征聚合与解码。}三个阶段的输出通过自顶向下的精炼过程融合：最高阶段的粗粒度全局特征上采样后与中间阶段特征拼接，逐级精炼，最终聚合特征经线性层解码为逐像素高斯参数（包含 SH 系数以建模视角相关外观）。MVP 还引入 opacity 正则化 $R_\alpha = \frac{1}{N_G}\sum_{j=1}^{N_G}|\sigma(\alpha_j \cdot s_j)|$（$s_j$ 为随机方向的 SH 基），抑制跨视角的 opacity 伪影。在 DL3DV 上的实验中，MVP 在 16/32/64/128 视角设置下一致超越 Long-LRM 和 iLRM，推理时间在单 GPU 上不到 1 秒（128 视角输入），零样本迁移至 Tanks and Temples 和 Mip-NeRF 360 时同样展现优异泛化能力。

\paragraph{C3G: 可学习查询 token 的紧凑场景表示}

MVP 通过双层次注意力提升了可扩展性，但在高斯数量上仍采用密集像素对齐。C3G（2025.12）~\cite{2512_04021}从表示效率的角度出发，提出与 GS-LRM/Long-LRM 的密集像素对齐高斯根本不同的紧凑范式——仅用 $N{=}2048$ 个可学习查询 token 经 Transformer 自注意力聚合多视角特征后，每个 token 解码为一个三维高斯。与 §\ref{sec:13_1}~节中 TokenGS 的可学习 token 方案类似，C3G 的核心优势是高斯数量与输入分辨率完全解耦，仅 2K 高斯即可表征完整场景。C3G 采用 VGGT 作为默认视觉编码器，训练中引入渐进低通滤波策略（RAIN-GS，从 $s{=}300$ 退火至 $s{=}0.3$）稳定初期位置预测。令人惊讶的是，仅通过光度重建损失训练，C3G 的查询 token 自发学会关注\emph{几何有意义的区域}——注意力图显示每个 token 精准对齐到多视角中的同一物体区域，形成了无监督的跨视角对应关系。基于这一涌现特性，C3G 进一步提出 Any-feature 3D 提升方法：冻结已学到的注意力权重，仅训练新的值投影层，即可将任意 2D 特征（如 DINO、CLIP 特征）提升为多视角一致的 3D 特征，实现场景理解与新视角合成的联合任务。在 DL3DV 上的实验中，C3G 在仅用 2048 个高斯的情况下实现了有竞争力的新视角合成质量，同时显著降低了内存开销。

\paragraph{VGGT-X: 稠密视角扩展的内存优化与位姿精炼}

C3G 展示了紧凑表示在大规模重建中的潜力。VGGT-X（2025.09）~\cite{2509_25191}将三维基础模型 VGGT 从稀疏视角扩展到稠密视角新视角合成，解决两个关键挑战：计算可扩展性和位姿精度不足。在可扩展性方面，VGGT-X 提出三项内存优化：(1) 移除 VGGT 24 层中仅第 4、11、17、23 层用于密集预测的冗余中间输出（VGGT$-$），将图像吞吐量从 150 提升至 600；(2) 除 MLP 头外全面切换为 BFloat16 精度（VGGT$--$），峰值 GPU 内存降低 74\%；(3) 异步帧处理——将 $N$ 张输入分为 $N/S$ 个 chunk 顺序处理，控制峰值内存。在位姿精度方面，VGGT-X 设计自适应全局对齐：先用 XFeat 特征匹配获取图像对应关系，再以极线距离的直方图密度估计自适应加权 $w_k = (f(e_{m,k}) / \text{Avg}(f(e_{m,k})))^\beta$（$\beta{=}0.5$），根据初始位姿精度的中位数自适应调整学习率。3DGS 训练阶段采用 MCMC-3DGS（比标准 3DGS 对噪声位姿更鲁棒）并联合优化残差相机参数。在 Mip-NeRF 360（最长 311 帧）、Tanks and Temples（最长 1106 帧）和 CO3Dv2 上，VGGT-X 将无 COLMAP 稠密新视角合成的性能大幅提升至接近有 COLMAP 初始化的水平。

\paragraph{UniQueR: 混合查询的遮挡感知稀疏推断}

VGGT-X 专注于稠密视角的扩展与优化。UniQueR（2026.03）~\cite{2603_22851}将重建建模为稀疏三维查询推断问题，与 GS-LRM/Long-LRM 的像素对齐范式和 C3G 的可学习 token 范式不同，UniQueR 的核心创新在于使网络能推断包括遮挡区域在内的场景结构——用比密集替代方案少一个数量级的基元实现更优的渲染质量和几何精度。

\textbf{混合查询初始化。}与 C3G 的纯可学习 token 不同，UniQueR 采用混合策略：一半查询从 DINOv2 编码器加交替注意力 Transformer 预测的非度量点图中采样，提供与观测表面对齐的粗几何结构；另一半作为可学习锚点均匀分布在三维空间中，用于探索和重建未被直接观测到的区域。这一设计使 UniQueR 在遮挡补全上具有天然优势——随机初始化的查询在新视角监督的压力下被鼓励放置在输入视角之外的区域。

\textbf{解耦注意力与高斯生成。}查询 Transformer 通过解耦注意力（先查询到图像 token 的交叉注意力，再查询间的自注意力）将计算复杂度从 $O((Q+L)^2)$ 降至 $O(QL+Q^2)$。每个查询预测一个形变偏移 $\delta q_i$ 和 $K$ 个高斯偏移 $\delta g_{ik}$，以残差方式生成 $Q \times K$ 个高斯基元：$G = \{p_i + \delta q_i + \delta g_{ik}\}_{i=1,k=1}^{Q,K}$。训练中，监督视角是输入视角的超集（如 3 张输入 + 3 张 held-out），强制查询学习推断未观测区域的几何。

在 Mip-NeRF 360 和 VR-NeRF 上的实验表明，UniQueR 在渲染质量和几何精度上均优于像素对齐方法（AnySplat）和可学习 token 方法（C3G），且预测的高斯可直接作为标准 3DGS 逐场景优化的初始化，进一步提升性能。

\begin{takeaway}
大规模重建模型将前馈 3DGS 从小场景扩展到 360 度全景重建，关键突破在于混合查询初始化和解耦注意力机制，使网络能以更少基元推断包括遮挡区域在内的完整场景结构。
\end{takeaway}

%=====================================================================


\section{特定领域前馈方法}
\label{sec:13_5}

前馈 3DGS 的通用范式可以适配到多种特定领域，本节综述事件相机、射频信号、人体渲染、体积可视化、视觉语言导航等场景中的专业方法。

\subsection{特殊传感器与信号建模}
\label{ss:13_5a}

本节讨论三条技术路线：（1）事件相机与特殊传感器；（2）射频信号建模；（3）人体渲染与通信。

\subsubsection{事件相机与特殊传感器}

本路线介绍将前馈 3DGS 适配到事件相机等特殊传感器的重建方法。EvGGS（2024.05）~\cite{2405_14959}是首个基于事件相机的可泛化三维重建框架，从纯事件输入前馈重建场景为三维高斯。事件相机输出异步事件流（像素级亮度变化信号），与传统相机的 RGB 帧有本质差异——事件数据稀疏、无绝对颜色信息。EvGGS 设计了三级级联架构来逐步弥合这一模态鸿沟：深度估计模块（UNet 骨干从事件流预测深度图）、强度重建模块（以深度特征为辅助输入从事件流重建灰度图）和高斯回归模块（从深度+强度特征预测三维高斯参数）。三个模块在 3DGS 渲染框架下协作训练——渲染损失同时回传到三个模块，使深度估计获得三维感知、强度重建获得几何引导、高斯回归获得更好的输入特征——形成互相促进的正反馈。为适应事件数据的特殊性，EvGGS 将标准 3DGS 的球谐系数替换为单一灰度值（事件数据无彩色信息）。作者还构建了 Ev3DS 数据集（合成+真实事件数据，含多种材质物体及灰度图、深度图、相机位姿等标定标签），渲染速度达 195 FPS（比 EventNeRF 的 0.045 FPS 快约 4000 倍）。

\subsubsection{射频信号建模}

本路线介绍将 3DGS 扩展到射频（RF）频段的信号建模方法。将 3DGS 从可见光扩展到射频（RF）频段，需要重新定义"颜色"和"渲染"的含义——RF 信号无形无色，其观测量是接收信号强度（RSSI）、信道状态信息（CSI）或功率角度谱（PAS）等物理量。以下两个方法分别解决跨频段泛化和接收器泛化两个核心问题。

XFreq-GS（2026.05）~\cite{2605_11432}解决跨频段无线辐射场重建问题——传统 RF-3DGS 方法仅在单一载波频率上建模，而实际无线系统（如 6G 网络）需要在 1--94 GHz 的宽频率范围内进行信道建模。XFreq-GS 的核心设计是频率解耦：三维高斯基元共享几何属性（位置、协方差），但配备频率自适应的 RF 属性——每个高斯的"辐射"不再是 RGB 颜色，而是可随载波频率变化的电磁传播参数。渲染时，3DGS 的 $\alpha$-混合被重新解释为 RF 信号的角度域叠加，输出功率角度谱（PAS）地图。在包含 14,460 个样本（1,446 个发射器位置 $\times$ 10 个载波频率）的数据集上，XFreq-GS 在 PAS 重建和未见频率泛化上均一致优于 WRF-GS+ 等 RF-3DGS 基线。

RxGS（2026.05）~\cite{2605_24290}解决接收器可泛化的 RF 数据合成问题——先前 RF-3DGS 方法为每个接收器训练独立模型，$N$ 个接收器需要 $N$ 个模型且无法预测训练中未见过的接收器位置。RxGS 的关键洞察来自一个数学证明：在 3DGS- RF 渲染框架中，接收器位置仅通过辐射系数（radiance coefficients）影响渲染结果，所有其他属性（几何）均与接收器无关。基于这一分解，RxGS 采用两阶段设计：第一阶段学习接收器无关的共享三维高斯几何；第二阶段冻结几何，学习以接收器位置为条件的方向辐射——通过互补的全局模块（捕获平均传播效应）和局部模块（建模逐散射体的残差）实现。单个统一模型即可支持多个接收器（BLE RSSI、RFID、WiFi CSI 三种模态），泛化到训练中未见过的接收器位置，训练、推理和存储成本降低达一个数量级。所有实验在单张 NVIDIA H100 GPU 上完成。

\subsubsection{人体渲染与通信}

本路线介绍前馈 3DGS 在人体渲染与实时通信中的应用方法。人体渲染与前馈 3DGS 的交叉需要解决特殊挑战：人体的非刚性形变、稀疏视角下的深度歧义、以及对实时性（视频通信）和真实性（不能使用 enrolled 固定外观）的双重需求。

GPS-Gaussian+（2024.11）~\cite{2411_11363}提出可泛化的像素级高斯 splatting 方法，在源视角上定义高斯参数映射并直接回归高斯属性，无需任何微调或优化即可实现即时新视角合成。该方法在人体和人体-场景数据上训练，核心设计包括：(1) 共享图像编码器从两幅邻近源视角提取特征；(2) 极线注意力机制（epipolar attention）沿极线搜索跨视角对应关系估计深度，在深度监督缺失时保持几何一致性——这避免了 DenseSplat 等方法对 MVS 代价体的依赖；(3) 深度残差模块在缺少深度监督时细化初始深度预测。在 DyNeRF 和 ENeRF-outdoor 数据集上的实验表明，GPS-Gaussian+ 在稀疏 4--6 视角设置下实现实时渲染，并已部署到 10 相机实时演示系统中。

Splat-SAP（2025.11）~\cite{2511_22704}利用像素级 point map 重建表示几何，通过两阶段学习策略解决大稀疏双目输入的挑战。第一阶段基于 MASt3R 的 ViT 骨干提取双目图像 token，通过迭代亲和力学习将尺度不变的 point map 转换到真实尺度空间——这一步骤无需 3D 监督，仅通过渲染损失自监督训练。第二阶段将初始 point map 投影到目标视角形成高斯平面锚点，通过立体匹配约束精炼几何，在精炼后的几何上构建高斯基元进行渲染。整个粗到细管线仅用渲染损失训练，泛化到工业相机、GoPro 和手机三种不同相机类型。

VoluMe（2025.07）~\cite{2507_21311}首次从单个 2D 网络摄像头实时预测三维高斯重建，目标是实现真实的 3D 视频通话——每帧独立条件化确保重建忠实再现输入视频（"真实性"属性），同时泛化到新视角产生真实感结果。模型基于 Splatter Image 的 UNet 架构（减小骨干以适配实时性），但引入多项改进：(1) 优化通道选择——只预测对重建质量贡献最大的高斯属性通道，减少冗余预测；(2) 抖动损失（jitter loss）抑制帧间时序闪烁；(3) 单应性（homography）裁剪替代 CLIFF 参数裁剪，提升面部区域的预测稳定性。在 Cafca、Ava-256 和 FFHQ 三个多视角人脸数据集上，VoluMe 在保持真实性的前提下实现了有竞争力的重建质量，并已在商用硬件（RTX 3090 + 单个网络摄像头）上部署为实时 3D 视频通话原型。

\subsection{视觉应用与联合优化}
\label{ss:13_5b}

本节讨论三条技术路线：（1）体积可视化与物体表示；（2）视觉语言导航；（3）光流-深度联合优化。

\subsubsection{体积可视化与物体表示}

本路线介绍将前馈 3DGS 扩展到科学可视化与通用物体表示的方法。前馈 3DGS 的对象通常是自然场景的 RGB 图像，以下两个方法将范式扩展到科学可视化（体数据）和通用物体表示（多模态输入）。

VVGT（2026.04）~\cite{2604_12217}提出前馈体积可视化框架（Visual Volume-Grounded Transformer），从多视角的体渲染图像直接预测三维高斯表示，无需逐场景优化。其核心挑战在于体数据没有明确的"表面"——传统 DVR（直接体渲染）在密集体素网格上逐射线积分，计算和存储开销巨大。VVGT 的双 Transformer 网络（DTN）包含两个互补分支：2D Transformer 从多视角渲染图像提取外观特征，3D Transformer 从初始体素表示捕获空间结构。两者通过 Volume Geometry Forcing 机制——一种极线交叉注意力——将多视角观测整合到分布式三维高斯基元中。VVGT 在 $4096^3$ 温度场和湍流模拟数据上训练，$512^3$ 子场景上测试，实现零样本泛化，将密集体数据转换为稀疏高斯表示后支持近瞬时交互式探索，渲染速度比传统 DVR 快数个数量级。

Object-X（2025.06）~\cite{2506_04789}提出多模态三维物体表示框架，将图像、点云、文本等多模态信息在三维体素网格中几何锚定。其核心设计是将每个物体编码为结构化潜在表示（Structured Latent Representation, SLat）：先将物体体素化为规范 3D 网格，在体素内聚合局部图像特征和点云几何特征，再通过稀疏 Transformer 编码器压缩为紧凑嵌入。该嵌入具有双重能力：(1) 通过 U-3DGS 解码器解码为三维高斯进行物体重建；(2) 直接作为描述符用于视觉定位、场景对齐和单图重建等下游任务，无需任务特定微调。存储需求比传统 3DGS 少 3--4 个数量级（体素分辨率可调，典型配置为 $32^3$ 或 $64^3$），在 3RScan 数据集上的场景对齐和物体重建任务中均展现优异性能。

\subsubsection{视觉语言导航}

本路线介绍将可泛化 3DGS 引入视觉语言导航任务的方法。UnitedVLN（2024.11）~\cite{2411_16053}将可泛化 3DGS 引入视觉语言导航（VLN-CE），解决连续环境中导航智能体的未来环境感知问题——智能体在执行导航指令时需要预见被遮挡或未探索区域的外观，以做出更好的路径决策。

UnitedVLN 的核心思路是构建"特征云"（feature cloud），将预训练 CLIP-ViT-B 提取的网格特征与神经点位置关联，形成同时编码外观和语义的三维神经表示。关键设计包括两种互补策略：(1) 搜索-查询采样（STQ）——先在特征云中搜索与目标位置最邻近的神经点，再以查询方式高效获取其特征，避免全局渲染的昂贵开销；(2) 分离-联合渲染（STU）——图像和语义特征先分别通过独立的 3DGS 和 NeRF 渲染管线渲染，再在特征层面联合融合，兼顾渲染质量和语义一致性。在大规模 HM-3D 数据集（800 个训练场景）上预训练后，UnitedVLN 在 R2R-CE 和 RxR-CE 两个 VLN-CE 基准上显著提升导航成功率，证明高保真的未来环境渲染能有效改善导航决策。

\subsubsection{光流-深度联合优化}

本路线介绍将光流与深度信息联合优化以提升稀疏视角前馈重建质量的方法。JointSplat（2025.06）~\cite{2506_03872}解决稀疏视角前馈 3DGS 中深度估计与光流匹配的互补利用问题。核心观察是：光流和深度包含互补的几何信息——光流在纹理丰富区域可靠但在遮挡区域失效，深度估计在弱纹理区域更稳定但在重复图案区域易出错。然而，简单融合两者的信息会引入新问题：遮挡区域的光流是无效的，直接使用会产生误导梯度。

JointSplat 提出概率优化机制：(1) 光流分支预测流匹配概率掩码，在训练期间根据匹配概率动态调节深度和流之间的信息融合比例——匹配概率高的区域以光流为主，概率低的区域以深度为主；(2) 多视角深度一致性损失通过比较不同视角深度预测的重投影误差，识别并抑制不确定区域的误导梯度。骨干采用 DepthSplat 的 ViT-B 单目分支，输入分辨率 $256 \times 256$，在 3 张 NVIDIA L40 GPU 上训练 400K 迭代。在 RealEstate10K 和 ACID 上的实验表明，JointSplat 一致超越 DepthSplat 等最佳方法，尤其在弱纹理和重复图案区域改进显著。

\begin{takeaway}
特定领域前馈方法将通用 3DGS 范式适配到事件相机、视觉语言导航等专业化场景，光流-深度联合优化等跨模态策略有效弥补了单一模态的固有缺陷，展现了前馈重建框架的广泛适用性。
\end{takeaway}

%=====================================================================

\section{近期研究（2026 年 8--11 月）}
\paragraph{Swimm3R：先恢复水下几何，再联合拟合场景与介质}
2026.08~\cite{2608_00950}。普通 SfM 会把水体散射造成的对比度损失和颜色变化当成观测误差。Swimm3R 将 MASt3R-SfM 适配到水下图像，先估计恢复后的稠密点云、相机位姿与光学参数，再据此联合拟合水下 Beta 原语和显式介质场。散射感知的几何梯度与边缘感知的逆深度正则共同限制几何和水体效应的相互补偿。

在 FLSea 的几何评价里，相对作者列出的 MASt3R-SfM 基线，AbsRel 从 0.2247 降至 0.2132。对保留帧的渲染比较固定场景地图，但用各测试图像细化相机位姿 300 步，并对各方法沿用同一协议。部分位姿比较依赖重建参考轨迹，Beta 原语的渲染速度也低于标准三维高斯，宜分别看待位姿、保留帧画质和渲染成本。

\paragraph{QuerySplat：用独立查询分别回归几何与外观}
2026.08~\cite{2608_01186}。逐像素分配高斯会把原语数量绑定到输入视图与像素，而以场景查询统一回归所有属性又可能把空间结构与高频颜色混在一起。QuerySplat 使用双分支查询解码器，几何分支借助预训练的视觉几何模型提供坐标系、相机估计和三维特征，外观分支专门恢复辐射属性。输出仍是可直接泼溅的高斯，变化发生在高斯的分配和属性回归上。

论文在 DL3DV 的新视角合成比较中报告，相对该协议内最佳无位姿和需位姿基线，平均 PSNR 分别提高 2.30 dB 和 1.04 dB。这是针对论文所用视图和基线的比较。可扩展性仍需结合查询数、输入视图数和显存开销来判断。

\paragraph{StreamSplat：用三维缓存承接连续输入}
2026.08~\cite{2608_01659}。StreamSplat 用已标定的视图逐批构建三维高斯。在采用 DepthSplat 几何骨干与 ReSplat 式解码结构的基础上，方法将历史特征保存于有容量限制的场景坐标体素缓存。下一批输入到达时，缓存几何投影到当前视图以引导深度估计，邻近历史特征再影响高斯原语回归，因此每批输入后都能解码出可渲染场景。

原文在 DL3DV、RealEstate10K 和 ScanNet 上考察稀疏输入及长轨迹，对照 DepthSplat、AnySplat、ReSplat、Long-LRM 和使用官方权重的 OF³GS。长序列实验使用 256、512、1024 视图。方法针对静态场景并依赖位姿精度，不显式处理运动物体或光照改变。固定体素大小和缓存预算也在内存与几何细节之间形成取舍。

\paragraph{UniqueSplat：随查询视角调整前馈高斯}
2026.08~\cite{2608_02145}。固定高斯场从所有目标视角渲染同一组高斯。UniqueSplat 通过视角条件超网络，把目标视角相关的信息注入预测过程，在共享场景信息的基础上为每个查询生成适配的高斯。这给大角度视点变化提供了额外的调整自由度，也使表示的使用方式从一次预测后复用，变为由查询视角参与预测。

论文分别在 RealEstate10K、ACID 和 DTU 上评价新视角合成。作者在局限部分指出，高斯仍按输入视图像素对齐预测，未观测区域在渲染时仍可能出现伪影。这个限制要与已观测区域的画质改善分别评价。

\paragraph{InfiniSplat：从固定像素格改为沿表面采样高斯}
2026.08~\cite{2608_02437}。单图重建若在每个像素中心放置原语，较大基线下容易显露源相机采样的空隙。InfiniSplat 用单目深度或带测距输入的几何先验决定不规则支撑点，再向隐式解码器查询图像特征和高斯属性，并对初始高斯作有界更新。预测不需要逐场景拟合。

原文将 RGB 与 LiDAR 条件的模型分别同前馈基线比较，涉及 ETH3D、ScanNet++、Tanks and Temples 与 DL3DV。两种输入不能视作相同观测条件。论文展示了大面积反遮挡、错误深度先验和极端视角外推时的缺失或伪影，不能从一帧恢复完全不可见的结构。

\paragraph{FlexSplat：用变形查询汇聚未标定对象图像}
2026.08~\cite{2608_07937}。FlexSplat 对未标定的对象多视图联合训练几何变换器和高斯解码器，以预测深度建立空间锚定，再让多视图变形注意力汇聚证据。查询数限定高斯预算，与输入分辨率解耦，也不需要为每个输出原语建立点云对应。

作者在 GSO 将其与使用已知位姿的对象重建器比较，优势主要在感知画质与省去标定，PSNR 并非全面超过最强有位姿方法。模型训练与测试限于对象中心捕获，大场景背景和大的深度或位姿偏差仍会影响重建，平移扰动也有待完整评测。

\paragraph{LocusGS：让前馈查询在三维空间分工}
2026.08~\cite{2608_12825}。固定数量的全局查询能限制高斯输出预算，但一个查询生成的高斯可能散布在不相邻的区域。LocusGS 为每个输出查询维护三维中心和尺度，沿解码层更新。锚点既参与从多视图标定图像中选择相关证据，也约束该查询所解码高斯的局部范围，从而让不同查询负责互补区域。

作者在 RealEstate10K 的两视图设置，以相同查询数和高斯预算对照 TokenGS，在 PSNR、SSIM 和 LPIPS 上报告改善；DL3DV 另测从四视图训练到不同测试视图数的泛化。射线聚合仍依赖已标定相机，只有中心与标量半径的锚点也难以描述细长或高度各向异性的结构。

\paragraph{SplatGuide：从同一高斯重建提取三种生成条件}
2026.08~\cite{2608_16863}。缺少输入相机位姿时，二维视频生成先验难以独立维持跨视角几何。SplatGuide 用冻结的前馈重建骨干从静态场景参考图像估计位姿和高斯，再为新视图提供像素级渲染、特征 token 与可见性驱动的参考视图选择。最终渲染结果由扩散生成器细化，高斯在这里承担条件和几何接口。

论文区分域内和域外基准，并通过参考视图选择与条件消融评价各信号。三种条件来自同一重建骨干，因此纹理稀少、视角基线过宽或物体运动导致几何失准时，它们可能同时失效。目标视角按指定相机位姿查询，输入参考图无已知位姿这一条件不应与目标轨迹的定义混为一谈。

\paragraph{SeeU：用跨视图语义证据细化高斯}
2026.08~\cite{2608_22740}。SeeU 从带相机参数的稀疏图像估计初始高斯，再以跨视图熵感知的语义嵌入为条件，用 Transformer 在高斯空间回归残差。它改的是初始高斯的表示与细化阶段，输入图像依然提供场景观测约束。

作者在 RealEstate10K 和 ACID 的域内协议中用两张参考图渲染三个目标视角。RealEstate10K 训练后直接评估 DTU 时，SeeU 的 PSNR 为 16.31 dB，HiSplat 为 16.05 dB；在跨域 ACID 对比中，HiSplat 的 LPIPS 略优。两组结果适合分别讨论跨域几何与感知画质，不能据此推断到未知输入位姿。

\paragraph{利用逐视图高斯子集过滤瞬态干扰}
2026.08~\cite{2608_26951}。逐视图预测的前馈模型会为各输入视角生成相应的一组高斯。该工作利用这些集合的来源信息，将某个视图的高斯暂时排除，观察其余视图对候选重建内容的支持，再通过渲染检查删除候选是否能改善重建。模型权重保持不变，处理发生在推断后对候选高斯的筛选阶段。

论文在 RobustNeRF 与 NeRF On-the-go 上比较 4、8 和 16 个带干扰的输入视图，并检查干净场景中的画质。验证每个候选都需要渲染，因此计算开销随候选数和输入视图数增长；在多视角里一致出现的内容也不能仅靠单视图来源判为干扰。

\paragraph{ReconSplat：用可投影的高斯潜变量约束跨视图扩散}
2026.08~\cite{2608_28895}。稀疏宽基线输入中，可见结构与未观测区域需要不同的重建证据。ReconSplat 先从有位姿图像预测承载潜变量的高斯，再将这组三维潜变量投影到各目标视图，给多视角扩散提供一致的空间条件。模型同时支持 RGB 与深度外推，并从预测中恢复显式点云。

RealEstate10K 的内插与外推使用不同目标视图；论文报告外推中的多数指标优于所列基线，但 SSIM 例外。作者指出初始高斯的浮游伪影、缺失结构或不精确几何会传递到扩散条件。未观测区的生成结果可用于观察外推能力，其几何真实性还需要独立观测检验。

\paragraph{DualDiff3D：用结构与外观扩散先验迭代修正高斯}
2026.09~\cite{2609_01516}。宽基线高斯渲染的伪影既可能来自几何缺失，也可能来自外观细节不足。DualDiff3D 通过双分支扩散处理结构和外观，以结构感知注意力耦合两者，再执行渲染、修复和重新优化循环。置信度掩码用于选择修复结果可参与后续高斯优化的区域。

作者在 DL3DV 和 LLFF 分别评价逐图修复与最终高斯重建，并与常规 3DGS 和 DIFIX 系列比较。扩散器采用单步设计，但整个方法仍有额外算力开销，输入极度稀疏时也可能退化。因此应把单步修复和完整场景优化的耗时分开。

\paragraph{用三维高斯扰动构造视频修复训练对}
2026.09~\cite{2609_03657}。为学习能修复跨视图伪影的视频先验，通常要从多种欠观测条件反复拟合场景，再渲染训练视频。该工作直接扰动已拟合高斯的缩放、旋转和保留状态，生成带空间对应的退化视频，与同场景的干净渲染配对。免优化指构造退化场景的过程，视频先验仍需训练。

论文在多个稀疏观测设置比较二维退化和不同三维扰动，另用扩展模型评价深度与下游策略表现。三维扰动提供的是一种训练数据生成近似，对特定相机、退化分布和重建问题的覆盖仍应在目标数据上验证。

\paragraph{AVSplat：在多视图全局聚合之前引入辅助视图}
2026.09~\cite{2609_05925}。输入视图增多本应改善覆盖，但直接用全局注意力融合时，证据未必集中在相互支持的观测上。AVSplat 在全局聚合前让每个输入视图先从少量辅助视图获取场景上下文，再以占据情况调节体素融合权重，缓和致密输入下属性被平均得过于平滑的问题。

作者在 Mip-NeRF 360、VR-NeRF 和 DL3DV 比较稀疏与致密输入，报告相对所比较的无位姿前馈模型，增加视图时的结果更稳定。该结果说明输入规模利用率有改善，具体到显存和时延，还需按实际分辨率与输入量比较。

\paragraph{VoxelTTO：用体素聚合消除像素冗余，按相机先验做测试时校正}
2026.09~\cite{2609_21498}。像素对齐法会让不同视角为同一表面重复预测基元。VoxelTTO 将密集特征聚合到共享体素，用稀疏三维卷积和可分裂解码器回归高斯属性。论文还用随机实心体积渲染替换训练和推断中的常规高斯光栅化。

作者在 DTU、Replica 与 Tanks and Temples 上评估重建。测试时优化通过修改骨干网络中的轻量 LoRA 层，让预测相机与给定相机先验对齐。这一项依赖可用的相机监督，需要把纯前馈配置与带相机先验及测试时优化的配置分开评价。

\paragraph{GrapeSplat：在占用体素而非图像网格上生成高斯}
2026.09~\cite{2609_23182}。从每幅图固定位置解码高斯会随输入视图增加而产生重复原语。GrapeSplat 将无位姿图像的预测几何和特征提升到共同三维空间，在可压缩的体素坐标中汇总。稀疏解码器从占用体素生成初始高斯，使规模更多取决于占据空间，而非输入像素数。

论文在 NRGBD、7-Scenes、ScanNet++ 和 DL3DV 使用不同数量的上下文图像评测，并在附录扩展到 DTU 与 ETH3D。输入视图数变化和场景尺度变化是不同的泛化轴。补充材料的 DTU 示例仍有背景变体在相同区域失败，因此不能将体素汇总当成对所有上游几何错误的修正。

\paragraph{Spackle：把单图大视角补全交给高斯残差分支}
2026.09~\cite{2609_30941}。仅凭一张图渲染偏离很远的视角时，新露出的区域和原本可见的区域会争夺固定原语预算。Spackle 冻结预训练前馈模型，用渲染结果中的缺损掩码限定残差分支的训练区域，并把生成补全用作伪监督。推断时合并基础与残差高斯，无需再执行生成式补图。

论文在 Middlebury、Booster、WildRGBD、ETH3D、LLFF 和 Tanks and Temples 的大视点偏移协议下比较补全和原可见区域表现。新增区域的合理外观不代表恢复了真实遮挡面，补全质量受缺损掩码、伪监督与基础重建质量影响。

\paragraph{AESplat：将可见基本颜色与方向相关外观分开预测}
2026.09~\cite{2609_36693}。无位姿前馈高斯需要同时处理几何和外观。AESplat 直接用输入 RGB 构造球谐的零阶分量，只将高阶项交给带三维归纳偏置的预测模块。观测颜色提供基色，高阶分量负责补偿随方向变化的剩余项。

按原文对 NAS3R 的比较，RealEstate10K 与 ACID 的 PSNR 分别提高 0.803 与 0.847 dB。这个增益针对新视图合成。原文指出，完全不在输入视图中的内容难以可靠重建，用单次 RGB 近似方向无关颜色在强曝光差异或非朗伯反射下也有偏差。

\paragraph{StereoGaussians：从单对标定双目图像预测附近视角}
2026.09~\cite{2609_38592}。StereoGaussians 在冻结的立体网络中提取视差和多尺度特征，再用高斯预测头构建米制场景。第二层高斯与扩展画布给反遮挡和原视野外区域留出表达容量。SceneSplat-Stereo 用静态高斯教师渲染目标视图，提供成对训练监督。

论文在真实 StereoNVS 的八个留出场景与由已优化教师渲染的 ReplicaGS 分开评测；后者并非独立实拍目标。前层深度固定限制了对立体估计错误的修正，教师伪影仍会传入。原文评测集中在相近相机，不能推断大角度反遮挡效果。

\paragraph{3DROID：把机器人工作空间高斯锚定到米制坐标并单独测量可靠性}
2026.10~\cite{2610_01744}。视觉前馈模型生成的高斯可以合成图像，尺度与坐标却可能不适合机器人工作空间。3DROID 将精化相机外参与有位姿预测结合，并用覆盖感知的渲染深度、双目几何和机器人运动学表面一致性另外评价重建。相机外参提供对全局坐标的约束，度量几何精度仍需测量。

在 110 个场景的保留 ext2R 相机上，相对默认无位姿配置，YoNoSplat 的 PSNR 从 12.57 提升到 15.81 dB，ZipSplat 从 12.75 提升到 16.27 dB。该比较同时改变外参来源和位姿条件。数据目前以静态帧为主，几何评价限于可由立体匹配或机器人表面支持的位置，未纳入腕部相机。

\paragraph{LoCoSplat：把三维场景聚合简化为固定局部均值}
2026.10~\cite{2610_04351}。LoCoSplat 保留以深度和图像特征构造体素锚点的二维阶段，把学习式三维稀疏网络换成固定的局部特征平均，并通过逐点感知机与高斯预测头解码。高斯生成仍依赖二维阶段给出的深度和已知输入位姿，改动集中在锚点附近的三维上下文融合。

在 RealEstate10K 的 6、12、24 视图设置中，作者把方法与按相同协议重新训练的 VolSplat 对比。在 24 视图设置中，作者报告 30.57 dB，配对差为 3.3 dB。其他训练语料和分辨率的发布模型在表中单独标注，不能从该实验推断无位姿下的速度和画质。

\paragraph{DensiTok：在几何潜空间补齐未观测的中间相机}
2026.10~\cite{2610_07958}。前馈高斯骨干看到的真实视图很少时，DensiTok 在它的几何令牌空间增加输入。层级自适应 VAE 压缩各层几何令牌，相机条件扩散变换器一步预测缺失中间视图的潜码，再解码为骨干重建头能用的令牌。这避免先生成完整图像再重跑骨干编码器。

论文对 AnySplat、WorldMirror 和 Depth Anything 3 在 RealEstate10K 与 DL3DV 上比较画质、位姿、深度或运行时间。虚拟相机只在首末观测相机之间插值，论文也在这个区间取评测视图。由此得到的结果无法说明训练轨迹之外的视图外推表现。

\paragraph{DeltaSplat：把图像残差变成高斯属性的迭代更新}
2026.10~\cite{2610_09853}。DeltaSplat 用相机射线与当前渲染深度把二维残差定位到三维，学习逐步更新已有高斯的位置、不透明度与颜色。迭代是在前馈网络内运行的轻量细化模块，无需为每个场景反向优化，适合接在已有无位姿重建骨干之后。

作者在 DL3DV 的 140 场景测试集上比较 6、12 和 24 个输入视图。已知内参但无已知外参时，6 视图的 PSNR 为 26.64 dB，较所用前馈骨干提升 1.75 dB；论文另列全未标定设置。输入分辨率为 224 乘 224，高分辨率的计算成本和精度仍需单独评估。

\paragraph{CompactSplat：按局部纹理复杂度分配前馈高斯}
2026.10~\cite{2610_13177}。逐像素预测使平坦区域消耗大量相似高斯。CompactSplat 用纹理引导的区域划分配合分层金字塔采样，让复杂和简单区域采用不同密度，并允许推断时控制稀疏度。改进发生在原语空间分配，输出仍可由高斯渲染器消费。

论文在 RealEstate10K 的非像素对齐方法对照中报告，CompactSplat-L 使用约 72.6k 个高斯，LPIPS 为 0.117；TokenGS 用约 66k 个高斯且 PSNR 更高，LPIPS 为 0.147。画质随指标与预算变化。作者还指出划分依赖正交网格，主要消除单视角内冗余，跨视角的重复结构仍有压缩空间。

\paragraph{Holi-Spatial 360：用匹配约束稀疏全景重建并生成空间标注}
2026.10~\cite{2610_19502}。稀疏全景相距较远时，图像修复容易改变房间原有结构。Holi-Spatial 360 用 DA3 的匹配给高斯优化提供几何锚定，再用适配的图像与视频模型补充连续轨迹视图。重建和 RGB-D 验证还被用于生成三维实例、描述与空间问答，构成可用于空间推理的数据。

10 个 RealSee3D 场景的 LiDAR 对照中，各方法采用相同 DA3 初始化，2 厘米阈值的 F-score 从 DA3 的 0.2207 提高到 0.2617。该检验覆盖的是这些室内场景的几何，未观测区域的生成内容仍受对应关系可靠性和参考覆盖限制。

\paragraph{InfiniSplat v2：固定几何中心，用可学习属性筛掉低质量观测}
2026.10~\cite{2610_20341}。InfiniSplat v2 用冻结的 DA3 从图像和可选相机条件估计深度与位姿，再把点云转成高斯。前馈属性分支通过渲染监督学习每个点的贡献，后续用输入视图的渲染反馈调整属性。固定中心将几何定位与外观学习分开，学习的不透明度可抑制几何不一致的点。

论文分别评价有位姿和无位姿配置，报告六个数据集上平均 LPIPS 与 DISTS 的改善，并在 ETH3D 和 DTU 做消融。抑制错误高斯有助于画质，但固定中心无法修复严重的几何位置误差，像素对齐产生的重复原语和未观测空洞仍存在。

\paragraph{Anchor3DGS：以共享三维锚点替代逐像素生成}
2026.10~\cite{2610_21634}。Anchor3DGS 从未标定输入构造几何支架，以固定数量的三维锚点聚合可见多视图信息，在锚点之间交换局部结构上下文，再把锚点解码为局部高斯。锚点预算不再随像素或视图数直接增长，并通过熵引导将更多容量分配给复杂区域。

论文以 3 和 6 个输入视图在 VR-NeRF、MipNeRF360 与独立的 DL3DV 子集评价，在 MipNeRF360 与 DL3DV 上报告超过 2 dB 的 PSNR 增益。该数字对应论文所列的稀疏设置和基线；极少视图中的结构仍受支架质量、锚点分辨率和遮挡影响。

\paragraph{在无位姿前馈串联中校正坐标约定和尺度}
2026.10~\cite{2610_23039}。直接把相机估计网络的输出接到高斯预测器时，坐标定义或平移尺度不一致会引入误差。该工作保持 VGGT 和 TokenGS 冻结，在接口处转换坐标约定并做平移尺度校正。作者还单独比较从全部场景视图估计位姿和仅用稀疏上下文的情况。

在 140 个 DL3DV 场景上，全视图估姿配置报告 22.47 dB，四视图输入配置报告 22.44 dB，所复现的 COLMAP 条件为 22.98 dB。两种观测权限应分开列出，不能把全视图位姿结果直接归到四图重建协议。

\paragraph{FineSplat：用视图分组和三维分组保留小块细节}
2026.10~\cite{2610_25011}。细粒度图像块能保留细节，但直接对高分辨率多视图做全局注意力代价很高。FineSplat 先在少数视图组成的小组内处理 4 乘 4 像素图块，再把初始高斯渲染回参考视图，用残差反馈组织三维空间分组，让不同视图组的证据沿三维对应关系交互。

论文给出 DL3DV 的 6、12 和 24 视图结果及跨域比较。在无位姿评价中，论文冻结高斯并对每个目标相机位姿优化 200 步进行对齐；有真实位姿时不做额外对齐。这是测量画质的对齐协议，需与模型的前馈推断时间分开理解。

\paragraph{RLSplat：按预算直接生成高斯并限制解码成本}
2026.10~\cite{2610_27712}。在密集预测后剪枝只能降低最终高斯数，前端仍会为被删原语支付预测成本。RLSplat 的分配策略先确定各锚点分别生成多少高斯，再从指定总预算直接解码，以局部交叉注意力细化。预算因此同时作用于表示规模和解码成本。

作者在 ACID 按每视图 1k、2k 和 4k 个高斯与预算可控方法比较，在 8 视图、每视图 1k 个时相对 F4Splat 报告 1.22 dB 的 PSNR 优势。图像编码仍有固定开销，不能把解码成本随预算下降外推成整个系统同幅加速。实验还限于 256 乘 256 分辨率与一种跨域设置。

\paragraph{SG-Splat：在位姿估计有误时分阶段修正重建}
2026.10~\cite{2610_35361}。SG-Splat 将几何估计与高斯属性预测分开。模型先用光流引导细化位姿，再利用残差修正几何。双路径训练混合真实和预测位姿并使用不同的优化方式，使网络既能保留位姿监督，又能适应推断时的位姿误差。

论文在 DL3DV 和 RealEstate10K 比较稀疏与较密输入的新视角合成及相对位姿。该方法需要已知相机内参，高斯仍按像素对齐预测，输入视图重叠和较高分辨率都可能带来额外开销。

\paragraph{GeoGSplat：先以锚点预测空间结构，再在固定预算内解码高斯}
2026.10~\cite{2610_36760}。GeoGSplat 结合锚点与三平面特征作为有结构的场景表示，分层解码时先在锚点附近预测稠密三维点，再回归高斯。这样把固定容量的非像素对齐预测连接到显式几何，而无需让每张输入图都直接生成一个像素对齐的原语集合。

论文设置 4096 个锚点，每个锚点预测 16 个点，共 65536 个高斯，并在 DTU 与 NRGBD 评价合成和重建。这是固定的原语容量。比较不同视图数和复杂度的场景时，需要同时考虑预算能否覆盖稀少结构及训练和推断成本。

\paragraph{ThermalSplat：从无位姿稀疏热图预测高斯}
2026.11~\cite{2611_15435}。热图的纹理和辐射响应不同于 RGB。ThermalSplat 微调无位姿几何骨干，以软 Census 损失监督对亮度变化较稳健的局部结构，再用边界损失约束温度突变，从两张稀疏热图前馈得到可渲染高斯。

逐场景热重建基线的对照按各方法预设的保留视图计算指标，另一组泛化评估使用未见场景里所有可用评价视图。两种协议对应不同问题。论文明确假设静态场景，宽基线少图的未观测区域也仍有几何歧义。

\paragraph{PaveSplat：用残差旋转校正表面对齐的协方差}
2026.11~\cite{2611_26045}。沿预测表面对齐的解析协方差给出了合理的高斯形状初值，但局部估计还可能偏离真实切平面。PaveSplat 保留 SurfSplat 类的解析初值，将切面内的尺度与残差四元数用于局部校正。残差从单位旋转初始化，学习的变化集中在每个表面基元的协方差方向。

论文在 RealEstate10K 使用两张输入、256 乘 256 训练及标准测试划分，并在相同训练配置下比较不同渲染尺度。跨尺度表现与训练配方和表面深度质量有关，应与完全不同的数据域或原语预算分开比较。

\paragraph{FFGaussType：把已知字形轮廓转为可复用高斯资产}
2026.11~\cite{2611_26431}。FFGaussType 面向已知轮廓的挤出字形与 SVG 图标。方法从采样锚点预测二维特征平面，通过跨资产共享的解码器把平面特征转成带法向的神经高斯。相比逐资产拟合，学习的对象是从给定轮廓几何到高斯资产的转换。

论文按字体家族或图标来源拆分数据，报告全测试集法向误差 1.62 度，来自已准备锚点的转换耗时为 9.1 毫秒。该时间不含轮廓锚点的前处理。论文仅研究具有均一材质的挤出资产，着色不处理自阴影或与周边场景的光传输。


\section*{本章小结}
\addcontentsline{toc}{section}{本章小结}

\begin{futurebox}{本章总结与展望}{ch13_summary}
本章综述了可泛化前馈重建这一 3DGS 领域的前沿方向，涵盖以下核心内容：
\begin{itemize}
\item \textbf{多视角前馈重建}（第~\ref{sec:13_1}~节，71 篇）：从像素对齐到体素对齐、从固定密度到自适应分配、从单次推断到迭代精炼，形成了完整的技术谱系。DepthSplat、MVSGaussian 等确立了基础管线；SparseSplat、GlobalSplat、TokenGS 等突破了像素对齐的表示瓶颈；ReSplat、GIFSplat、Diff3R 等引入迭代精炼策略提升质量；生成先验的融合（MVSplat360、GeoNVS、ConFixGS）进一步弥合了前馈推断与逐场景优化之间的质量差距。
\item \textbf{无标定与自由视角}（第~\ref{sec:13_2}~节，26 篇）：NoPoSplat 的坐标规范化、Splatt3R 和 FreeSplatter 的基础模型迁移、AnySplat 的单次全能推断，逐步消除了对已知相机位姿的依赖，使前馈 3DGS 从受控采集走向随意拍摄场景。
\item \textbf{单图到 3D 前馈方法}（第~\ref{sec:13_3}~节，8 篇）：通过教师-学生蒸馏（studentSplat）、视频扩散潜在空间（Wonderland）、几何引导生成（G4Splat）和视觉语言模型辅助（CATSplat），在极端信息缺失条件下实现了从单张图像到三维场景的映射。
\item \textbf{大规模重建模型}（第~\ref{sec:13_4}~节，9 篇）：从 GS-LRM 的 2--4 视角到 Long-LRM 的 32 视角再到 VGGT-X 的 1000+ 视角，Mamba-Transformer 混合架构、迭代精炼和测试时训练等技术使前馈重建突破了输入规模的限制。
\item \textbf{特定领域前馈方法}（第~\ref{sec:13_5}~节，11 篇）：从事件相机（EvGGS）到射频信号（XFreq-GS、RxGS）、从人体通信（GPS-Gaussian+、VoluMe）到体积可视化（VVGT），展示了前馈 3DGS 范式在多个领域的广泛适用性。
\end{itemize}

可泛化前馈重建面临的挑战与未来方向包括：
\begin{itemize}
\item \textbf{质量-效率的根本权衡}：前馈方法的渲染质量仍落后于逐场景优化方法。虽然迭代精炼（ReSplat、GIFSplat）和优化感知训练（ForeSplat）在缩小差距，但如何在保持单次前馈效率的同时逼近优化方法的质量极限，仍是核心挑战。
\item \textbf{表示效率}：像素对齐方法产生的高斯数量随输入增长爆炸，自适应分配（SparseSplat、GlobalSplat）和解耦表示（TokenGS、LGTM）是重要方向，但紧凑性与渲染质量之间的平衡仍需进一步探索。
\item \textbf{长序列与大规模场景}：虽然 Long-LRM 和 MVP 将输入扩展到数十张视角，但前馈方法在完整城市级别的大规模重建上仍面临计算瓶颈。如何在保持全局一致性的同时处理数千张图像的输入，是需要解决的关键问题。
\item \textbf{动态场景的前馈重建}：本章主要关注静态场景的前馈重建。将前馈范式扩展到动态场景——从时序输入直接预测动态三维高斯表示——是一个重要但尚处于早期阶段的方向（Fuse-and-Refine 的初步尝试表明这一方向的潜力）。
\item \textbf{多模态统一}：当前方法针对不同模态（RGB、事件、射频等）独立设计适配策略，缺乏统一的前馈框架。如何构建一个通用的多模态前馈 3DGS 模型，使其能够处理各种输入类型并输出统一的三维表示，是具有挑战性的开放问题。
\end{itemize}
\end{futurebox}
